% Lutte contre les troubles liés à la régulation du taux d'hormones sexuelles chez l'homme et chez la femme et la stérilité — SVT Tle D — Cours signé OnBuch+
% Programme officiel MINESEC. Compiler avec : tectonic N12-regulation-hormones-sexuelles-sterilite.tex
\newcommand{\DOCMATIERE}{SVT}
\newcommand{\DOCNIVEAU}{Tle D}
\newcommand{\DOCMODULE}{Module II : L'Éducation à la Santé — La santé reproductive}
\newcommand{\DOCLECON}{N12}
\newcommand{\DOCTITRE}{Lutte contre les troubles liés à la régulation du taux d'hormones sexuelles chez l'homme et chez la femme et la stérilité}
% OnBuch+ — préambule « cours signés » ENRICHI (compilé avec tectonic / XeLaTeX)
% Système visuel « Pop » d'OnBuch+ : fond crème (jamais blanc), bordures encre
% épaisses, ombres « dures » décalées, titres Archivo Black en capitales.
% Version MATHÉMATIQUES : graphes (pgfplots/tikz), boîtes pédagogiques
% (définition, propriété, méthode, attention, le savais-tu, exercice résolu,
% exercice, TP, bilan), page de garde et signature OnBuch+ ancrée en bas.
%
% Macros attendues dans main.tex AVANT \input{preamble} :
%   \DOCMATIERE  \DOCNIVEAU  \DOCMODULE  \DOCLECON (numéro)  \DOCTITRE
\documentclass[11pt]{article}

\usepackage{fontspec}
\usepackage[french]{babel}
\usepackage[a4paper,margin=2cm,top=3.4cm,bottom=2.8cm,headheight=1.4cm,footskip=1.3cm]{geometry}
\usepackage{amsmath,amssymb,amsfonts}
\usepackage{mathtools} % \xmapsto, \coloneqq... souvent utilisés par les modèles
\usepackage{cancel} % simplifications barrées (bilans de forces)
% \abs{z} (module, valeur absolue) et \norm{u} (norme), utilisés par les modèles
\providecommand{\abs}{}\let\abs\relax\DeclarePairedDelimiter{\abs}{\lvert}{\rvert}
\providecommand{\norm}{}\let\norm\relax\DeclarePairedDelimiter{\norm}{\lVert}{\rVert}
\usepackage[table]{xcolor} % [table] : fournit aussi \rowcolors (alternance de lignes)
\usepackage{pagecolor}
\usepackage{tikz}
\usetikzlibrary{arrows.meta,positioning,calc,shapes.geometric,shapes.misc,shapes.arrows,decorations.pathmorphing,decorations.pathreplacing,decorations.markings,patterns,shadows,mindmap,trees,fit,backgrounds,matrix,intersections,angles,quotes}
\usepackage{pgfplots}
\usepackage[version=4]{mhchem} % équations nucléaires \ce{^{235}_{92}U}
\usepackage[european,siunitx]{circuitikz} % schémas électriques
\pgfplotsset{compat=1.18}
\usepgfplotslibrary{fillbetween,groupplots} % aires entre courbes (calcul intégral)
\usepackage{enumitem}
\usepackage{fancyhdr}
% [most] charge toutes les bibliothèques tcolorbox (listings, minted, raster,
% documentation...) et gonfle inutilement la mémoire TeX sur un document long
% (~300 boîtes) : on ne charge que ce qui est réellement utilisé.
\usepackage[skins,breakable]{tcolorbox}
\usepackage{graphicx}
\usepackage{colortbl}
\usepackage{booktabs}
\usepackage{tabularx}
\usepackage{diagbox} % case d'en-tête diagonale des tableaux à double entrée
% Colonnes extensibles centrée / gauche / droite (C, L, R), souvent utilisées par les modèles
\newcolumntype{C}{>{\centering\arraybackslash}X}
\newcolumntype{L}{>{\raggedright\arraybackslash}X}
\newcolumntype{R}{>{\raggedleft\arraybackslash}X}
\usepackage{array}
\usepackage{multicol}
\usepackage{siunitx}
% parse-numbers=false : accepte aussi \SI{2,5 \times 10^{-3}}{...}
\sisetup{locale=FR,output-decimal-marker={,},group-minimum-digits=4,parse-numbers=false}
\DeclareSIUnit\ohm{\ensuremath{\symup{\Omega}}} % U+2126 absent de la police
\DeclareSIUnit\division{div} % réglages d'oscilloscope (ms/div, V/div)
\DeclareSIUnit\tr{tr} % tours (vitesse de rotation, ex. \SI{1500}{\tr\per\minute})
\DeclareSIUnit\molar{M} % concentration molaire (ex. \SI{3}{\molar}, \SI{1}{\micro\molar})
\DeclareSIUnit\an{an} % année (le modèle confond parfois avec \year, un compteur TeX) — ex. \SI{5}{\kilo\meter\per\an}
\DeclareSIUnit\FCFA{FCFA} % franc CFA (ex. \SI{500}{\FCFA\per\kilo\gram})
\DeclareSIUnit\degreeBrix{\textdegree Brix} % teneur en sucre d'un sirop/jus (réfractométrie)
\DeclareSIUnit\month{mois} % le modèle utilise parfois \month, un compteur TeX, comme unité de durée
\usepackage{qrcode}
% \degree : commande fréquemment utilisée par les modèles pour un angle ou une
% température en dehors de \SI{}{} (non fournie par les paquets chargés).
\providecommand{\degree}{^{\circ}}
% \qed (fin de démonstration), utilisé par les modèles sans amsthm
\providecommand{\qed}{\hfill$\square$}
% \barre : double barre des valeurs interdites dans un tableau de variations (tkz-tab)
\providecommand{\barre}{\ensuremath{\|}}
% environnement proof (amsthm non chargé) utilisé par les modèles
\ifdefined\proof\else\newenvironment{proof}[1][Démonstration]{\par\noindent\textit{#1.}\ }{\qed\par}\fi
\usepackage{lastpage}
\usepackage{hyperref}
\hypersetup{hidelinks}

% ============================================================
% Palette « Pop » OnBuch+ — mêmes hex que la classe P (plus_ui.dart)
% ============================================================
\definecolor{poporange}{HTML}{F59321}
\definecolor{popdark}{HTML}{E07A0C}
\definecolor{popcream}{HTML}{FFF6E8}
\definecolor{popink}{HTML}{1C1714}
\definecolor{poppurple}{HTML}{7A5AE0}
\definecolor{popgreen}{HTML}{2E9E6A}
\definecolor{poppink}{HTML}{E0457B}
\definecolor{popblue}{HTML}{2F7DE1}
\definecolor{popgold}{HTML}{C98A12}
\definecolor{popmuted}{HTML}{8A7F76}
\definecolor{popline}{HTML}{EADFD2}
\definecolor{popgray}{HTML}{8A7F76}
\colorlet{popgrey}{popgray}
% Alias tolérants (le modèle nomme parfois les couleurs différemment)
\colorlet{popwhite}{white}
\colorlet{popblack}{popink}
\colorlet{popred}{poppink}
\colorlet{popyellow}{popgold}
% Teintes claires (fonds de tableaux/encadrés) — le modèle nomme parfois
% les couleurs avec un suffixe L (ex. popblueL) sans les avoir définies.
\colorlet{poporangeL}{poporange!20}
\colorlet{popdarkL}{popdark!20}
\colorlet{poppurpleL}{poppurple!20}
\colorlet{popgreenL}{popgreen!20}
\colorlet{poppinkL}{poppink!20}
\colorlet{popblueL}{popblue!20}
\colorlet{popgoldL}{popgold!20}
\colorlet{popredL}{popred!20}
\colorlet{popgrayL}{popgray!20}
% Teintes claires pour fonds de tableaux / zones
\colorlet{popblueL}{popblue!10!white}
\colorlet{poporangeL}{poporange!14!white}
\colorlet{popgreenL}{popgreen!12!white}
\colorlet{poppurpleL}{poppurple!12!white}
\colorlet{poppinkL}{poppink!12!white}
\colorlet{popgoldL}{popgold!14!white}

\pagecolor{popcream}

% ============================================================
% Typographie
% ============================================================
\defaultfontfeatures{Ligatures=TeX}
\setmainfont{PlusJakartaSans-Medium.ttf}[Path=../../../fonts/,BoldFont=PlusJakartaSans-Bold.ttf]
% Maths en sans-serif (Fira Math) avec chiffres et lettres droites de Plus
% Jakarta, pour que les formules se fondent dans le texte.
\usepackage[math-style=ISO,bold-style=ISO]{unicode-math}
\setmathfont{FiraMath-Regular.otf}[Path=../../../fonts/]
\setmathfont{PlusJakartaSans-Medium.ttf}[Path=../../../fonts/,range={up/num,up/{Latin,latin},\mathup}]
\usepackage{icomma} % virgule décimale sans espace en mode math
% Caractères mathématiques Unicode absents des polices de texte (Plus Jakarta,
% Archivo Black) : rendus par leur équivalent mathématique, en texte comme en
% formule — sinon ils s'affichent en carré vide (ex. « Arithmétique dans ℤ »).
\usepackage{newunicodechar}
\newunicodechar{ℝ}{\ensuremath{\mathbb{R}}}
\newunicodechar{ℤ}{\ensuremath{\mathbb{Z}}}
\newunicodechar{ℂ}{\ensuremath{\mathbb{C}}}
\newunicodechar{ℕ}{\ensuremath{\mathbb{N}}}
\newunicodechar{ℚ}{\ensuremath{\mathbb{Q}}}
\newunicodechar{𝔻}{\ensuremath{\mathbb{D}}}
\newunicodechar{α}{\ensuremath{\alpha}}
\newunicodechar{β}{\ensuremath{\beta}}
\newunicodechar{γ}{\ensuremath{\gamma}}
\newunicodechar{δ}{\ensuremath{\delta}}
\newunicodechar{ε}{\ensuremath{\varepsilon}}
\newunicodechar{θ}{\ensuremath{\theta}}
\newunicodechar{λ}{\ensuremath{\lambda}}
\newunicodechar{μ}{\ensuremath{\mu}}
\newunicodechar{σ}{\ensuremath{\sigma}}
\newunicodechar{τ}{\ensuremath{\tau}}
\newunicodechar{φ}{\ensuremath{\varphi}}
\newunicodechar{ψ}{\ensuremath{\psi}}
\newunicodechar{ω}{\ensuremath{\omega}}
\newunicodechar{Σ}{\ensuremath{\Sigma}}
% majuscules produites par \MakeUppercase dans les titres (α-aminés -> Α)
\newunicodechar{Α}{\ensuremath{\alpha}}
\newunicodechar{Β}{\ensuremath{\beta}}
\newunicodechar{Γ}{\ensuremath{\Gamma}}
\newunicodechar{Δ}{\ensuremath{\Delta}}
\newunicodechar{Ε}{\ensuremath{\varepsilon}}
\newunicodechar{Θ}{\ensuremath{\Theta}}
\newunicodechar{Λ}{\ensuremath{\Lambda}}
\newunicodechar{Μ}{\ensuremath{\mu}}
\newunicodechar{Τ}{\ensuremath{\tau}}
\newunicodechar{Φ}{\ensuremath{\Phi}}
\newunicodechar{Ψ}{\ensuremath{\Psi}}
\newunicodechar{Ω}{\ensuremath{\Omega}}
\newunicodechar{↦}{\ensuremath{\mapsto}}
\newunicodechar{∈}{\ensuremath{\in}}
\newunicodechar{∉}{\ensuremath{\notin}}
\newunicodechar{∧}{\ensuremath{\wedge}}
\newunicodechar{⇔}{\ensuremath{\Leftrightarrow}}
\newunicodechar{⟺}{\ensuremath{\Longleftrightarrow}}
\newunicodechar{⇒}{\ensuremath{\Rightarrow}}
\newunicodechar{⟹}{\ensuremath{\Longrightarrow}}
\newunicodechar{∘}{\ensuremath{\circ}}
\newunicodechar{∪}{\ensuremath{\cup}}
\newunicodechar{∩}{\ensuremath{\cap}}
\newunicodechar{≡}{\ensuremath{\equiv}}
\newunicodechar{⊂}{\ensuremath{\subset}}
\newunicodechar{⊥}{\ensuremath{\perp}}
\newunicodechar{∃}{\ensuremath{\exists}}
\newunicodechar{∀}{\ensuremath{\forall}}
\newunicodechar{∛}{\ensuremath{\sqrt[3]{\,}}}
\newunicodechar{∜}{\ensuremath{\sqrt[4]{\,}}}
\newunicodechar{⃗}{\ensuremath{{}^{\to}}}
\newunicodechar{ⁿ}{\ifmmode^{n}\else\textsuperscript{\ensuremath{n}}\fi}
\newunicodechar{ˣ}{\ifmmode^{x}\else\textsuperscript{\ensuremath{x}}\fi}
\newunicodechar{ᵇ}{\ifmmode^{b}\else\textsuperscript{\ensuremath{b}}\fi}
\newunicodechar{ʳ}{\ifmmode^{r}\else\textsuperscript{\ensuremath{r}}\fi}
\newunicodechar{ᵘ}{\ifmmode^{u}\else\textsuperscript{\ensuremath{u}}\fi}
\newunicodechar{ʸ}{\ifmmode^{y}\else\textsuperscript{\ensuremath{y}}\fi}
\newunicodechar{ᵃ}{\ifmmode^{a}\else\textsuperscript{\ensuremath{a}}\fi}
\newunicodechar{ᵉ}{\ifmmode^{e}\else\textsuperscript{\ensuremath{e}}\fi}
\newunicodechar{ᵏ}{\ifmmode^{k}\else\textsuperscript{\ensuremath{k}}\fi}
\newunicodechar{ᵖ}{\ifmmode^{p}\else\textsuperscript{\ensuremath{p}}\fi}
\newunicodechar{ᴺ}{\ifmmode^{N}\else\textsuperscript{\ensuremath{N}}\fi}
\newunicodechar{ᶜ}{\ifmmode^{c}\else\textsuperscript{\ensuremath{c}}\fi}
\newunicodechar{ʰ}{\ifmmode^{h}\else\textsuperscript{\ensuremath{h}}\fi}
\newunicodechar{ᵠ}{\ifmmode^{q}\else\textsuperscript{\ensuremath{q}}\fi}
\newunicodechar{ₙ}{\ifmmode_{n}\else\textsubscript{\ensuremath{n}}\fi}
\newunicodechar{ₐ}{\ifmmode_{a}\else\textsubscript{\ensuremath{a}}\fi}
\newunicodechar{ᵢ}{\ifmmode_{i}\else\textsubscript{\ensuremath{i}}\fi}
\newunicodechar{ᵣ}{\ifmmode_{r}\else\textsubscript{\ensuremath{r}}\fi}
\newunicodechar{ₖ}{\ifmmode_{k}\else\textsubscript{\ensuremath{k}}\fi}
\newunicodechar{ᵦ}{\ifmmode_{\beta}\else\textsubscript{\ensuremath{\beta}}\fi}
\newunicodechar{ₓ}{\ifmmode_{x}\else\textsubscript{\ensuremath{x}}\fi}
\newfontfamily\popdisplay{ArchivoBlack-Regular.ttf}[Path=../../../fonts/]
\newfontfamily\poplogo{SpaceGrotesk-Bold.ttf}[Path=../../../fonts/]
\newfontfamily\popsemi{PlusJakartaSans-SemiBold.ttf}[Path=../../../fonts/]
\newfontfamily\popxbold{PlusJakartaSans-ExtraBold.ttf}[Path=../../../fonts/]
\newfontfamily\popxboldls{PlusJakartaSans-ExtraBold.ttf}[Path=../../../fonts/,LetterSpace=3.0]

\setlist[enumerate]{leftmargin=1.7em,itemsep=5pt,topsep=4pt}
\setlist[itemize]{leftmargin=1.7em,itemsep=4pt,topsep=4pt,label={\color{poporange}\textbullet}}
\setlength{\parindent}{0pt}
\setlength{\parskip}{5pt}
\renewcommand{\arraystretch}{1.25}

% pgfplots : style Pop par défaut
\pgfplotsset{
  every axis/.append style={axis line style={popink,line width=1pt},tick style={popink},
    grid style={popline,line width=0.5pt},label style={font=\small},tick label style={font=\footnotesize}},
  popcurve/.style={poporange,line width=1.8pt,no markers,smooth},
}
% Même style disponible dans un \draw TikZ simple (hors pgfplots)
\tikzset{popcurve/.style={poporange,line width=1.8pt}}
\tikzset{popgrid/.style={popline,very thin}}
\colorlet{popcurve}{poporange} % le nom du style est parfois utilisé comme couleur

% ============================================================
% Pill / squiggle
% ============================================================
\newcommand{\poppill}[3]{%
  \tikz[baseline=-0.55ex]\node[draw=popink,line width=1.2pt,rounded corners=2.6pt,%
    fill=#2,inner xsep=6.5pt,inner ysep=3pt]{%
    \popxboldls\fontsize{7}{7}\selectfont\color{#3}\MakeUppercase{#1}};%
}
\newcommand{\popsquiggle}[1]{%
  \tikz[baseline=-0.4ex]{
    \draw[#1,line width=2.6pt,line cap=round]
      plot [smooth] coordinates {(0,0.09) (0.34,0.01) (0.68,0.21) (1.02,0.01) (1.36,0.10)};
  }%
}
% Mot-clé surligné (marqueur orange sous le texte)
\newcommand{\cle}[1]{{\popxbold\color{popdark}#1}}

% ============================================================
% En-tête / pied
% ============================================================
\pagestyle{fancy}
\fancyhf{}
\renewcommand{\headrulewidth}{0pt}
\renewcommand{\footrulewidth}{0pt}
\lhead{\raisebox{-0.15ex}{\poplogo\fontsize{13}{13}\selectfont\color{popink}OnBuch\color{poporange}+}}
\rhead{\poppill{\DOCMATIERE\ \textbullet\ \DOCNIVEAU\ \textbullet\ Leçon \DOCLECON}{white}{popink}}
\lfoot{\popxbold\fontsize{7.6}{7.6}\selectfont\color{popmuted}\MakeUppercase{Cours signé OnBuch+}\ \textbullet\ {\popsemi\DOCTITRE}}
\rfoot{\tikz[baseline=-0.6ex]\node[rounded corners=8.5pt,fill=popink,minimum height=17pt,minimum width=17pt,inner xsep=5pt,inner sep=0]{\popxbold\fontsize{8}{8}\selectfont\color{popcream}\thepage\,/\,\pageref*{LastPage}};}

% ============================================================
% Titres
% ============================================================
\newcommand{\chaptitre}[2]{%
  \begin{tcolorbox}[enhanced,colback=poporange,colframe=popink,boxrule=2.6pt,arc=11pt,
      drop shadow=popink,left=15pt,right=15pt,top=12pt,bottom=11pt,before skip=3pt,after skip=18pt]
    \poppill{#2}{popink}{popcream}\par\vspace{9pt}
    {\raggedright\hyphenpenalty=10000\exhyphenpenalty=10000\popdisplay\fontsize{22}{24}\selectfont\color{popcream}\MakeUppercase{#1}\par}
  \end{tcolorbox}
}
% Section numérotée : pastille orange avec le numéro + titre Archivo
\newcounter{popsec}
\newcommand{\coursec}[1]{%
  \stepcounter{popsec}\setcounter{popsub}{0}%
  \par\vspace{16pt}\noindent
  \begin{minipage}{\linewidth}
  \tikz[baseline=-0.7ex]\node[draw=popink,line width=1.6pt,fill=poporange,rounded corners=4pt,minimum size=22pt,inner sep=0]{\popdisplay\fontsize{12}{12}\selectfont\color{popcream}\thepopsec};%
  \hspace{8pt}{\popdisplay\fontsize{15}{17}\selectfont\color{popink}\MakeUppercase{#1}}\par
  \vspace{4pt}\noindent{\color{poporange}\rule{36pt}{3.6pt}}
  \end{minipage}\par\nopagebreak\vspace{8pt}
}
\newcounter{popsub}
\newcommand{\courssub}[1]{%
  \stepcounter{popsub}%
  \par\vspace{9pt}\noindent{\popxbold\fontsize{11.5}{13}\selectfont\color{popink}{\color{poporange}\thepopsec.\thepopsub}\hspace{6pt}#1}\par\nopagebreak\vspace{4pt}
}

% ============================================================
% Boîtes pédagogiques — même recette : fond blanc, bordure épaisse colorée,
% ombre dure encre, badge plein. Titre optionnel [#1].
% ============================================================
\tcbset{popbase/.style={breakable,enhanced,colback=white,boxrule=2.3pt,arc=9pt,
  shadow={3pt}{-3pt}{0pt}{popink},left=14pt,right=14pt,top=11pt,bottom=11pt,before skip=11pt,after skip=15pt,
  fontupper=\popsemi\fontsize{10.6}{14.5}\selectfont\color{popink},
  fontlower=\popsemi\fontsize{10.6}{14.5}\selectfont\color{popink}}}
% Coche dessinée (le glyphe ✓ n'existe pas dans les polices du document)
\newcommand{\popcheck}[1]{\tikz[baseline=-0.5ex]\draw[#1,line width=1.6pt,line cap=round,line join=round](-0.13,0.0)--(-0.04,-0.09)--(0.13,0.1);}
\providecommand{\coche}{\popcheck{popgreen}}
\providecommand{\resultat}[1]{\par\smallskip\noindent\fcolorbox{popgreen}{popgreenL}{\parbox{\dimexpr\linewidth-2\fboxsep-2\fboxrule\relax}{\textbf{Résultat.} #1}}\par\smallskip}
\newcommand{\popboxhead}[3]{\poppill{#1}{#2}{white}\ifx\relax#3\relax\else\hspace{6pt}{\popxbold\fontsize{10.6}{12}\selectfont\color{popink}#3}\fi\par\vspace{7pt}}

\newtcolorbox{aretenir}[1][]{popbase,colframe=popgreen,before upper={\popboxhead{À retenir}{popgreen}{#1}}}
\newtcolorbox{exemplebox}[1][]{popbase,colframe=poporange,before upper={\popboxhead{Exemple}{poporange}{#1}}}
\newtcolorbox{definition}[1][]{popbase,colframe=popblue,before upper={\popboxhead{Définition}{popblue}{#1}}}
\newtcolorbox{propriete}[1][]{popbase,colframe=poppurple,before upper={\popboxhead{Propriété}{poppurple}{#1}}}
\newtcolorbox{methode}[1][]{popbase,colframe=popgold,before upper={\popboxhead{Méthode}{popgold}{#1}}}
\newtcolorbox{attention}[1][]{popbase,colframe=poppink,before upper={\popboxhead{Attention}{poppink}{#1}}}
\newtcolorbox{experience}[1][]{popbase,colframe=popblue,colback=popblueL,before upper={\popboxhead{Expérience}{popblue}{#1}}}
% « Le savais-tu ? » : carte violette pleine (contexte, culture, Cameroun)
\newtcolorbox{savaistu}[1][]{popbase,colback=poppurple,colframe=popink,
  fontupper=\popsemi\fontsize{10.4}{14.2}\selectfont\color{white},
  before upper={\poppill{Le savais-tu\,?}{white}{poppurple}\ifx\relax#1\relax\else\hspace{6pt}{\popxbold\color{white}#1}\fi\par\vspace{7pt}}}
% Exercice résolu : énoncé en haut, \tcblower puis la solution sur fond crème
\newcounter{popexo}\newcounter{popexr}
\newtcolorbox{exoresolu}[1][]{popbase,colframe=poporange,colbacklower=popcream,
  segmentation style={popink,line width=1.2pt,dashed},
  before upper={\stepcounter{popexr}\popboxhead{Exercice résolu \thepopexr}{poporange}{#1}},
  before lower={\poppill{Solution}{popink}{popcream}\par\vspace{6pt}}}
% Exercice (non résolu en place) — la correction est dans la partie Corrigés
\newtcolorbox{exercice}[1][]{popbase,colframe=popink,
  before upper={\stepcounter{popexo}\popboxhead{Exercice \thepopexo}{popink}{#1}}}
% Corrigé
\newtcolorbox{corrige}[1][]{popbase,colframe=popgreen,colback=popgreenL,
  before upper={\popboxhead{Corrigé}{popgreen}{#1}}}
% Objectifs / prérequis / bilan
\newtcolorbox{objectifs}{popbase,colframe=popink,colback=poporangeL,
  before upper={\popboxhead{Objectifs de la leçon}{popink}{}}}
\newtcolorbox{prerequis}{popbase,colframe=popink,colback=popblueL,
  before upper={\popboxhead{Prérequis}{popblue}{}}}
\newtcolorbox{bilan}{popbase,colframe=popink,colback=white,boxrule=2.8pt,
  before upper={\popboxhead{Fiche bilan}{poporange}{L'essentiel à connaître pour l'examen}}}
% Figure encadrée avec légende
\newtcolorbox{popfigure}[1][]{enhanced,breakable=false,colback=white,colframe=popline,boxrule=1.4pt,arc=8pt,
  left=10pt,right=10pt,top=10pt,bottom=8pt,before skip=10pt,after skip=12pt,halign=center,
  lower separated=false,halign lower=center,
  fontlower=\popsemi\fontsize{9}{11.5}\selectfont\color{popmuted},
  #1}
\newcommand{\legende}[1]{\par\vspace{6pt}{\popsemi\fontsize{9}{11.5}\selectfont\color{popmuted}#1}}

% ============================================================
% Signature OnBuch+
% ============================================================
\newcommand{\poplogoinline}[1]{{\poplogo\fontsize{#1}{#1}\selectfont\color{popink}OnBuch\color{poporange}+}}
\newcommand{\signatureonbuch}{%
  \begin{tcolorbox}[enhanced,colback=popink,colframe=popink,boxrule=2.2pt,arc=10pt,
      drop shadow=poporange,left=14pt,right=14pt,top=10pt,bottom=10pt,before skip=0pt,after skip=0pt]
    \tikz[baseline=-0.7ex]\node[circle,fill=popgreen,draw=popcream,line width=1.3pt,minimum size=22pt,inner sep=0]{\popcheck{white}};%
    \hspace{9pt}\begin{minipage}[c]{0.62\linewidth}
      {\popxbold\fontsize{12}{13}\selectfont\color{popcream}Signé\ \textbullet\ {\poplogo OnBuch\color{poporange}+}}\par\vspace{2pt}
      {\raggedright\popsemi\fontsize{8}{10.5}\selectfont\color{popline}\MakeUppercase{Équipe pédagogique OnBuch+}\\\MakeUppercase{Conforme au programme officiel MINESEC}\par}
    \end{minipage}\hfill
    {\poplogo\fontsize{22}{22}\selectfont\color{popcream}OnBuch\color{poporange}+}
  \end{tcolorbox}%
}

% ============================================================
% Page de garde
% #1 titre · #2 matière · #3 niveau · #4 module · #5 n° leçon · #6 durée ·
% #7 description / objectifs (texte ou itemize)
% ============================================================
\newcommand{\pagegarde}[7]{%
  \thispagestyle{empty}
  \newgeometry{a4paper,margin=1.7cm,top=1.6cm,bottom=1.4cm}%
  \begin{tikzpicture}[remember picture,overlay]
    \fill[poporange!18] ([xshift=-3.2cm,yshift=-3.6cm]current page.north east) circle (4.6cm);
    \fill[poppurple!14] ([xshift=2.4cm,yshift=7.6cm]current page.south west) circle (3.8cm);
  \end{tikzpicture}%
  {\poplogo\fontsize{30}{30}\selectfont\color{popink}OnBuch\color{poporange}+}\hfill
  \raisebox{0.6ex}{\poppill{Cours signé \textbullet\ Premium}{popink}{popcream}}\par
  \vspace{1pt}\popsquiggle{poppurple}\par
  \vspace{8pt}
  \begin{tcolorbox}[enhanced,colback=poporange,colframe=popink,boxrule=2.8pt,arc=13pt,
      drop shadow=popink,left=18pt,right=18pt,top=12pt,bottom=13pt,after skip=10pt]
    \poppill{#2 \textbullet\ #3}{popink}{popcream}\hspace{5pt}\poppill{#4}{white}{popink}\par\vspace{8pt}
    {\popdisplay\fontsize{14}{14}\selectfont\color{popink}LEÇON #5}\par\vspace{4pt}
    {\raggedright\hyphenpenalty=10000\exhyphenpenalty=10000\popdisplay\fontsize{25}{27}\selectfont\color{popcream}\MakeUppercase{#1}\par}
    \vspace{8pt}
    {\popxbold\fontsize{9.5}{9.5}\selectfont\color{popink}\MakeUppercase{Durée conseillée : #6}}
  \end{tcolorbox}
  \begin{tcolorbox}[enhanced,colback=white,colframe=popink,boxrule=2.2pt,arc=10pt,
      drop shadow=popink,left=16pt,right=16pt,top=10pt,bottom=10pt,after skip=8pt]
    \poppill{Dans cette leçon}{poppurple}{white}\par\vspace{5pt}
    \popsemi\fontsize{9.3}{12.6}\selectfont\color{popink}#7
  \end{tcolorbox}
  \noindent\begin{minipage}[t]{0.487\linewidth}
    \begin{tcolorbox}[enhanced,colback=popgreen,colframe=popink,boxrule=2.2pt,arc=10pt,
        drop shadow=popink,left=11pt,right=11pt,top=10pt,bottom=10pt,height=3.1cm,valign=center]
      \tcbox[colback=white,colframe=popink,boxrule=1.3pt,arc=5pt,boxsep=0pt,nobeforeafter,
        left=3pt,right=3pt,top=3pt,bottom=3pt,valign=center]{\qrcode[height=1.9cm]{https://play.google.com/store/apps/details?id=com.luvvix.onbuch&hl=fr}}%
      \hspace{8pt}\begin{minipage}[c]{0.5\linewidth}
        {\popdisplay\fontsize{11}{12.5}\selectfont\color{white}\MakeUppercase{Télécharge OnBuch}}\par\vspace{4pt}
        {\raggedright\popsemi\fontsize{8.4}{11}\selectfont\color{white}Gratuit \textbullet\ Google Play\\Tuteur IA, annales, concours, résultats\par}
      \end{minipage}
    \end{tcolorbox}
  \end{minipage}\hfill
  \begin{minipage}[t]{0.487\linewidth}
    \begin{tcolorbox}[enhanced,colback=poppurple,colframe=popink,boxrule=2.2pt,arc=10pt,
        drop shadow=popink,left=11pt,right=11pt,top=10pt,bottom=10pt,height=3.1cm,valign=center]
      \tikz[baseline=-0.7ex]\node[circle,fill=white,draw=popink,line width=1.3pt,minimum size=1.6cm,inner sep=0]{\popdisplay\fontsize{24}{24}\selectfont\color{poppurple}+};%
      \hspace{8pt}\begin{minipage}[c]{0.6\linewidth}
        {\popdisplay\fontsize{11}{12.5}\selectfont\color{white}\MakeUppercase{Passe à OnBuch+}}\par\vspace{4pt}
        {\raggedright\popsemi\fontsize{8.4}{11}\selectfont\color{white}Tous les cours signés, Tuteur IA illimité, fiches PDF — onglet OnBuch+ de l'app\par}
      \end{minipage}
    \end{tcolorbox}
  \end{minipage}
  \vspace{8pt}
  \popcontenu
  \vfill
  \signatureonbuch
  \restoregeometry
  \newpage
}

% Rangée « Ce cours contient » (page de garde)
\newcommand{\poptile}[3]{% #1 couleur #2 gros texte #3 légende
  \begin{tcolorbox}[enhanced,colback=#1,colframe=popink,boxrule=2pt,arc=9pt,shadow={2.5pt}{-2.5pt}{0pt}{popink},
      left=6pt,right=6pt,top=9pt,bottom=9pt,halign=center,valign=center,height=1.75cm,width=\linewidth,nobeforeafter]
    {\popdisplay\fontsize{12.5}{13}\selectfont\color{white}\MakeUppercase{#2}}\par\vspace{3pt}
    {\centering\popsemi\fontsize{7.8}{9.5}\selectfont\color{white}#3\par}
  \end{tcolorbox}}
\newcommand{\popcontenu}{%
  \par\noindent{\popxboldls\fontsize{7.5}{7.5}\selectfont\color{popmuted}CE COURS SIGNÉ CONTIENT}\par\vspace{7pt}
  \noindent\begin{minipage}[t]{0.235\linewidth}\poptile{popblue}{Cours}{complet, illustré, pas à pas}\end{minipage}\hfill
  \begin{minipage}[t]{0.235\linewidth}\poptile{poporange}{Méthodes}{et exercices résolus}\end{minipage}\hfill
  \begin{minipage}[t]{0.235\linewidth}\poptile{poppink}{Exercices}{type examen + corrigés}\end{minipage}\hfill
  \begin{minipage}[t]{0.235\linewidth}\poptile{popgreen}{Fiche bilan}{l'essentiel à retenir}\end{minipage}\par}

% Sommaire
\newtcolorbox{sommaire}{enhanced,colback=white,colframe=popink,boxrule=2.2pt,arc=10pt,
  drop shadow=popink,left=18pt,right=18pt,top=13pt,bottom=13pt,before skip=6pt,after skip=20pt,
  before upper={\poppill{Sommaire}{poppurple}{white}\par\vspace{9pt}\popsemi\fontsize{10.6}{14.5}\selectfont\color{popink}}}

% Signature de fin de cours — poussée tout en bas de la dernière page
\newcommand{\findecours}{%
  \par\vfill\nopagebreak
  \begin{center}\popsquiggle{poporange}\end{center}\vspace{4pt}
  \signatureonbuch
}
\AtBeginDocument{\def\llbracket{\mathopen{[\mkern-3.5mu[}}\def\rrbracket{\mathclose{]\mkern-3.5mu]}}} % intervalles d'entiers (glyphe ⟦ absent de la police)
% Glyphes absents de Fira Math : redéfinis à partir de glyphes disponibles
\AtBeginDocument{%
  \def\setminus{\mathbin{\backslash}}\let\smallsetminus\setminus
  \def\vdots{\mathord{\vbox{\baselineskip3.2pt\lineskiplimit0pt\kern4pt\hbox{.}\hbox{.}\hbox{.}}}}%
  \def\ddots{\mathinner{\mkern1mu\raise6.5pt\hbox{.}\mkern2mu\raise3.6pt\hbox{.}\mkern2mu\raise0.7pt\hbox{.}\mkern1mu}}%
  \def\triangle{\mathord{\scalebox{0.9}{$\symup{\Delta}$}}}%
  \def\nabla{\mathord{\raisebox{\depth}{\scalebox{1}[-1]{$\symup{\Delta}$}}}}%
  \def\checkmark{\popcheck{popdark}}%
  \def\star{\mathbin{\ast}}\def\bigstar{\mathord{\ast}}%
  \def\diamond{\mathbin{\scalebox{0.6}{$\Diamond$}}}%
  \def\bigcup{\mathop{\mathchoice{\raisebox{-0.3em}{\scalebox{1.7}{$\cup$}}}{\scalebox{1.25}{$\cup$}}{\cup}{\cup}}\displaylimits}%
  \def\bigcap{\mathop{\mathchoice{\raisebox{-0.3em}{\scalebox{1.7}{$\cap$}}}{\scalebox{1.25}{$\cap$}}{\cap}{\cap}}\displaylimits}%
  \def\lesseqqgtr{\mathrel{\lessgtr}}\def\bigoplus{\mathop{\oplus}\displaylimits}%
}
\providecommand{\ssi}{si et seulement si} % abréviation fréquente
\AtBeginDocument{\providecommand{\wideparen}{\overparen}} % arc de cercle

%% FIGFIX-BEGIN v4 — correctifs globaux des figures (docs/onbuch-generation/tools/figfix)
\usepackage{adjustbox}\usepackage{etoolbox}
% toute figure TikZ/pgfplots plus large que la ligne est réduite à la largeur disponible
\BeforeBeginEnvironment{tikzpicture}{\begin{adjustbox}{max width=\linewidth}}
\AfterEndEnvironment{tikzpicture}{\end{adjustbox}}
% légendes pgfplots lisibles par-dessus les courbes
\pgfplotsset{every axis/.append style={legend style={fill=white,fill opacity=0.92,text opacity=1,draw=black!15,inner sep=3pt}}}
% étiquettes d'arêtes (syntaxe edge["..."]) avec fond blanc
\tikzset{every edge quotes/.append style={fill=white,inner sep=1.2pt}}
%% FIGFIX-END
\begin{document}

\pagegarde{Lutte contre les troubles liés à la régulation du taux d'hormones sexuelles chez l'homme et chez la femme et la stérilité}{SVT}{Tle D}{Module II : L'Éducation à la Santé — La santé reproductive}{N12}{18 h}{Cette leçon aborde la régulation hormonale de la reproduction chez l'homme et la femme, les moyens de la maîtriser (contraception, contragestion), les dérèglements hormonaux et la prise en charge de la stérilité, y compris les techniques de PMA. Elle s'appuie sur l'interprétation d'expériences classiques et de courbes hormonales pour construire les savoirs et savoir-faire requis au Bac.
\vspace{4pt}
\begin{multicols}{2}\small
\begin{itemize}
\item À la fin de cette leçon, je sais expliquer le rôle de l'axe hypothalamo-hypophysaire dans la régulation de la testostérone à partir d'expériences de castration, greffe et parabiosis.
\item Je sais tracer et analyser les courbes de variation des hormones ovariennes (œstradiol, progestérone) et hypophysaires (FSH, LH) au cours du cycle menstruel.
\item Je sais réaliser un schéma fonctionnel de la rétroaction négative chez l'homme et des rétrocontrôles positif/négatif chez la femme.
\item Je sais expliquer le mécanisme d'action des pilules contraceptives (œstro-progestatives, microprogestatives) et contragestives (pilule du lendemain) et interpréter les courbes hormonales d'une utilisatrice.
\item Je sais identifier les causes principales de dérèglements hormonaux (hyper/hypogonadisme, SOPK, etc.), leurs conséquences cliniques et les traitements hormonaux adaptés.
\item Je sais interpréter un spermogramme (concentration, mobilité, morphologie) et en déduire une éventuelle cause d'infertilité masculine.
\end{itemize}\end{multicols}}

\begin{sommaire}
\begin{multicols}{2}
\begin{enumerate}
\item 1. Régulation de la testostérone chez l'homme
\item 2. Cycles sexuels et régulation hormonale chez la femme
\item 3. Maîtrise de la reproduction : contraception et contragestion
\item 4. Dérèglements des hormones sexuelles : causes, conséquences, traitement, prévention
\item 5. Stérilité et procréation médicalement assistée (PMA)
\item Activité d'intégration
\item Méthodes et exercices résolus
\item Exercices
\item Corrigés des exercices
\item Fiche bilan
\end{enumerate}
\end{multicols}
\end{sommaire}

\begin{objectifs}
\begin{itemize}
\item À la fin de cette leçon, je sais expliquer le rôle de l'axe hypothalamo-hypophysaire dans la régulation de la testostérone à partir d'expériences de castration, greffe et parabiosis.
\item Je sais tracer et analyser les courbes de variation des hormones ovariennes (œstradiol, progestérone) et hypophysaires (FSH, LH) au cours du cycle menstruel.
\item Je sais réaliser un schéma fonctionnel de la rétroaction négative chez l'homme et des rétrocontrôles positif/négatif chez la femme.
\item Je sais expliquer le mécanisme d'action des pilules contraceptives (œstro-progestatives, microprogestatives) et contragestives (pilule du lendemain) et interpréter les courbes hormonales d'une utilisatrice.
\item Je sais identifier les causes principales de dérèglements hormonaux (hyper/hypogonadisme, SOPK, etc.), leurs conséquences cliniques et les traitements hormonaux adaptés.
\item Je sais interpréter un spermogramme (concentration, mobilité, morphologie) et en déduire une éventuelle cause d'infertilité masculine.
\item Je sais décrire les principales techniques de procréation médicalement assistée (FIV, ICSI, don de gamètes) et discuter leurs avantages et limites éthiques.
\item Je sais concevoir un support de sensibilisation (affiche, dépliant) sur l'infertilité adapté au contexte camerounais.
\end{itemize}
\end{objectifs}

\begin{prerequis}
\begin{itemize}
\item Notion de système endocrinien : glandes, hormones, rétroaction (Première).
\item Structure et fonction des gonades (testicules, ovaires) et de l'hypophyse (Première).
\item Bases de la méiose et de la gamétogenèse (Première).
\item Lecture de graphiques (courbes, axes, unités) et calcul de pourcentages.
\item Vocabulaire de base en éthique biomédicale (consentement, don, anonymat).
\end{itemize}
\end{prerequis}

\newpage

\coursec{Régulation de la testostérone chez l'homme}

Au quartier Mvog-Ada, Marie, 28 ans, découvre après deux ans de mariage qu'elle ne tombe pas enceinte malgré des rapports réguliers. Son mari, Paul, 30 ans, a récemment subi une analyse de sperme qui révèle une faible mobilité des spermatozoïdes (asthénospermie) et une concentration inférieure à \SI{15}{millions/mL} (oligospermie). Le couple s'interroge : est-ce un problème hormonal, une infection ancienne, ou une cause génétique ? Le centre de santé propose un bilan hormonal complet (testostérone, LH, FSH, inhibine) et, si besoin, une prise en charge par procréation médicalement assistée. Cette situation illustre l'importance cruciale de comprendre la régulation hormonale de la fonction testiculaire pour diagnostiquer et traiter les troubles de la fertilité masculine.

\courssub{Rappel anatomique et acteurs de l'axe hypothalamo-hypophyso-gonal}

La régulation de la testostérone repose sur un axe neuroendocrine précis : l'axe \cle{hypothalamo-hypophyso-gonal (HHG)}. Trois structures dialoguent par signaux chimiques (hormones) transportés par le sang.

\begin{definition}[Axe hypothalamo-hypophyso-gonal (HHG)]
Ensemble des interactions hormonales entre l'hypothalamus (zone nerveuse), l'hypophyse antérieure (adénohypophyse) et les gonades (testicules chez l'homme), assurant la régulation de la stéroïdogenèse et de la gamétogenèse.
\end{definition}

\begin{itemize}
    \item \textbf{Hypothalamus :} Synthétise la \cle{GnRH} (Gonadotropin-Releasing Hormone), peptide libéré de façon \cle{pulsatile} (toutes les \SI{60}{\minute} à \SI{90}{\minute} chez l'homme adulte) dans le système porte hypophysaire.
    \item \textbf{Hypophyse antérieure :} Répond à la GnRH par la sécrétion de deux gonadotrophines :
        \begin{itemize}
            \item \cle{LH} (Luteinizing Hormone) : cible les \cle{cellules de Leydig} (interstitielles) $\rightarrow$ stimule la synthèse de testostérone.
            \item \cle{FSH} (Follicle-Stimulating Hormone) : cible les \cle{cellules de Sertoli} (dans les tubules séminifères) $\rightarrow$ soutient la spermatogenèse et sécrète l'inhibine.
        \end{itemize}
    \item \textbf{Testicules :}
        \begin{itemize}
            \item Cellules de Leydig : transforment le cholestérol en \cle{testostérone} (stéroïde) sous l'effet de la LH.
            \item Cellules de Sertoli : nourrissent les cellules germinales, sécrètent l'\cle{inhibine} (protéine) et la protéine ABP (Androgen Binding Protein) qui concentre la testostérone dans le tubule.
        \end{itemize}
\end{itemize}

\begin{popfigure}
\begin{tikzpicture}[
    node distance=1.2cm and 1.5cm,
    box/.style={rectangle, draw=popdark, thick, fill=popcream, rounded corners, minimum width=2.5cm, minimum height=1.2cm, align=center},
    arrow/.style={-Stealth, thick},
    inhib/.style={-Circle, thick, poporange},
    plus/.style={font=\bfseries\large, popgreen},
    minus/.style={font=\bfseries\large, poporange}
]
% Nodes
\node[box, align=center] (hypo){Hypothalamus\\ \textbf{GnRH}};
\node[box, right=of hypo, align=center] (hippo){Hypophyse antérieure\\ \textbf{LH \quad FSH}};
\node[box, below right=of hippo, align=center] (leydig){Cellules de Leydig\\ \textbf{Testostérone}};
\node[box, below left=of hippo, align=center] (sertoli){Cellules de Sertoli\\ \textbf{Inhibine, ABP}};
\node[box, below=1.5cm of leydig, align=center] (target){Cibles :\\ Spermatogenèse, caractères sexuels, métabolisme};

% Arrows Stimulation
\draw[arrow, popblue] (hypo.east) -- node[above, plus] {+} (hippo.west);
\draw[arrow, popblue] (hippo.south east) -- node[right, plus] {LH +} (leydig.north east);
\draw[arrow, popblue] (hippo.south west) -- node[left, plus] {FSH +} (sertoli.north west);
\draw[arrow, popblue] (leydig.south) -- node[right, plus] {T +} (target.north east);
\draw[arrow, popblue] (sertoli.south) -- node[left, plus] {Inhibine/ABP +} (target.north west);

% Feedback Testosterone on Hypophysis and Hypothalamus
\draw[inhib, bend right=20] (leydig.north west) to node[above left, minus] {-} (hippo.south west);
\draw[inhib, bend left=20] (leydig.north) to node[above, minus] {-} (hypo.south east);

% Feedback Inhibin on Hypophysis (specific FSH)
\draw[inhib, bend left=15] (sertoli.north east) to node[above right, minus] {-} (hippo.south east);
\end{tikzpicture}
\legende{Schéma fonctionnel de l'axe hypothalamo-hypophyso-gonal chez l'homme. Flèches pleines \textcolor{popblue}{$\blacktriangleright$} = stimulation. Flèches cerclées \textcolor{poporange}{$\otimes$} = inhibition (rétroaction négative). La testostérone inhibe la LH et la GnRH. L'inhibine inhibe spécifiquement la FSH.}
\end{popfigure}

\courssub{Mise en évidence expérimentale de la régulation (Démarche d'investigation)}

La compréhension de cet axe repose sur des expériences classiques de physiologie expérimentale (rongeurs, mammifères) transposables à l'homme. Nous les analysons selon la démarche : \textbf{Observation $\rightarrow$ Hypothèse $\rightarrow$ Expérience $\rightarrow$ Interprétation $\rightarrow$ Conclusion}.

\begin{experience}[Castration (ablation des testicules) chez le rat mâle adulte]
\textbf{Matériel :} Rats mâles adultes, groupes témoin (opérés sham) et castrés. Dosage sanguin de la testostérone, LH, FSH à J0, J1, J7, J21.
\textbf{Protocole :} Ablation chirurgicale des deux testicules sous anesthésie. Prélèvements sanguins sériels.
\textbf{Observations (résultats types) :}
\begin{itemize}
    \item \textbf{Testostérone :} Chute brutale dès 24h, devient indétectable ($< \SI{0,5}{nmol/L}$) à J7.
    \item \textbf{LH et FSH :} Augmentation progressive et significative (× 5 à × 10) à partir de J3-J7, se stabilisant à un plateau élevé.
\end{itemize}
\textbf{Interprétation :} La disparition de la source de testostérone lève l'inhibition exercée sur l'hypophyse et l'hypothalamus. L'élévation des gonadotrophines prouve qu'elles étaient maintenues à un niveau bas par la testostérone.
\textbf{Conclusion :} La testostérone exerce une \cle{rétroaction négative} sur la sécrétion de LH et FSH (et indirectement sur GnRH).
\end{experience}

\begin{experience}[Injection d'extraits testiculaires (riches en testostérone) chez le rat castré]
\textbf{Protocole :} Rats castrés depuis 3 semaines (taux LH/FSH hauts, T nul). Injection quotidienne d'extrait testiculaire (ou testostérone pure) pendant 7 jours. Dosages hormonaux.
\textbf{Observations :}
\begin{itemize}
    \item \textbf{Testostérone :} Remonte à des valeurs physiologiques ($\approx \SI{15}{nmol/L}$).
    \item \textbf{LH et FSH :} Chutent brutalement pour revenir aux valeurs du témoin intact (niveaux bas).
\end{itemize}
\textbf{Interprétation :} L'apport exogène de testostérone restaure le signal inhibiteur sur l'axe HHG.
\textbf{Conclusion :} La testostérone (androgène) est l'hormone effectrice de la rétroaction négative sur la LH et la FSH. L'inhibine (protéine) contenue dans l'extrait brut participe aussi à la baisse spécifique de la FSH.
\end{experience}

\begin{experience}[Ablation et greffe d'hypophyse / Section de la tige pituitaire]
\begin{itemize}
    \item \textbf{Hypophysectomie :} Ablation de l'hypophyse $\rightarrow$ atrophie testiculaire rapide, chute de la testostérone à zéro, arrêt de la spermatogenèse. \textbf{Preuve :} L'hypophyse est \emph{indispensable} à la fonction testiculaire (source de LH/FSH).
    \item \textbf{Greffe d'hypophyse (sous capsule rénale) :} Restauration de la fonction testiculaire (vascularisation par le receveur). \textbf{Preuve :} L'action est humorale (hormones dans le sang), non nerveuse.
    \item \textbf{Section de la tige pituitaire (stalk) :} Interrompt le flux sanguin porte hypothalamo-hypophysaire. $\rightarrow$ Chute de LH/FSH et atrophie testiculaire. \textbf{Preuve :} La GnRH hypothalamique rejoint l'hypophyse \emph{exclusivement par le sang porte} (voie humorale), pas par innervation directe.
\end{itemize}
\end{experience}

\begin{experience}[Parabiose (souris mâle castré + souris mâle intacte)]
\textbf{Protocole :} Chirurgie de jonction vasculaire (anastomose) entre un mâle castré (taux de testostérone nul, LH/FSH élevés) et un mâle intact (taux de testostérone normal, LH/FSH bas). Circulation sanguine commune.
\textbf{Observation :} Chez le mâle castré, le taux de LH et FSH chute rapidement pour rejoindre des valeurs normales ; l'hypophyse réduit sa taille (perte de l'hypertrophie post-castration).
\textbf{Interprétation :} La testostérone sécrétée par le testicule du mâle intact passe dans la circulation commune, atteint l'hypophyse et l'hypothalamus du mâle castré et rétablit la rétroaction négative.
\textbf{Conclusion :} Confirmation du caractère \cle{humoral} (transport sanguin) de la régulation hormonale et de l'effet inhibiteur de la testostérone sur les gonadotrophines.
\end{experience}

\begin{aretenir}[Synthèse des preuves expérimentales]
\begin{enumerate}
    \item \textbf{Castration :} Preuve de l'existence d'une rétroaction négative (disparition hormone $\rightarrow$ hausse gonadotrophines).
    \item \textbf{Injection extrait testiculaire :} Preuve que la testostérone (et inhibine) est le médiateur de cette rétroaction.
    \item \textbf{Hypophysectomie/Greffe :} Preuve du rôle indispensable et humorale de l'hypophyse (LH/FSH).
    \item \textbf{Section tige pituitaire :} Preuve de la voie porte sanguine pour la GnRH.
    \item \textbf{Parabiose :} Preuve du transport sanguin des hormones régulatrices (testostérone).
\end{enumerate}
\end{aretenir}

\courssub{Modélisation quantitative : Courbes de variations hormonales}

L'analyse de courbes est une compétence clé au Baccalauréat. Voici la modélisation des expériences ci-dessus.

\begin{methode}[Tracer et analyser les courbes de rétroaction négative]
\begin{enumerate}
    \item \textbf{Axes :} Abscisse = Temps (jours). Ordonnées = Concentration (nmol/L pour T ; UI/L pour LH/FSH). Échelle logarithmique souvent utile pour LH/FSH.
    \item \textbf{Repères :} Marquer le jour de la castration (flèche $\downarrow$), le jour de début d'injection (flèche $\uparrow$).
    \item \textbf{Courbe Testostérone :} Plateau haut $\rightarrow$ chute brutale (castration) $\rightarrow$ plateau bas (zéro) $\rightarrow$ remontée progressive (injection) $\rightarrow$ plateau physiologique.
    \item \textbf{Courbes LH / FSH :} Plateau bas $\rightarrow$ montée lente (latence 2-3 j) vers plateau haut (castration) $\rightarrow$ chute rapide (injection) $\rightarrow$ retour plateau bas.
    \item \textbf{Relation inverse :} Souligner le miroir entre T et LH/FSH (rétroaction négative).
\end{enumerate}
\end{methode}

\begin{popfigure}
\begin{tikzpicture}
\begin{axis}[
    width=\linewidth, height=10cm,
    xlabel={Temps (jours)}, ylabel={Concentration (échelle relative)},
    xmin=0, xmax=35, ymin=0, ymax=1.2,
    xtick={0,7,14,21,28,35},
    xticklabels={J0,{J7 (Castration)},J14,{J21 (Inj. T)},J28,J35},
    legend style={at={(0.5,1.02)}, anchor=south, legend columns=3, font=\small},
    grid=major, grid style={popline},
    axis lines=left,
    tick align=outside,
    enlargelimits=false
]
% Testosterone
\addplot[popcurve, popblue, very thick, domain=0:7, samples=50] {1};
\addplot[popcurve, popblue, very thick, domain=7:21, samples=50] {0.05};
\addplot[popcurve, popblue, very thick, domain=21:35, samples=100] {0.05 + 0.95*(1-exp(-(x-21)/3))};
% LH
\addplot[popcurve, poporange, very thick, dashed, domain=0:7, samples=50] {0.2};
\addplot[popcurve, poporange, very thick, dashed, domain=7:21, samples=100] {0.2 + 0.8*(1-exp(-(x-7)/4))};
\addplot[popcurve, poporange, very thick, dashed, domain=21:35, samples=100] {1 - 0.8*exp(-(x-21)/2)};
% FSH
\addplot[popcurve, popgreen, very thick, dotted, domain=0:7, samples=50] {0.15};
\addplot[popcurve, popgreen, very thick, dotted, domain=7:21, samples=100] {0.15 + 0.7*(1-exp(-(x-7)/5))};
\addplot[popcurve, popgreen, very thick, dotted, domain=21:35, samples=100] {0.85 - 0.7*exp(-(x-21)/3)};

% Vertical lines for events
\draw[popred, thick, dashed] (axis cs:7,0) -- (axis cs:7,1.2) node[above, font=\tiny, popred] {Castration};
\draw[poppurple, thick, dashed] (axis cs:21,0) -- (axis cs:21,1.2) node[above, font=\tiny, poppurple] {Début injection T};
\legend{Testostérone, LH, FSH}
\end{axis}
\end{tikzpicture}
\legende{Évolution des taux de testostérone (T), LH et FSH avant/après castration (J7) et après injection d'extrait testiculaire (J21). Notez la relation inverse caractéristique de la rétroaction négative.}
\end{popfigure}

\begin{exemplebox}[Valeurs de référence chez l'homme adulte (Exemple chiffré)]
\begin{center}
\begin{tabularx}{\linewidth}{|l|X|X|}\hline
\rowcolor{popblueL}
\textbf{Hormone} & \textbf{Intervalle de référence (sanguin)} & \textbf{Unité} \\ \hline
Testostérone totale & 10 -- 35 & nmol/L (ou 300 -- 1000 ng/dL) \\ \hline
LH & 1,5 -- 9 & UI/L \\ \hline
FSH & 1,5 -- 12 & UI/L \\ \hline
Inhibine B & 80 -- 250 & pg/mL \\ \hline
SHBG (Globuline liant les hormones sexuelles) & 10 -- 50 & nmol/L \\ \hline
\end{tabularx}
\end{center}
\emph{Note :} Les valeurs varient selon l'âge, l'heure du prélèvement (pic matinal pour T et LH) et le laboratoire. La testostérone \emph{biodisponible} (libre + liée à l'albumine) est plus pertinente cliniquement que la testostérone totale si SHBG est anormal.
\end{exemplebox}

\courssub{Spécificité des rétroactions : Testostérone vs Inhibine}

\begin{attention}[Piège fréquent : Confondre les cibles de l'inhibition]
\begin{itemize}
    \item \textbf{Testostérone (stéroïde) :} Diffuse dans le sang, traverse la barrière hémato-encéphalique. Inhibe \textbf{GnRH (hypothalamus)} et \textbf{LH (hypophyse)}. Action majeure sur l'axe.
    \item \textbf{Inhibine (protéine, cellules de Sertoli) :} N'atteint pas l'hypothalamus facilement. Inhibe \textbf{spécifiquement la FSH} au niveau hypophysaire (action paracrine/endocrine sélective).
    \item \textbf{Conséquence clinique :} Chez l'homme âgé ou en cas de lésion tubulaire (Sertoli), l'inhibine chute $\rightarrow$ FSH augmente \emph{seule} (FSH élevée, LH normale, T normale). C'est un marqueur précoce d'atteinte tubulaire.
\end{itemize}
\end{attention}

\courssub{Complément Bac : La pulsatilité de la GnRH et de la LH}

\begin{savaistu}[Pourquoi la pulsatilité change tout en thérapeutique]
La GnRH n'est pas sécrétée en continu mais par \textbf{pulsations} (environ 1 pulse / \SI{60}{\minute}--\SI{90}{\minute} chez l'homme). Cette pulsatilité est \emph{indispensable} :
\begin{itemize}
    \item \textbf{Stimulation continue (agonistes GnRH à action prolongée, ex : Triptoréline, Leuproréline) :} $\rightarrow$ \textbf{Down-récepteurs} (désensibilisation) $\rightarrow$ \textbf{Chute brutale de LH/FSH puis de Testostérone} (castration chimique). Utilisé dans le cancer de la prostate hormono-dépendant, l'endométriose, la puberté précoce, ou en PMA (protocole long FIV).
    \item \textbf{Stimulation pulsatile (pompe à GnRH) :} Mime la physiologie $\rightarrow$ Maintien/restauration de la sécrétion gonadotrope physiologique. Utilisé pour traiter l'hypogonadisme hypogonadotrope (ex : syndrome de Kallmann) pour induire la spermatogenèse.
\end{itemize}
\textbf{Au Cameroun :} Les protocoles de stimulation ovarienne en FIV (Fécondation In Vitro) au Centre Mère-Enfant de la Fondation Chantal Biya ou aux cliniques privées de Douala/Yaoundé utilisent systématiquement des agonistes (ou antagonistes) de la GnRH pour contrôler la pulsatilité et éviter l'ovulation prématurée.
\end{savaistu}

\begin{exoresolu}[Interprétation d'un bilan hormonal masculin]
\textbf{Énoncé :} Paul, 30 ans (le mari de Marie), consulte pour infertilité. Son spermogramme montre : concentration \SI{8}{millions/mL} (norme OMS 2021 $> \SI{16}{millions/mL}$), mobilité progressive \SI{25}{\percent} (norme $> \SI{30}{\percent}$). Le bilan hormonal donne : Testostérone = \SI{8}{nmol/L} (bas), LH = \SI{12}{UI/L} (élevé), FSH = \SI{25}{UI/L} (très élevé), Inhibine B = \SI{30}{pg/mL} (bas).
\textbf{Questions :}
\begin{enumerate}
    \item Caractériser le profil hormonal (hypo/hyper-gonadotrophie).
    \item Identifier le niveau de la lésion (hypothalamo-hypophysaire vs testiculaire).
    \item Expliquer la dissociation LH/FSH (FSH $>$ LH).
    \item Proposer une étiologie probable.
\end{enumerate}

\tcblower
\textbf{Solution détaillée :}
\begin{enumerate}
    \item \textbf{Profil :} \cle{Hypogonadisme hypergonadotrope} (ou testiculaire primitif). T bas, LH/FSH hauts. La gonade ne répond pas à la stimulation hypophysaire, qui compense par hypersecrétion (perte de rétroaction négative).
    \item \textbf{Niveau lésion :} \textbf{Testiculaire.} L'hypophyse fonctionne (elle sécrète LH/FSH en excès), mais la cible (testicule) est défaillante. Si lésion hypophysaire : T bas, LH/FSH bas ou inappropriément normaux (hypogonadisme hypogonadotrope).
    \item \textbf{Dissociation FSH $>$ LH :} La FSH est spécifiquement inhibée par l'\textbf{Inhibine B} (cellules de Sertoli). L'inhibine B est très basse (\SI{30}{pg/mL}) $\rightarrow$ levée sélective de l'inhibition sur FSH. La LH est inhibée par la Testostérone (qui est basse mais pas nulle) et par l'estradiol (issu de l'aromatisation de T). L'atteinte des tubules (Sertoli) est donc plus sévère que l'atteinte interstitielle (Leydig).
    \item \textbf{Étiologies probables :}
        \begin{itemize}
            \item \textbf{Génétique :} Syndrome de Klinefelter (47,XXY) : cause n°1 d'infertilité masculine non obstructive. Phénotype : taille grande, gynécomastie, petits testicules fermes. Caryotype indispensable.
            \item \textbf{Acquise :} Orchite post-oreillons (bilatérale), radiothérapie/chimiothérapie, varicocèle bilatérale sévère non traitée, toxiques (alcool, pesticides agricoles), cryptorchidie bilatérale non opérée à temps.
        \end{itemize}
\end{enumerate}
\textbf{Conclusion pour le couple :} Si caryotype normal, une biopsie testiculaire (TESE) peut rechercher des spermatozoïdes pour ICSI (Injection IntraCytoplasmique de Spermatozoïde). Si Klinefelter : ICSI possible si foyers de spermatogenèse résiduels, avec conseil génétique (risque transmission).
\end{exoresolu}

\begin{exercice}[Analyse de cas clinique : Hypogonadisme secondaire]
M. K., 25 ans, consulte pour absence de caractères sexuels secondaires (peau imberbe, voix aiguë, testicules \SI{4}{mL}), fatigue, absence de libido. Bilan : Testostérone = \SI{2}{nmol/L} ; LH = \SI{0,5}{UI/L} ; FSH = \SI{0,8}{UI/L} ; Prolactine = \SI{1200}{mUI/L} (N $< \SI{400}{mUI/L}$). IRM hypophyse : micro-adénome de \SI{6}{mm}.
\begin{enumerate}
    \item Classer le type d'hypogonadisme.
    \item Expliquer le mécanisme physiopathologique reliant l'adénome prolactinome à l'hypogonadisme.
    \item Citer le traitement de première intention et son effet attendu sur l'axe HHG.
\end{enumerate}
\end{exercice}

\begin{corrige}[Exercice n°1]
\begin{enumerate}
    \item \textbf{Hypogonadisme hypogonadotrope (secondaire / central).} T bas, LH/FSH bas (inappropriément bas par rapport au T bas). L'hypophyse ne stimule pas la gonade.
    \item \textbf{Mécanisme :} L'adénome sécrète de la \textbf{Prolactine (PRL)} en excès (hyperprolactinémie). La PRL stimule la sécrétion de dopamine par les neurones TIDA de l'hypothalamus, ce qui inhibe la sécrétion pulsatile de \textbf{GnRH}. $\rightarrow$ Baisse GnRH $\rightarrow$ Baisse LH/FSH $\rightarrow$ Atrophie testiculaire $\rightarrow$ Baisse Testostérone.
    \item \textbf{Traitement :} \textbf{Agonistes dopaminergiques} (Bromocriptine, Cabergoline). Action : Stimulent récepteurs D2 sur lactotrophes $\rightarrow$ Inhibent sécrétion PRL $\rightarrow$ Réduction taille adénome $\rightarrow$ Levée inhibition sur GnRH $\rightarrow$ Restauration pulsatilité LH/FSH $\rightarrow$ Remontée Testostérone $\rightarrow$ Virilisation, fertilité retrouvée. Chirurgie (transsphénoïdale) si résistance ou macro-adénome compressif.
\end{enumerate}
\end{corrige}

\coursec{Cycles sexuels et régulation hormonale chez la femme}

Dans la section précédente, nous avons vu que chez l'homme, la sécrétion de testostérone est maintenue à un taux relativement \emph{constant} grâce à un rétrocontrôle négatif exercé sur le complexe hypothalamo-hypophysaire : l'axe fonctionne comme un thermostat, en régime continu. Chez la femme, la situation est radicalement différente : les sécrétions hormonales sont \emph{cycliques}. Une jeune fille de Douala qui consulte au centre de santé parce que ses règles surviennent « tous les 28 jours environ » illustre cette périodicité : l'ovaire et l'utérus accomplissent des cycles coordonnés, pilotés par un dialogue hormonal entre l'hypothalamus, l'hypophyse et les ovaires. Comprendre ce dialogue, c'est comprendre l'ovulation, la fécondité, mais aussi la contraception et certaines stérilités que nous étudierons dans les sections suivantes.

\courssub{Le cycle sexuel féminin : double cycle ovarien et utérin}

\begin{definition}[Cycle sexuel féminin]
Le \cle{cycle sexuel} de la femme est l'ensemble des modifications périodiques de l'appareil génital, d'une durée moyenne de \textbf{28 jours} (variant normalement de 21 à 35 jours), qui se répètent de la puberté à la ménopause, sauf pendant la grossesse. Par convention, le \cle{jour 1} du cycle est le premier jour des règles. Le cycle comporte deux volets imbriqués : le \cle{cycle ovarien} (dans les ovaires) et le \cle{cycle utérin} (dans l'endomètre).
\end{definition}

\textbf{Le cycle ovarien} comprend trois temps :

\begin{itemize}
\item La \cle{phase folliculaire} (jours 1 à 14 environ) : sous l'action de la FSH, une cohorte de follicules ovariens se développe ; l'un d'eux devient le \cle{follicule dominant} (follicule de De Graaf), qui sécrète de plus en plus d'œstradiol.
\item L'\cle{ovulation} (jour 14 environ) : le follicule mature se rompt et libère l'ovocyte II, capté par le pavillon de la trompe.
\item La \cle{phase lutéale} (jours 15 à 28) : les restes du follicule se transforment en \cle{corps jaune} (corps lutéum), qui sécrète progestérone et œstradiol. En l'absence de fécondation, le corps jaune dégénère (corps blanc) vers le jour 26.
\end{itemize}

\textbf{Le cycle utérin} en est le reflet dans la muqueuse utérine (endomètre) :

\begin{itemize}
\item La \cle{phase menstruelle} (jours 1 à 5) : la chute des hormones ovariennes provoque la desquamation de l'endomètre : ce sont les \cle{règles} (perte de 30 à 80 mL de sang).
\item La \cle{phase proliférative} (jours 5 à 14) : sous l'action de l'œstradiol, l'endomètre se régénère et s'épaissit (de 0,5 mm à 3--5 mm).
\item La \cle{phase sécrétoire} (jours 15 à 28) : sous l'action de la progestérone, l'endomètre devient glandulaire et vascularisé, prêt à accueillir un éventuel embryon.
\end{itemize}

\begin{aretenir}[Correspondance des phases]
Phase folliculaire ovarienne $\leftrightarrow$ phases menstruelle puis proliférative utérines ; phase lutéale ovarienne $\leftrightarrow$ phase sécrétoire utérine. L'ovaire commande, l'utérus exécute : ce sont les hormones ovariennes qui orchestrent le cycle utérin.
\end{aretenir}

\courssub{Les hormones en jeu : ovariennes et hypophysaires}

Quatre acteurs principaux dialoguent en permanence :

\begin{itemize}
\item \textbf{Hormones hypophysaires (gonadotrophines)} : la \cle{FSH} (hormone folliculo-stimulante) déclenche le recrutement et la croissance des follicules ; la \cle{LH} (hormone lutéinisante) déclenche l'ovulation par son pic et transforme le follicule rompu en corps jaune (lutéinisation).
\item \textbf{Hormones ovariennes} : l'\cle{œstradiol}, sécrété par le follicule dominant, stimule la prolifération de l'endomètre ; la \cle{progestérone}, sécrétée par le corps jaune, transforme l'endomètre en muqueuse de nidation. S'y ajoutent l'\cle{inhibine B}, sécrétée par les follicules en croissance, qui inhibe sélectivement la FSH, et l'\cle{AMH} (hormone anti-müllérienne), sécrétée par les petits follicules antraux.
\end{itemize}

L'hypothalamus, comme chez l'homme, libère la \cle{GnRH} par pulses, qui commande la sécrétion de FSH et de LH par l'antéhypophyse.

\begin{savaistu}[L'AMH, témoin fidèle de la réserve ovarienne]
L'AMH est sécrétée par les follicules antraux et pré-antraux. Sa concentration sanguine \emph{ne varie pratiquement pas au cours du cycle}, contrairement à l'œstradiol ou à la progestérone. Elle sert donc de \textbf{marqueur de la réserve ovarienne} : un dosage d'AMH faible (par exemple $< 5$ pmol/L) indique un stock de follicules appauvri. Dans les centres de procréation médicalement assistée, ce dosage aide à prédire la réponse aux traitements de stimulation ovarienne.
\end{savaistu}

\courssub{Mise en évidence expérimentale du contrôle hormonal du cycle}

\begin{experience}[Bilan des expériences classiques : ablation, greffe, injection, section, parabiose]
On utilise des mammifères modèles (rate, lapin) pour dissocier les rôles de l'hypothalamus, de l'hypophyse et des ovaires.

\begin{enumerate}
\item \textbf{Ovariectomie (ablation des ovaires)} : disparition des cycles, atrophie utérine, \textbf{hausse marquée de la FSH et de la LH}. $\rightarrow$ Les ovaires freinent l'hypophyse (rétrocontrôle négatif).
\item \textbf{Greffe d'ovaire} chez une femelle ovariectomisée (même sans réinnervation) : reprise des cycles, développement utérin, \textbf{baisse de la FSH et de la LH} vers la normale. $\rightarrow$ Le frein est hormonal (œstradiol, progestérone, inhibine), non nerveux.
\item \textbf{Injection d'extraits ovariens} (riches en œstradiol/progestérone) chez une femelle intacte ou ovariectomisée : \textbf{baisse de la FSH et de la LH}. $\rightarrow$ Confirmation du rétrocontrôle négatif par les hormones ovariennes.
\item \textbf{Hypophysectomie (ablation de l'hypophyse)} : arrêt des cycles, atrophie ovarienne et utérine. $\rightarrow$ L'hypophyse est indispensable à l'activité ovarienne.
\item \textbf{Section de la tige pituitaire} (liaison hypothalamus-hypophyse) : arrêt des cycles, atrophie des gonades, chute des gonadotrophines. $\rightarrow$ L'hypothalamus commande l'hypophyse via un facteur sanguin (la GnRH).
\item \textbf{Parabiose} (jonction vasculaire de deux femelles, l'une intacte, l'autre ovariectomisée) : les cycles de l'intacte s'imposent à l'ovariectomisée (dont l'utérus se développe, FSH/LH baissent). $\rightarrow$ Preuve formelle que le signal de rétrocontrôle circule dans le \textbf{sang} (facteur humoral).
\end{enumerate}

\textbf{Synthèse} : L'axe hypothalamo-hypophysaire est stimulé par la GnRH (hypothalamus) et freiné par les hormones ovariennes (rétrocontrôle négatif). L'expérience de parabiose écarte tout rôle nerveux direct dans ce rétrocontrôle.
\end{experience}

\begin{attention}[Ne pas conclure trop vite]
Après ovariectomie, la montée de FSH/LH ne prouve pas que les ovaires « commandent » le cycle : elle prouve seulement qu'ils \emph{freinent} l'hypophyse. C'est l'ensemble des expériences (ablation, greffe, injections d'extraits ovariens qui font chuter les gonadotrophines, section de la tige, parabiose) qui établit le rétrocontrôle négatif \emph{humoral}. Au Bac, distingue toujours observation, interprétation et conclusion.
\end{attention}

\courssub{Analyse des courbes hormonales sur 28 jours}

\begin{popfigure}
\begin{center}
\begin{tikzpicture}
\begin{axis}[
width=0.98\linewidth, height=8.2cm,
xlabel={Temps (jours)}, ylabel={Taux plasmatique (u.a.)},
xmin=0, xmax=28, ymin=0, ymax=100,
xtick={0,7,14,21,28}, ytick=\empty,
legend style={at={(0.02,0.97)}, anchor=north west, font=\small, draw=popline},
grid=major, grid style={popline!40},
axis line style={popdark}
]
\addplot[popcurve, popgreen, very thick, domain=0:28, samples=200]
{38*exp(-((x-13.5)^2)/1.2)+14*exp(-((x-6)^2)/60)+6};
\addlegendentry{Œstradiol}
\addplot[popcurve, poporange, very thick, domain=0:28, samples=200]
{4+70*exp(-((x-21)^2)/18)*(x>13.5)};
\addlegendentry{Progestérone}
\addplot[popcurve, popblue, very thick, domain=0:28, samples=200]
{12+78*exp(-((x-13.5)^2)/0.8)};
\addlegendentry{LH}
\addplot[popcurve, poppurple, very thick, domain=0:28, samples=200]
{10+14*exp(-((x-5)^2)/40)+16*exp(-((x-13.5)^2)/2.5)};
\addlegendentry{FSH}
\draw[popdark, dashed] (axis cs:14,0) -- (axis cs:14,100);
\node[popdark, font=\small, anchor=south west] at (axis cs:14.3,88) {Ovulation};
\end{axis}
\end{tikzpicture}
\end{center}
\legende{Variations des taux plasmatiques de FSH, LH, œstradiol et progestérone au cours d'un cycle de 28 jours (unités arbitraires). Le pic de LH (jour 13--14) précède l'ovulation ; la progestérone ne s'élève qu'en phase lutéale ; la chute des hormones ovariennes en fin de cycle déclenche les règles.}
\end{popfigure}

L'examen de ces courbes révèle une chronologie rigoureuse :

\begin{enumerate}
\item \textbf{Phase folliculaire précoce (jours 1--7)} : la FSH, légèrement élevée en début de cycle, recrute les follicules. Ceux-ci sécrètent de l'œstradiol et de l'inhibine B, qui exercent un \cle{rétrocontrôle négatif} : la FSH diminue, ce qui assure la sélection du seul follicule dominant.
\item \textbf{Fin de phase folliculaire (jours 8--13)} : l'œstradiol croît de façon exponentielle. Lorsqu'il dépasse un \cle{seuil critique} (environ 800 pmol/L) pendant plus de 48 h, le rétrocontrôle \emph{bascule} : il devient \textbf{positif}.
\item \textbf{Pic de LH (jour 13--14)} : l'hypophyse libère brutalement une grande quantité de LH (et, dans une moindre mesure, de FSH). C'est le déclencheur de l'ovulation.
\item \textbf{Phase lutéale (jours 15--26)} : le corps jaune sécrète progestérone (et œstradiol) ; l'ensemble exerce à nouveau un \cle{rétrocontrôle négatif} puissant : FSH et LH restent basses, aucun nouveau follicule ne peut être recruté.
\item \textbf{Chute prémenstruelle (jours 26--28)} : faute de fécondation, le corps jaune régresse ; progestérone et œstradiol s'effondrent. Ce retrait hormonal provoque le spasme des artères spirales de l'endomètre, sa nécrose et les règles. La levée du frein permet à la FSH de remonter : un nouveau cycle commence.
\end{enumerate}

\begin{definition}[Rétrocontrôle positif]
Un \cle{rétrocontrôle positif} est un mécanisme par lequel une hormone \emph{stimule indirectement sa propre sécrétion} : ici, l'œstradiol à concentration élevée et prolongée stimule la libération de GnRH et la sensibilité de l'hypophyse, ce qui déclenche le pic de LH ; or la LH favorise la maturation finale du follicule et l'ovulation, marquant la fin de la sécrétion massive d'œstradiol par le follicule dominant. C'est l'unique exemple majeur de rétroaction positive en endocrinologie humaine.
\end{definition}

\begin{propriete}[Le pic de LH, signal de l'ovulation]
Le pic de LH dure \textbf{24 à 36 heures} et \textbf{précède l'ovulation de 10 à 12 heures}. C'est donc le meilleur marqueur prédictif de l'ovulation : les tests d'ovulation vendus en pharmacie détectent précisément ce pic dans les urines.
\end{propriete}

\begin{exemplebox}[Ordres de grandeur à connaître]
Pic de LH : 30 à 100 UI/L (contre 5 à 15 UI/L en début de cycle) ; œstradiol au pic pré-ovulatoire : 800 à 1500 pmol/L ; progestérone en phase lutéale : 30 à 100 nmol/L (contre moins de 5 nmol/L en phase folliculaire).
\end{exemplebox}

\begin{methode}[Analyser une courbe hormonale de cycle]
\begin{enumerate}
\item Repérer le \textbf{jour 1} : début des règles, toutes les hormones sont basses.
\item Localiser le \textbf{pic de LH} : il est brutal, étroit, vers le jour 13--14 ; il annonce l'ovulation 10 à 12 h plus tard.
\item Identifier la \textbf{montée de progestérone} : elle ne commence qu'après le pic de LH et culmine vers le jour 21 ; elle signe l'existence d'un corps jaune, donc d'une ovulation.
\item Repérer la \textbf{chute prémenstruelle} (jours 26--28) : effondrement conjoint de la progestérone et de l'œstradiol, annonçant les règles.
\item Conclure : un cycle \emph{ovulatoire} montre un pic de LH suivi d'une phase lutéale progestative ; un cycle sans pic de LH ni progestérone est un \emph{cycle anovulatoire}.
\end{enumerate}
\end{methode}

\begin{attention}[Piège classique du Bac]
La progestérone \textbf{n'augmente qu'après l'ovulation} : elle est le produit du corps jaune. Une progestérone élevée en phase folliculaire est anormale et signe une pathologie (par exemple un kyste lutéinisé) ou une erreur de datation du cycle. Ne jamais écrire que « la progestérone prépare l'ovulation » : c'est l'œstradiol qui, par rétrocontrôle positif, la déclenche.
\end{attention}

\courssub{Schéma fonctionnel des rétrocontrôles sur l'axe hypothalamo-hypophyso-ovarien}

\begin{popfigure}
\begin{center}
\begin{tikzpicture}[
box/.style={draw=popdark, thick, rounded corners, fill=popcream, minimum width=3.2cm, minimum height=1cm, align=center, font=\small},
fbpos/.style={-{Stealth[length=3mm]}, very thick, popgreen},
neg/.style={-{Stealth[length=3mm]}, very thick, poppink, dashed},
stim/.style={-{Stealth[length=3mm]}, thick, popblue}
]
\node[box, align=center] (hypo) at (0,4.5){\textbf{Hypothalamus}\\ GnRH (pulses)};
\node[box, align=center] (hypo2) at (0,2.2){\textbf{Antéhypophyse}\\ FSH et LH};
\node[box, align=center] (ov) at (0,-0.6){\textbf{Ovaire}\\ Follicule puis corps jaune};
\node[box, fill=popgreenL, align=center] (e2) at (5.2,0.4){Œstradiol\\ (follicule dominant)};
\node[box, fill=poporangeL, align=center] (p4) at (5.2,-2.2){Progestérone + Inhibine\\ (corps jaune, follicules)};
\draw[stim] (hypo) -- node[right, font=\small]{stimule} (hypo2);
\draw[stim] (hypo2) -- node[right, font=\small]{croissance folliculaire, ovulation} (ov);
\draw[stim] (ov.east) -- (e2.west);
\draw[stim] (ov.east) -- (p4.west);
\draw[fbpos] (e2.north) .. controls (4.5,4.5) .. (hypo.east) node[midway, above, font=\small, popgreen]{\textbf{Rétrocontrôle $+$} (seuil critique) $\to$ pic de LH};
\draw[neg] (p4.west) .. controls (-4,-2.2) and (-4,2.2) .. (hypo2.west) node[midway, left, font=\small, poppink]{\textbf{Rétrocontrôle $-$} (phase lutéale)};
\draw[neg] (e2.west) .. controls (-3.5,0.4) and (-3.5,1.2) .. (hypo2.south west) node[midway, below left, font=\small, poppink]{\textbf{Rétrocontrôle $-$} (faible dose)};
\end{tikzpicture}
\end{center}
\legende{Schéma fonctionnel de la régulation des sécrétions ovariennes. Flèches bleues : stimulations ; flèche verte : rétrocontrôle positif de l'œstradiol à seuil critique (déclenche le pic de LH via l'hypothalamus) ; flèches rouges discontinues : rétrocontrôles négatifs (œstradiol à faible dose sur l'hypophyse, progestérone et inhibine sur l'hypophyse).}
\end{popfigure}

\begin{exoresolu}[Interprétation d'un dosage hormonal]
Chez une patiente de 29 ans consultant pour cycles irréguliers à l'hôpital de district de Bafoussam, on dose les hormones au jour 3 du cycle puis au jour 21. Résultats : jour 3, FSH = 7 UI/L, œstradiol = 150 pmol/L, progestérone = 2 nmol/L ; jour 21, progestérone = 4 nmol/L, LH sans pic décelé au jour 14. Interpréter ces résultats.
\tcblower
\textbf{Analyse} : au jour 3, les valeurs sont normales pour un début de cycle (hormones basses, FSH légèrement au-dessus de la LH pour le recrutement). Au jour 21, en pleine phase lutéale théorique, la progestérone devrait être comprise entre 30 et 100 nmol/L si une ovulation avait eu lieu ; or elle n'est que de 4 nmol/L, valeur typique de phase folliculaire. De plus, aucun pic de LH n'a été décelé au moment attendu (jour 14).\par
\textbf{Conclusion} : l'absence de pic de LH et l'absence d'élévation de la progestérone en milieu de cycle indiquent un \textbf{cycle anovulatoire} : aucun follicule n'a ovulé, donc aucun corps jaune ne s'est formé. Ce type de trouble, fréquent, est une cause d'infertilité que l'on peut traiter par inducteurs de l'ovulation (section 5).
\end{exoresolu}

\begin{savaistu}[Complément Bac : la kisspeptine, chef d'orchestre du pic de LH]
Des neurones hypothalamiques particuliers produisent la \cle{kisspeptine}, un peptide qui stimule puissamment les neurones à GnRH. L'œstradiol à seuil critique active ces neurones à kisspeptine : c'est ainsi que le rétrocontrôle positif déclenche la décharge de GnRH puis le pic de LH. Chez l'être humain, des mutations du gène de la kisspeptine ou de son récepteur provoquent une absence de puberté (hypogonadisme hypogonadotrope), preuve de son rôle central (connaissance admise, non exigible en détail).
\end{savaistu}

\begin{aretenir}[L'essentiel de la section]
Le cycle féminin est piloté par un double dialogue : l'hypophyse (FSH, LH) commande l'ovaire ; l'ovaire (œstradiol, progestérone, inhibine) rétrocontrôle l'axe hypothalamo-hypophysaire. Le rétrocontrôle est \textbf{négatif} la plupart du temps, mais devient \textbf{positif} lorsque l'œstradiol dépasse un seuil critique : c'est le déclencheur du pic de LH, donc de l'ovulation. La progestérone, produit du corps jaune, n'apparaît qu'après l'ovulation et sa chute provoque les règles. Cette logique de régulation est la clé de la contraception hormonale et de la prise en charge des stérilités, objets des sections suivantes.
\end{aretenir}

\coursec{Maîtrise de la reproduction : contraception et contragestion}

Dans la section précédente, nous avons établi que le cycle ovarien de la femme est piloté par un dialogue hormonal précis : la FSH stimule la croissance folliculaire, le pic de LH déclenche l'ovulation, et les hormones ovariennes (\oe{}stradiol, progestérone) exercent des rétrocontrôles sur le complexe hypothalamo-hypophysaire. Cette connaissance des mécanismes naturels de régulation ouvre une possibilité considérable : si l'on comprend comment le corps déclenche l'ovulation, on peut apprendre à la bloquer volontairement. C'est exactement le principe de la \cle{contraception hormonale}. Dans un pays comme le Cameroun, où la planification familiale est un enjeu majeur de santé publique (espacement des naissances, lutte contre les grossesses précoces et la mortalité maternelle), comprendre ces méthodes est à la fois un savoir scientifique exigible au Baccalauréat et un savoir-être citoyen.

\courssub{Définitions et principes généraux}

\begin{definition}[Contraception et contragestion]
La \cle{contraception} est l'ensemble des méthodes employées \emph{avant} ou indépendamment du rapport sexuel pour empêcher de façon \textbf{temporaire et réversible} la fécondation ou la nidation.

La \cle{contragestion} est une intervention réalisée \emph{après} un rapport sexuel non protégé (ou après échec d'une méthode contraceptive) pour empêcher la grossesse \textbf{avant l'implantation} de l'œuf dans l'endomètre. On parle couramment de « contraception d'urgence ».
\end{definition}

\begin{attention}[Contraception, contragestion, avortement : ne pas confondre]
La contraception agit \emph{avant} la fécondation, la contragestion agit \emph{avant la nidation} (implantation, vers le 6\ieme{}--7\ieme{} jour post-fécondation). L'interruption volontaire de grossesse (IVG) intervient \emph{après} la nidation : ce n'est ni une contraception ni une contragestion. Au Baccalauréat, cette distinction chronologique est souvent testée : situe toujours le moment d'action de la méthode sur la chronologie rapport $\rightarrow$ fécondation $\rightarrow$ nidation.
\end{attention}

L'idée directrice de la contraception hormonale est simple et élégante : pendant la grossesse, les taux élevés et \emph{constants} d'\oe{}strogènes et de progestérone bloquent l'ovulation par rétrocontrôle négatif. La pilule \emph{mimique} cet état hormonal : elle fait « croire » à l'hypothalamus et à l'hypophyse que les ovaires travaillent déjà, ce qui supprime la sécrétion de FSH et de LH et donc l'ovulation.

\courssub{La pilule œstro-progestative (pilule combinée)}

\textbf{Observation de départ.} Chez une femme prenant quotidiennement un comprimé associant un \oe{}strogène de synthèse (l'\cle{éthinylœstradiol}) et un progestatif de synthèse, les dosages sanguins montrent des taux d'hormones ovariennes bas et stables, et l'échographie ne révèle aucune ovulation. Comment des hormones qui, dans le cycle naturel, \emph{déclenchent} l'ovulation peuvent-elles la \emph{bloquer} ?

\begin{experience}[Mise en évidence expérimentale du blocage de l'ovulation]
\textbf{Protocole} : on prélève quotidiennement du sang chez deux groupes de femmes pendant 28 jours. Le groupe témoin ne prend aucun traitement ; le groupe test prend une pilule combinée (éthinylœstradiol + progestatif) du 1\ier{} au 21\ier{} jour, puis un placebo pendant 7 jours.

\textbf{Résultats} :
\begin{itemize}
\item Groupe témoin : pic d'\oe{}stradiol vers le 12\ieme{} jour, pic de LH et de FSH vers le 14\ieme{} jour, puis élévation de la progestérone en phase lutéale ; ovulation constatée.
\item Groupe test : taux de FSH et de LH bas et constants, \textbf{aucun pic de LH}, taux d'\oe{}stradiol et de progestérone endogènes bas ; aucune ovulation.
\end{itemize}

\textbf{Interprétation} : l'apport exogène et continu d'\oe{}strogène et de progestatif exerce un \emph{rétrocontrôle négatif permanent} sur l'hypothalamus (diminution des pulsations de GnRH) et sur l'hypophyse (chute de FSH et de LH). Sans FSH, pas de maturation folliculaire complète ; sans pic de LH, pas d'ovulation possible.
\end{experience}

\begin{propriete}[Mécanisme d'action de la pilule œstro-progestative]
La pilule combinée supprime le pic de LH par \textbf{rétroaction négative continue} de l'éthinylœstradiol (associé au progestatif) sur le complexe hypothalamo-hypophysaire. Elle agit à trois niveaux complémentaires :
\begin{enumerate}
\item \textbf{Niveau central} : blocage de la sécrétion de FSH et de LH $\Rightarrow$ inhibition de la folliculogenèse et \textbf{blocage de l'ovulation} (effet principal) ;
\item \textbf{Niveau du col de l'utérus} : le progestatif \textbf{épaissit la glaire cervicale}, qui devient imperméable aux spermatozoïdes ;
\item \textbf{Niveau de l'endomètre} : la muqueuse utérine reste \textbf{atrophique}, impropre à la nidation.
\end{enumerate}
\end{propriete}

\begin{attention}[Piège classique de rédaction]
Ne dis jamais que « la pilule contient de la progestérone » : elle contient un \emph{progestatif de synthèse} (molécule voisine, plus stable et active par voie orale). De même, le saignement qui survient pendant les 7 jours d'arrêt (ou de placebo) n'est \textbf{pas} de vraies règles : c'est une \emph{hémorragie de privation} due à la chute brutale des hormones exogènes, sans rapport avec un cycle ovulatoire.
\end{attention}

\begin{popfigure}
\begin{center}
\begin{tikzpicture}
\begin{axis}[
width=0.95\linewidth, height=7.5cm,
xlabel={Temps (jours)}, ylabel={Taux plasmatique (u. a.)},
xmin=0, xmax=28, ymin=0, ymax=10,
xtick={0,7,14,21,28},
legend style={at={(0.5,1.02)}, anchor=south, legend columns=4, font=\small},
grid=major, grid style={popline!40},
]
\addplot[popcurve, poppink, domain=0:21, samples=100] {3.2};
\addplot[popcurve, poppink, domain=21:28, samples=30] {0.6};
\addlegendentry{Hormones exogènes (pilule)}
\addplot[popcurve, popblue, domain=0:28, samples=100] {1.1 + 0.15*sin(deg(x))};
\addlegendentry{LH}
\addplot[popcurve, poppurple, domain=0:28, samples=100] {1.6 + 0.1*sin(deg(2*x))};
\addlegendentry{FSH}
\addplot[popcurve, poporange, domain=0:28, samples=100] {0.7 + 0.08*sin(deg(1.5*x))};
\addlegendentry{\oe{}stradiol et progestérone endogènes}
\addplot[popline, dashed, domain=0:28, samples=2] {0};
\draw[popline, dashed] (axis cs:21,0) -- (axis cs:21,10);
\node[font=\small, popmuted, anchor=west] at (axis cs:21.3,9.3) {arrêt 7 jours};
\end{axis}
\end{tikzpicture}
\end{center}
\legende{Évolution schématique des taux hormonaux chez une femme sous pilule œstro-progestative sur 28 jours. Les hormones exogènes (rose) sont apportées à taux constant du 1\ier{} au 21\ier{} jour ; la FSH (violet) et la LH (bleu) restent basses, \textbf{sans pic pré-ovulatoire} ; les hormones ovariennes endogènes (orange) restent basses et plates : il n'y a ni folliculogenèse aboutie, ni ovulation, ni corps jaune. Pendant les 7 jours d'arrêt, la chute des hormones exogènes provoque une hémorragie de privation.}
\end{popfigure}

\begin{methode}[Interpréter une courbe hormonale chez une utilisatrice de pilule]
Face à un document de type « dosages hormonaux sur 28 jours », procède en étapes :
\begin{enumerate}
\item Repère l'axe des abscisses (jours du cycle) et identifie chaque courbe grâce à la légende ;
\item Cherche le \textbf{pic de LH} vers le 14\ieme{} jour : son \emph{absence} est la signature d'une anovulation ;
\item Vérifie la \textbf{platitude} des courbes d'\oe{}stradiol et de progestérone endogènes : des taux bas et constants indiquent l'absence de follicule mûr et de corps jaune ;
\item Observe les taux de FSH : une FSH supprimée confirme le rétrocontrôle négatif exercé par les hormones exogènes ;
\item Conclus en reliant chaque observation au mécanisme : rétrocontrôle négatif continu $\Rightarrow$ pas de pic de LH $\Rightarrow$ pas d'ovulation.
\end{enumerate}
\end{methode}

\courssub{La pilule microprogestative (progestatif seul)}

Certaines femmes ne peuvent pas prendre d'\oe{}strogènes (antécédents vasculaires, allaitement, tabagisme après 35 ans). On utilise alors une pilule contenant \emph{uniquement un progestatif}, prise \textbf{en continu, sans interruption}.

\begin{propriete}[Mécanisme d'action de la pilule microprogestative]
Le progestatif seul agit principalement de façon \emph{périphérique} :
\begin{itemize}
\item \textbf{épaississement de la glaire cervicale} : obstacle majeur à la progression des spermatozoïdes (effet principal) ;
\item \textbf{modification de l'endomètre} : muqueuse peu réceptive à la nidation ;
\item effet central inconstant : l'ovulation est bloquée chez une proportion variable des utilisatrices selon la molécule (forte inhibition avec le désogestrel, partielle avec le lévonorgestrel) ; elle est donc \emph{souvent conservée}.
\end{itemize}
\end{propriete}

\begin{attention}[Limites de la pilule microprogestative]
La pilule microprogestative \textbf{ne garantit pas l'anovulation} : la fécondation reste possible, et la contraception repose surtout sur la glaire cervicale. Elle exige une \emph{observance stricte} (prise à heure fixe, tolérance de retard de 3 h seulement pour le lévonorgestrel, 12 h pour le désogestrel). Elle ne doit pas être utilisée comme méthode unique chez les femmes présentant un risque de \textbf{grossesse ectopique} (extra-utérine), car le progestatif ralentit la mobilité tubaire, ce qui favorise l'implantation de l'œuf dans la trompe.
\end{attention}

\courssub{La contragestion (contraception d'urgence)}

Un rapport non protégé, un oubli de pilule, un préservatif rompu : dans ces situations, il reste possible d'agir \emph{après} le rapport mais \emph{avant} la nidation.

\begin{definition}[Contragestion]
La contragestion (ou contraception d'urgence) est une intervention réalisée après un rapport sexuel non protégé pour empêcher la grossesse \textbf{avant l'implantation} de l'œuf. Deux principes actifs sont utilisés :
\begin{itemize}
\item le \cle{lévonorgestrel} (progestatif à forte dose, « pilule du lendemain »), efficace surtout dans les 72 h ;
\item l'\cle{ulipristal} (acétate d'ulipristal), efficace jusqu'à 120 h.
\end{itemize}
\end{definition}

\textbf{Mécanisme d'action.} L'effet dépend du moment du cycle où le produit est pris :
\begin{itemize}
\item \emph{avant l'ovulation} : la forte dose de progestatif (lévonorgestrel) ou le blocage des récepteurs de la progestérone (ulipristal) \textbf{retarde ou bloque le pic de LH}, donc \textbf{retarde l'ovulation} ; les spermatozoïdes (durée de vie de 3 à 5 jours) ne rencontrent alors aucun ovocyte ;
\item \emph{après l'ovulation} : le produit \textbf{modifie l'endomètre} et le rend impropre à la nidation, empêchant l'implantation de l'œuf.
\end{itemize}

\begin{savaistu}[L'ulipristal, un modulateur sélectif des récepteurs de la progestérone]
L'acétate d'ulipristal est un \emph{modulateur sélectif des récepteurs de la progestérone} (SPRM) : il se fixe sur les récepteurs sans les activer, bloquant l'action de la progestérone naturelle. Il reste efficace jusqu'à \textbf{120 heures} (5 jours) après le rapport non protégé, alors que le lévonorgestrel perd rapidement son efficacité au-delà de 72 h. Au Cameroun, la contraception d'urgence est disponible en pharmacie et dans les centres de santé de planification familiale ; son accès rapide est un enjeu de prévention des grossesses non désirées, notamment chez les adolescentes. Attention toutefois : c'est une méthode de \emph{rattrapage}, qui ne protège ni des infections sexuellement transmissibles (VIH/Sida, hépatite B) ni des rapports ultérieurs du même cycle.
\end{savaistu}

\courssub{Autres méthodes hormonales et contraception d'urgence au cuivre}

Le même principe (apport continu d'hormones $\Rightarrow$ rétrocontrôle négatif $\Rightarrow$ anovulation et glaire épaissie) se décline sous plusieurs formes, qui diffèrent seulement par la \emph{voie d'administration} et la durée d'action :

\begin{center}
\begin{tabularx}{\linewidth}{|l|X|X|}
\hline
\rowcolor{popblueL} \textbf{Méthode} & \textbf{Principe} & \textbf{Durée d'action} \\
\hline
Patch contraceptif & Œstro-progestatif diffusé à travers la peau & 1 semaine par patch \\
\hline
Anneau vaginal & Œstro-progestatif diffusé localement & 3 semaines \\
\hline
Implant sous-cutané & Progestatif seul diffusé en continu & 3 à 5 ans \\
\hline
DIU hormonal & Progestatif libéré dans l'utérus : glaire épaisse, endomètre atrophique & 3 à 5 ans \\
\hline
\end{tabularx}
\end{center}

\begin{propriete}[Contraception d'urgence au cuivre (complément Bac)]
Le \cle{DIU au cuivre}, posé jusqu'à 5 jours après un rapport non protégé, est la méthode de contragestion la plus efficace. Il est \textbf{non hormonal} : les ions cuivre sont \emph{spermicides} (ils altèrent la mobilité et la viabilité des spermatozoïdes) et provoquent une réaction inflammatoire locale de l'endomètre à effet \textbf{anti-implantation}. Il peut ensuite être conservé comme contraception durable.
\end{propriete}

\begin{popfigure}
\begin{center}
\begin{tikzpicture}[font=\small, >=Stealth, node distance=6mm]
% Axe corps
\node[draw, rounded corners, fill=poppurpleL, text width=3.1cm, align=center] (hypo) {\textbf{Hypothalamus + hypophyse}\\ (centre de commande)};
\node[draw, rounded corners, fill=poppinkL, text width=2.6cm, align=center, below left=9mm and 2mm of hypo] (ovaire) {\textbf{Ovaires}\\ follicule, ovulation};
\node[draw, rounded corners, fill=poporangeL, text width=2.6cm, align=center, below right=9mm and 2mm of hypo] (uterus) {\textbf{Utérus}\\ endomètre};
\node[draw, rounded corners, fill=popgreenL, text width=2.6cm, align=center, below=9mm of hypo, yshift=-16mm] (col) {\textbf{Col de l'utérus}\\ glaire cervicale};
% Flèches physiologiques
\draw[->, popdark, thick] (hypo) -- node[left, align=center, text width=2.2cm]{FSH, LH\\ (naturellement)} (ovaire);
\draw[->, popdark, thick] (ovaire.east) -- node[above, text width=2.6cm, align=center]{\oe{}stradiol, progestérone} (uterus.west);
% Actions des méthodes
\draw[->, poppink, very thick] (hypo.north) ++(0,0.15) -- ++(0,0.9) node[above, align=center, text width=4.6cm]{\textbf{Pilule combinée} : rétrocontrôle négatif continu $\Rightarrow$ pas de pic de LH};
\draw[->, poppink, very thick] (ovaire.south) -- ++(0,-0.7) node[below, align=center, text width=3.4cm]{blocage de l'ovulation (pilule combinée, implant)};
\draw[->, poppink, very thick] (col.south) -- ++(0,-0.7) node[below, align=center, text width=3.8cm]{glaire épaissie, hostile aux spermatozoïdes (progestatifs, DIU hormonal)};
\draw[->, poppink, very thick] (uterus.south) -- ++(0,-0.7) node[below, align=center, text width=4.2cm]{endomètre atrophique, impropre à la nidation (progestatifs, contragestion, DIU Cu)};
\end{tikzpicture}
\end{center}
\legende{Schéma comparatif des sites d'action des méthodes hormonales et du DIU au cuivre. \textbf{1.} Niveau central (hypothalamus-hypophyse) : la pilule combinée supprime FSH et LH. \textbf{2.} Ovaire : anovulation. \textbf{3.} Col de l'utérus : glaire cervicale épaissie (effet principal des progestatifs seuls). \textbf{4.} Endomètre : atrophie et non-réceptivité (progestatifs, contragestion, DIU hormonal et DIU au cuivre).}
\end{popfigure}

\courssub{Efficacité contraceptive : l'indice de Pearl et l'observance}

Pour comparer objectivement les méthodes, on utilise l'\cle{indice de Pearl} : nombre de grossesses survenues pour 100 femmes utilisant la méthode pendant un an.

\begin{exemplebox}[Indices de Pearl de quelques méthodes]
\begin{itemize}
\item Pilule combinée en \emph{usage parfait} (aucun oubli) : environ \textbf{0,3 \%} ;
\item Pilule combinée en \emph{usage typique} (oublis, retards, vomissements) : environ \textbf{9 \%} ;
\item Implant et DIU hormonal : moins de \textbf{1 \%} (efficacité indépendante de l'observance) ;
\item DIU au cuivre : environ \textbf{0,6 \%}.
\end{itemize}
L'écart spectaculaire entre usage parfait (0,3 \%) et usage typique (9 \%) montre que le principal facteur d'échec de la pilule est \emph{l'oubli}, d'où l'importance de l'\textbf{éducation thérapeutique} : expliquer le mécanisme, l'heure de prise, la conduite à tenir en cas d'oubli (prendre le comprimé dès que possible, utiliser un préservatif les 7 jours suivants si l'oubli dépasse 12 h).
\end{exemplebox}

\begin{exoresolu}[Analyse de courbes hormonales chez deux femmes]
\textbf{Énoncé.} On a dosé quotidiennement, pendant 28 jours, la LH, la FSH, l'\oe{}stradiol et la progestérone plasmatiques chez deux femmes A et B. Chez A, on observe un pic d'\oe{}stradiol au 12\ieme{} jour, un pic de LH au 14\ieme{} jour, puis une élévation de la progestérone de 1 à 18 unités arbitraires entre le 15\ieme{} et le 25\ieme{} jour. Chez B, les quatre hormones restent basses et constantes ; aucun pic n'est observé. L'une des deux femmes prend une pilule œstro-progestative.
\begin{enumerate}
\item Identifier la femme sous pilule, en justifiant.
\item Expliquer le mécanisme responsable des résultats observés chez elle.
\item Préciser deux autres sites d'action de cette pilule.
\end{enumerate}
\tcblower
\textbf{Solution détaillée.}
\begin{enumerate}
\item \textbf{La femme B} prend la pilule. Justification : chez elle, l'absence de pic de LH vers le 14\ieme{} jour traduit l'absence d'ovulation, et la platitude des courbes d'\oe{}stradiol et de progestérone montre qu'il n'y a ni follicule mûr (pas de pic d'\oe{}stradiol) ni corps jaune (pas d'élévation de progestérone). Chez A, la succession pic d'\oe{}stradiol $\rightarrow$ pic de LH $\rightarrow$ progestérone élevée correspond à un cycle ovulatoire normal.
\item La pilule apporte en continu de l'éthinylœstradiol et un progestatif. Ces hormones exogènes exercent un \emph{rétrocontrôle négatif permanent} sur l'hypothalamus (chute des pulsations de GnRH) et sur l'hypophyse (chute de FSH et de LH). Sans FSH, la folliculogenèse ne s'achève pas ; sans pic de LH, l'ovulation est impossible ; sans ovulation, pas de corps jaune, donc la progestérone endogène reste basse.
\item La pilule agit aussi sur le \textbf{col de l'utérus} (épaississement de la glaire cervicale, qui bloque le passage des spermatozoïdes) et sur l'\textbf{endomètre} (atrophie le rendant impropre à la nidation).
\end{enumerate}
\end{exoresolu}

\begin{exercice}[Pour t'entraîner]
Une femme sous pilule microprogestative (progestatif seul) présente, lors d'un dosage, un pic de LH discret au 15\ieme{} jour du cycle. 1. Ce résultat est-il compatible avec l'efficacité de sa contraception ? Justifie. 2. Quel est alors le principal mécanisme qui la protège d'une grossesse ? 3. Pourquoi cette méthode est-elle déconseillée chez une femme ayant déjà eu une grossesse extra-utérine ?
\end{exercice}

\begin{corrige}[Exercice « Pour t'entraîner »]
1. Oui, ce résultat est compatible : la pilule microprogestative ne garantit pas l'anovulation ; l'ovulation est conservée dans une proportion importante de cycles (surtout avec le lévonorgestrel). 2. La protection repose principalement sur l'épaississement de la glaire cervicale, imperméable aux spermatozoïdes, complété par la modification de l'endomètre (non-réceptivité à la nidation). 3. Le progestatif ralentit les mouvements ciliaires et le péristaltisme de la trompe ; si une fécondation survient malgré tout, la migration de l'œuf vers l'utérus est retardée, ce qui augmente le risque d'implantation tubaire (grossesse ectopique), accident grave déjà survenu chez cette femme.
\end{corrige}

\begin{aretenir}[L'essentiel de la section]
\begin{itemize}
\item La \textbf{pilule œstro-progestative} bloque l'ovulation par rétrocontrôle négatif continu (pas de pic de LH), épaissit la glaire cervicale et atrophie l'endomètre ; signature biologique : courbes d'hormones \emph{basses et plates}.
\item La \textbf{pilule microprogestative} agit surtout sur la glaire cervicale ; l'ovulation peut être conservée (variable selon la molécule).
\item La \textbf{contragestion} (lévonorgestrel jusqu'à 72 h, ulipristal jusqu'à 120 h) retarde l'ovulation ou empêche la nidation ; le \textbf{DIU au cuivre} est une contragestion non hormonale (spermicide et anti-implantation) la plus efficace.
\item L'efficacité se mesure par l'\textbf{indice de Pearl} ; l'écart entre usage parfait (0,3 \%) et typique (9 \%) souligne le rôle décisif de l'observance et de l'éducation thérapeutique.
\end{itemize}
\end{aretenir}

\coursec{Dérèglements des hormones sexuelles : causes, conséquences, traitement, prévention}

Dans la section précédente, nous avons vu comment la \cle{contraception} et la \cle{contragestion} exploitent volontairement les boucles de rétrocontrôle hormonal pour bloquer l'ovulation ou empêcher la nidation. Mais il arrive que ces boucles de régulation se dérèglent \emph{d'elles-mêmes}, sans intervention médicale : les gonades ne répondent plus correctement aux ordres hypophysaires, ou l'axe hypothalamo-hypophysaire cesse de commander. Il en résulte des troubles de la puberté, du cycle menstruel, de la fertilité, parfois du métabolisme général. Cette section a pour but d'analyser les principaux \cle{dérèglements hormonaux} de la reproduction, leurs causes, leurs conséquences, et les moyens de les traiter et de les prévenir — compétence essentielle pour comprendre la prise en charge de l'infertilité, fréquente au Cameroun comme ailleurs.

\courssub{Notion générale : l'hypogonadisme}

\begin{definition}[Hypogonadisme]
L'\cle{hypogonadisme} est l'insuffisance de production d'hormones sexuelles par les gonades (testicules chez l'homme, ovaires chez la femme). Il entraîne un déficit en \cle{testostérone} chez l'homme ou en \cle{œstrogènes} (principalement œstradiol) chez la femme, avec souvent une baisse ou un arrêt de la \cle{gamétogenèse}.
\end{definition}

L'intuition est simple : la gonade est une usine commandée par un chef d'orchestre (l'axe hypothalamo-hypophysaire). Si l'usine est en panne, la production chute \emph{malgré} des ordres abondants : c'est l'hypogonadisme \cle{primaire}. Si c'est le chef d'orchestre qui se tait, l'usine, pourtant saine, reste au ralenti faute d'ordres : c'est l'hypogonadisme \cle{secondaire} (ou central).

\begin{propriete}[Signature biologique des deux types d'hypogonadisme]
Grâce au rétrocontrôle négatif exercé par les hormones gonadiques sur l'hypothalamus et l'hypophyse :
\begin{itemize}
\item dans l'hypogonadisme \cle{primaire}, la chute de testostérone ou d'œstradiol lève le frein : la \cle{LH} et la \cle{FSH} sont \textbf{élevées} (hypogonadisme \emph{hypergonadotrope}) ;
\item dans l'hypogonadisme \cle{secondaire}, le défaut est en amont : la \cle{LH} et la \cle{FSH} sont \textbf{basses ou inappropriément normales} (hypogonadisme \emph{hypogonadotrope}).
\end{itemize}
\end{propriete}

\begin{methode}[Différencier un hypogonadisme primaire d'un hypogonadisme secondaire]
\begin{enumerate}
\item Doser dans le sang : \cle{testostérone} (homme) ou \cle{œstradiol} (femme), puis \cle{LH} et \cle{FSH} (au matin, à jeun).
\item Reporter les valeurs sur le schéma fonctionnel de la régulation (hypothalamus $\rightarrow$ GnRH $\rightarrow$ hypophyse $\rightarrow$ LH/FSH $\rightarrow$ gonade).
\item Si l'hormone gonadique est basse \textbf{et} LH/FSH élevées : la gonade est en cause $\Rightarrow$ hypogonadisme \textbf{primaire}.
\item Si l'hormone gonadique est basse \textbf{et} LH/FSH basses ou inappropriément normales : le défaut est hypothalamique ou hypophysaire $\Rightarrow$ hypogonadisme \textbf{secondaire}.
\item Confirmer éventuellement par un test à la GnRH (bolus IV : pas de réponse en LH/FSH si hypophysaire ; réponse exagérée si hypothalamique) ou une imagerie (IRM hypophysaire).
\end{enumerate}
\end{methode}

\begin{attention}[Piège classique au Bac]
Une LH élevée ne signifie jamais « trop de stimulation réussie » : elle signifie que l'hypophyse « crie » des ordres auxquels la gonade ne répond plus. Inversement, ne concluez jamais à un hypogonadisme secondaire sur le seul taux de testostérone bas : il faut \emph{toujours} croiser avec LH et FSH.
\end{attention}

\courssub{Hypogonadisme masculin : primaire et secondaire}

\textbf{Hypogonadisme primaire (testiculaire).} La lésion siège dans le testicule. Causes principales : le \cle{syndrome de Klinefelter} (caryotype 47,XXY : tubules séminifères hyalinisés, azoospermie, testostérone basse, LH et FSH très élevées) ; l'\cle{orchite} (inflammation testiculaire, notamment post-ourlienne — d'où l'importance du vaccin ROR contre les oreillons) ; les traumatismes, torsions testiculaires, cryptorchidies non corrigées précocement, radiothérapie ou chimiothérapie.

\textbf{Hypogonadisme secondaire (central).} La lésion siège à l'hypophyse ou à l'hypothalamus : \cle{adénome hypophysaire} (tumeur qui détruit les cellules gonadotropes par compression), \cle{syndrome de Kallmann} (déficit congénital en sécrétion de GnRH, associé à une anosmie — absence d'odorat), séquelles de chirurgie ou d'irradiation sellaire, anorexie mentale ou entraînement sportif intensif (hypogonadisme fonctionnel par inhibition de la pulsatilité GnRH).

\begin{exemplebox}[Seuils diagnostiques chiffrés (repères pour l'interprétation)]
Chez l'homme adulte, on évoque un hypogonadisme lorsque la \cle{testostérone totale} est inférieure à \SI{12}{nmol/L} (dosage à jeun, le matin entre 8 h et 10 h, confirmé par un second prélèvement), avec des signes cliniques (fatigue, baisse de la libido, infertilité, fonte musculaire). Chez la femme atteinte de SOPK, la testostérone totale dépasse souvent \SI{2,5}{nmol/L}. Chez un patient Klinefelter typique, on observe : testostérone $\approx$ \SI{8}{nmol/L}, LH $>$ \SI{20}{UI/L}, FSH $>$ \SI{30}{UI/L}.
\end{exemplebox}

\begin{popfigure}
\begin{center}
\begin{tabularx}{\linewidth}{|l|X|X|}
\hline
\rowcolor{popblueL} \textbf{Critère} & \textbf{Hypogonadisme primaire} & \textbf{Hypogonadisme secondaire} \\
\hline
Siège de la lésion & Testicule (ou ovaire) & Hypophyse ou hypothalamus \\
\hline
Exemples de causes & Klinefelter (47,XXY), orchite ourlienne, cryptorchidie, irradiation & Adénome hypophysaire, syndrome de Kallmann, anorexie, sport intensif \\
\hline
Testostérone / Œstradiol & Bas & Bas \\
\hline
LH et FSH & \textbf{Élevées} (perte du rétrocontrôle négatif) & \textbf{Basses ou inappropriément normales} (défaut de commande) \\
\hline
Traitement étiologique & Aucun (lésion irréversible le plus souvent) & Chirurgie de l'adénome, agonistes dopaminergiques, prise en charge nutritionnelle \\
\hline
Traitement substitutif & Testostérone (homme), œstroprogestatifs (femme) & Gonadotrophines (LH/FSH) ou GnRH pulsatile si désir de grossesse \\
\hline
Fertilité & Souvent compromise (azoospermie, épuisement folliculaire) & Souvent récupérable sous gonadotrophines exogènes \\
\hline
\end{tabularx}
\end{center}
\legende{Tableau comparatif des hypogonadismes primaire et secondaire : siège de la lésion, signature hormonale et stratégie thérapeutique. La clé du diagnostic différentiel est le couple « hormone gonadique basse + taux de LH/FSH », lu à la lumière du rétrocontrôle négatif.}
\end{popfigure}

\courssub{Hypogonadisme et troubles hormonaux chez la femme}

\textbf{Insuffisance ovarienne précoce (IOP).} Il s'agit d'un arrêt de fonctionnement des ovaires avant 40 ans : aménorrhée (absence de règles depuis 4 mois au moins), œstradiol bas, FSH élevée ($>$ \SI{25}{UI/L} à deux reprises espacées d'un mois) — c'est un hypogonadisme \emph{primaire} féminin. Causes : épuisement prématuré du stock folliculaire (origine génétique : délétion X, mutation FMR1 ; auto-immune : thyroïdite, surrénalite ; iatrogène : chimiothérapie alkylante, radiothérapie pelvienne, chirurgie ovarienne extensive). Conséquences : infertilité définitive le plus souvent, ostéoporose précoce, symptômes climatériques (bouffées de chaleur, sécheresse), risque cardiovasculaire accru. Traitement : substitution hormonale de la ménopause (œstrogènes + progestatif) jusqu'à l'âge habituel de la ménopause (\SI{50}{ans} env.) pour protéger l'os et le métabolisme.

\textbf{Syndrome des ovaires polykystiques (SOPK).} C'est le trouble endocrinien le plus fréquent de la femme en âge de procréer (5 à 10 \% des femmes). Il associe classiquement :
\begin{itemize}
\item une \cle{hyperandrogénie} (excès de testostérone : hirsutisme score de Ferriman-Gallwey $>$ 6, acné, alopécie androgénique) ;
\item une \cle{anovulation} chronique (cycles irréguliers $>$ \SI{35}{jours} ou absents, infertilité) ;
\item des ovaires augmentés de volume, à nombreux petits follicules bloqués en début de maturation (aspect « polykystique » à l'échographie : $\ge$ 12 follicules de 2 à \SI{9}{mm} par ovaire et/ou volume ovarien $>$ \SI{10}{mL}).
\end{itemize}
Le mécanisme central est une boucle vicieuse : la \cle{résistance à l'insuline} (fréquente en cas de surpoids/obésité androïde) élève l'insulinémie ; l'insuline (via son récepteur et le récepteur à l'IGF-1) stimule la production d'androgènes par les \cle{cellules de la thèque} ovarienne ; l'hyperandrogénie locale bloque la maturation folliculaire (arrêt au stade antral) ; l'absence d'ovulation supprime le pic de LH et la sécrétion de progestérone, entretenant l'hyperplasie thécale et l'hyper LH. La LH est souvent élevée (rapport LH/FSH $> 2$), la FSH normale ou basse, l'œstradiol normal (production continue par les nombreux follicules), et l'\cle{AMH} (hormone anti-müllérienne, reflet du stock de follicules antraux) est très élevée ($>$ \SI{35}{pmol/L} ou \SI{5}{ng/mL}).

\begin{attention}[Critères de Rotterdam — piège du diagnostic]
Le SOPK ne se diagnostique \textbf{pas} sur la seule échographie : des ovaires « polykystiques » existent chez 20 \% des femmes sans aucune maladie. Il faut \textbf{2 critères sur 3} : (1) hyperandrogénie clinique ou biologique ; (2) anovulation ou cycles irréguliers ; (3) ovaires polykystiques à l'échographie. Il faut aussi éliminer les étiologies mimétiques : hyperplasie surrénalienne congénitale (dosage 17-OH-progestérone), syndrome de Cushing, tumeur sécrétante androgènes, hyperprolactinémie, troubles thyroïdiens.
\end{attention}

\begin{popfigure}
\begin{center}
\begin{tikzpicture}
\begin{axis}[
  width=0.95\linewidth, height=7cm,
  ybar, bar width=12pt,
  symbolic x coords={LH, FSH, {\OE}stradiol, Testostérone, AMH},
  xtick=data,
  ylabel={Taux rapporté à la limite haute de la normale (valeur 1)},
  ymin=0, ymax=2.8,
  ymajorgrids, grid style={popline!40},
  legend style={at={(0.5,-0.25)}, anchor=north, legend columns=2, draw=none, font=\small},
  nodes near coords, every node near coord/.append style={font=\small, rotate=90, anchor=west, yshift=2pt},
  enlarge x limits=0.2,
]
\addplot[fill=popblue, draw=popdark] coordinates {(LH,1.9) (FSH,0.85) ({\OE}stradiol,0.95) (Testostérone,1.7) (AMH,2.4)};
\addplot[fill=popmuted!40, draw=popdark] coordinates {(LH,1.0) (FSH,1.0) ({\OE}stradiol,1.0) (Testostérone,1.0) (AMH,1.0)};
\legend{Femme atteinte de SOPK, Femme témoin (norme)}
\end{axis}
\end{tikzpicture}
\end{center}
\legende{Profil hormonal typique du SOPK (valeurs normalisées : 1 = limite supérieure de la normale). On observe une LH élevée (rapport LH/FSH $> 2$), une FSH normale ou basse, un œstradiol normal, une testostérone et une AMH nettement élevées. Ce profil reflète l'hyperandrogénie thécale et l'accumulation de follicules antraux immatures.}
\end{popfigure}

\textbf{Hyperprolactinémie.} La \cle{prolactine} (hormone hypophysaire antérieure de la lactation) inhibe, à dose élevée, la sécrétion pulsatile de \cle{GnRH} par l'hypothalamus (via stimulation des neurones dopaminergiques tubéro-infundibulaires). Une hyperprolactinémie (adénome hypophysaire à prolactine ou \cle{prolactinome}, certains médicaments : neuroleptiques, antidépresseurs, antiémétiques métoclopramide ; hypothyroïdie primitive ; grossesse) provoque donc un hypogonadisme \emph{secondaire} fonctionnel : \cle{aménorrhée} (absence de règles), \cle{galactorrhée} (écoulement de lait hors allaitement) et infertilité chez la femme ; baisse de la libido, dysfonction érectile, infertilité et parfois gynécomastie chez l'homme. Le traitement de référence est un agoniste dopaminergique (cabergoline, 0,5 à \SI{1}{mg} par semaine), qui normalise la prolactine, restaure l'axe gonadotrope et permet souvent le retour de la fertilité. La chirurgie est réservée aux macro-adénomes résistants ou compressifs.

\courssub{Conséquences cliniques des dérèglements hormonaux}

Au-delà de la sphère reproductive, les hormones sexuelles ont des effets systémiques majeurs :
\begin{itemize}
\item \cle{Infertilité} : anovulation (SOPK, hyperprolactinémie, IOP) ou azoospermie/oligospermie sévère (hypogonadismes masculins primaires ou secondaires non traités) ;
\item \cle{Troubles métaboliques} : dans le SOPK, la résistance à l'insuline expose au diabète de type 2, à la dyslipidémie (hypertriglycéridémie, HDL bas) et à la stéatose hépatique non alcoolique (NASH) ; le déficit androgénique masculin favorise l'obésité abdominale, la résistance à l'insuline et le syndrome métabolique ;
\item \cle{Ostéoporose} : les œstrogènes et la testostérone (via aromatisation en œstradiol) protègent le capital osseux en inhibant la résorption ostéoclastique ; leur carence prolongée (IOP, hypogonadisme masculin non traité, aménorrhée de l'athlète) fragilise le squelette (ostéopénie puis ostéoporose) ;
\item \cle{Troubles psychologiques et sociaux} : baisse de l'estime de soi, anxiété, dépression, retentissement majeur sur la vie de couple — particulièrement lourd dans notre contexte social camerounais où la stérilité est souvent vécue comme une stigmatisation, une « honte » familiale, source de violences conjugales ou de répudiation.
\end{itemize}

\courssub{Traitements}

\begin{itemize}
\item \cle{Substitution hormonale} : testostérone (gel transdermique \SI{50}{mg}/jour, injection d'undécanoate \SI{1000}{mg} tous les 3 mois) chez l'homme hypogonadique ne désirant pas de descendance immédiate ; œstrogènes + progestatifs (pilule ou THM) chez la femme en IOP. \textbf{Attention :} la testostérone exogène bloque la spermatogenèse par rétrocontrôle négatif sur LH/FSH — elle est formellement contre-indiquée si l'on recherche la fertilité.
\item \cle{Inducteurs d'ovulation} : le \cle{citrate de clomifène} (anti-œstrogène compétitif au niveau hypothalamique, lève le rétrocontrôle négatif $\rightarrow$ montée FSH/LH) ou le \cle{letrozole} (inhibiteur de l'aromatase, abaisse les œstrogènes $\rightarrow$ même effet, souvent plus efficace dans le SOPK) en première intention ; les \cle{gonadotrophines} injectables (FSH recombinante, hMG, hCG pour déclencher) en seconde intention ou dans les hypogonadismes secondaires (où l'ovaire est sain).
\item \cle{Agonistes et antagonistes de la GnRH} : ils désensibilisent (agonistes : busereline, triptoréline) ou bloquent (antagonistes : cétrolix, ganirelix) l'axe hypothalamo-hypophysaire (utilisés dans l'endométriose, les fibromes utérins symptomatiques, et les protocoles de stimulation contrôlée pour FIV).
\item \cle{Metformine} : dans le SOPK avec résistance à l'insuline avérée (HOMA-IR $>$ 2,5), elle améliore la sensibilité à l'insuline, abaisse l'insulinémie, réduit l'hyperandrogénie ovarienne et peut restaurer l'ovulation spontanée (connaissance admise au Bac).
\item Traitement étiologique : chirurgie transsphénoïdale d'un adénome hypophysaire compressif, agonistes dopaminergiques dans l'hyperprolactinémie, thyroxine si hypothyroïdie cause d'hyperprolactinémie.
\end{itemize}

\begin{savaistu}[Pistes d'adjuvants dans le SOPK]
La \cle{vitamine D} (déficit fréquent dans le SOPK) et l'\cle{inositol} (notamment le myo-inositol, seconde messager de l'insuline) sont étudiés comme adjuvants : ils amélioreraient la sensibilité à l'insuline, la qualité ovocytaire et le taux de grossesse. Ces compléments ne remplacent pas les traitements de référence (hygiène de vie, letrozole/clomifène, metformine), mais illustrent la recherche actuelle sur le volet métabolique de ce syndrome.
\end{savaistu}

\courssub{Prévention}

\begin{itemize}
\item \cle{Dépistage précoce} : devant des cycles irréguliers, une aménorrhée, un retard pubertaire (pas de caractères sexuels secondaires à \SI{14}{ans} fille, \SI{14}{ans} garçon), une infertilité de couple (12 mois de rapports non protégés), réaliser des dosages hormonaux ciblés (LH, FSH, œstradiol, testostérone, prolactine, AMH, TSH) et une échographie pelvienne (femme) ou testiculaire (homme) ;
\item \cle{Hygiène de vie} : maintien d'un poids sain (IMC 18,5--25) et activité physique régulière (une perte de 5 à 10 \% du poids corporel restaure souvent l'ovulation spontanée dans le SOPK avec surpoids/obésité) ; éviter tabac (toxique ovarien direct, ménopause précoce) et alcool, toxiques gonadiques professionnels (pesticides, solvants, métaux lourds) ;
\item \cle{Vaccination contre les oreillons} (ROR) : elle prévient l'orchite ourlienne, cause majeure d'hypogonadisme primaire et de stérilité masculine définitive ;
\item \cle{Éducation sexuelle et à la santé reproductive} : consultation précoce en cas de troubles du cycle, lutte contre les infections sexuellement transmissibles (Chlamydia, Gonocoque : cause majeure de stérilité tubaire par salpingite séquellaire), information sans tabou sur l'infertilité du couple (partage de la responsabilité, recours aux structures spécialisées comme le Centre de Procréation Médicalement Assistée de l'Hôpital Central de Yaoundé ou de la Maternité de l'Hôpital Laquintinie de Douala).
\end{itemize}

\begin{exoresolu}[Interpréter un bilan hormonal]
Un homme de 28 ans consulte au Centre Hospitalier d'Essos pour infertilité du couple. Bilan hormonal (prélèvement à 8 h, à jeun) : testostérone totale = \SI{7}{nmol/L} (norme : \num{12} à \SI{35}{nmol/L}) ; LH = \SI{1,5}{UI/L} (norme : \num{2} à \SI{9}{UI/L}) ; FSH = \SI{1}{UI/L} (norme : \num{1} à \SI{8}{UI/L}) ; prolactine = \SI{900}{mUI/L} (norme $<$ \SI{400}{mUI/L}).
\begin{enumerate}
\item S'agit-il d'un hypogonadisme primaire ou secondaire ? Justifier.
\item Proposer une hypothèse étiologique et le traitement de première intention.
\end{enumerate}
\tcblower
1) La testostérone est effondrée (hypogonadisme), mais la LH et la FSH sont \textbf{basses} (inférieures aux normes) : l'hypophyse ne « crie » pas alors que la gonade est au ralenti. Le défaut est donc en amont du testicule : il s'agit d'un hypogonadisme \textbf{secondaire} (hypogonadotrope). 2) La prolactine est très élevée (\SI{900}{mUI/L}) : l'hyperprolactinémie inhibe la sécrétion pulsatile de GnRH hypothalamique, ce qui éteint la sécrétion de LH/FSH. L'hypothèse étiologique forte est un \textbf{adénome hypophysaire à prolactine (prolactinome)}, à confirmer par IRM sellaire (recherche d'une micro- ou macro-adénome). Le traitement de première intention repose sur un agoniste dopaminergique (cabergoline), qui normalise la prolactine, restaure l'axe gonadotrope (LH, FSH, testostérone) et permet souvent le retour de la fertilité en quelques mois.
\end{exoresolu}

\begin{aretenir}[L'essentiel de la section]
\begin{itemize}
\item Hypogonadisme = insuffisance de production d'hormones sexuelles ; \textbf{primaire} (gonade en cause : LH/FSH élevées = hypergonadotrope) ou \textbf{secondaire} (axe en cause : LH/FSH basses = hypogonadotrope).
\item Le diagnostic différentiel repose \emph{toujours} sur le croisement hormone gonadique (testostérone / œstradiol) / gonadotrophines (LH, FSH), lu avec le schéma de rétrocontrôle négatif.
\item SOPK : 2 critères de Rotterdam sur 3 (hyperandrogénie + anovulation + ovaires polykystiques) ; mécanisme central = résistance à l'insuline $\rightarrow$ hyperinsulinisme $\rightarrow$ hyperandrogénie ovarienne $\rightarrow$ blocage folliculaire ; traitement : hygiène de vie (perte de poids), letrozole/clomifène, metformine.
\item Hyperprolactinémie : aménorrhée-galactorrhée (femme) ou hypogonadisme (homme) par inhibition de la GnRH ; traitement : cabergoline.
\item Prévention : dépistage précoce, poids/activité physique, vaccin ROR, éducation à la santé reproductive et lutte contre les IST.
\end{itemize}
\end{aretenir}

\begin{exercice}[À toi de jouer]
Une femme de 24 ans présente une aménorrhée secondaire depuis 8 mois, un hirsutisme (score de Ferriman-Gallwey = 12) et un surpoids (IMC = \SI{31}{kg/m^2}). Bilan hormonal (J3 du cycle induit par duphaston) : testostérone totale = \SI{3,1}{nmol/L} ; LH = \SI{14}{UI/L} ; FSH = \SI{6}{UI/L} ; œstradiol = \SI{150}{pmol/L} (normal) ; prolactine et TSH normales. Échographie pelvienne : utérus normal, ovaires de volume augmenté (droit \SI{12}{mL}, gauche \SI{11}{mL}) avec 15 follicules périphériques de 3 à \SI{8}{mm} par ovaire.
\begin{enumerate}
\item Le diagnostic de SOPK est-il justifié selon les critères de Rotterdam ? Justifier chaque critère.
\item Expliquer le rôle de la résistance à l'insuline dans la genèse de l'hyperandrogénie et de l'anovulation chez cette patiente.
\item Proposer une prise en charge en trois volets (hygiène de vie, médical, fertilité).
\end{enumerate}
\end{exercice}

\begin{corrige}[Exercice]
\begin{enumerate}
\item \textbf{Oui}, le diagnostic de SOPK est pleinement justifié : la patiente réunit les \textbf{3 critères de Rotterdam} :
      \begin{itemize}
      \item \textbf{Hyperandrogénie} : clinique (hirsutisme score 12) \textbf{et} biologique (testostérone \SI{3,1}{nmol/L} $>$ \SI{2,5}{nmol/L}) ;
      \item \textbf{Anovulation} : aménorrhée secondaire de 8 mois (cycles $>$ \SI{35}{jours} ou absents) ;
      \item \textbf{Ovaires polykystiques} : volume ovarien $>$ \SI{10}{mL} bilatéral et $\ge$ 12 follicules de 2--9 mm par ovaire (15 observés).
      \end{itemize}
      Le rapport LH/FSH = 14/6 $\approx$ 2,3 $>$ 2 conforte le diagnostic. Les causes mimétiques sont exclues (prolactine, TSH normales).
\item La résistance à l'insuline liée au surpoids (IMC 31) entraine un \textbf{hyperinsulinisme compensateur}. L'insuline en excès se lie aux récepteurs à l'IGF-1 sur les \textbf{cellules de la thèque} ovarienne et potentialise l'action de la LH, stimulant fortement l'enzyme \cle{cytochrome P450c17} (17$\alpha$-hydroxylase/17,20-lyase) $\rightarrow$ surproduction d'androgènes (androstènedione, testostérone). L'hyperandrogénie locale \textbf{bloque la maturation folliculaire} (arrêt au stade antral, pas de follicule dominant) $\rightarrow$ \textbf{anovulation}. L'absence de corps jaune $\rightarrow$ pas de progestérone $\rightarrow$ pas de rétrocontrôle négatif sur la LH $\rightarrow$ \textbf{LH chroniquement élevée} (rapport LH/FSH $>$ 2) qui entretient l'hyperplasie thécale. C'est un cercle vicieux.
\item \textbf{Prise en charge en trois volets :}
      \begin{enumerate}
      \item \textbf{Hygiène de vie (pierre angulaire) :} régime hypocalorique méditerranéen + activité physique \SI{150}{min}/semaine $\rightarrow$ objectif perte de 5 à 10 \% du poids initial (\num{4} à \SI{8}{kg}). Cela restaure souvent la sensibilité à l'insuline, abaisse les androgènes et induit le retour d'ovulations spontanées.
      \item \textbf{Traitement médical de l'hyperandrogénie/résistance à l'insuline :} \textbf{Metformine} (\SI{850}{mg} à \SI{1700}{mg}/jour) si HOMA-IR élevé ou intolérance au glucose. Si hirsutisme majeur : association pilule œstroprogestative (anti-androgénique : cyprotérone, drospirénone) \emph{si pas de désir de grossesse immédiat}.
      \item \textbf{Si désir de grossesse (après échec de 6 à 12 mois d'ovulation spontanée post-perte de poids) :} Induction de l'ovulation par \textbf{letrozole} (\SI{2,5}{mg} à \SI{7,5}{mg}/jour J3--J7) en première intention (meilleur taux de naissance vivante, moins d'effet anti-œstrogène sur l'endomètre que le clomifène), ou citrate de clomifène. Surveillance échographique (échographie de contrôle J12--J13) pour éviter la grossesse multiple. En cas d'échec (résistance au letrozole/clomifène) : gonadotrophines injectables (FSHr) en milieu spécialisé (PMA).
      \end{enumerate}
\end{enumerate}
\end{corrige}

\coursec{Stérilité et procréation médicalement assistée (PMA)}

Dans les sections précédentes, nous avons vu que les dérèglements hormonaux (insuffisance de testostérone, troubles de l'ovulation, hyperstimulation) perturbent la fonction reproductive. Mais même lorsque les axes hormonaux fonctionnent correctement, un couple peut ne pas parvenir à concevoir un enfant. C'est le drame de la \cle{stérilité}, problème majeur de santé reproductive qui touche, selon l'OMS, environ un couple sur six dans le monde — et au Cameroun, où elle est souvent vécue comme une injustice sociale et une source de stigmatisation, en particulier pour les femmes. Cette dernière section répond à trois questions : qu'est-ce que la stérilité ? Quelles en sont les causes chez l'homme et chez la femme ? Et comment la science médicale permet-elle aujourd'hui de la contourner grâce à la \cle{procréation médicalement assistée} (PMA) ?

\courssub{Définitions : stérilité et infertilité}

\begin{definition}[Stérilité et infertilité]
On parle de \cle{stérilité} (ou d'infécondité du couple) lorsqu'un couple n'obtient \textbf{aucune grossesse après 12 mois de rapports sexuels réguliers (2 à 3 par semaine) et non protégés}. On distingue :
\begin{itemize}
\item la \textbf{stérilité primaire} : le couple n'a jamais eu d'enfant ;
\item la \textbf{stérilité secondaire} : le couple a déjà eu un enfant (ensemble ou séparément) mais n'arrive plus à concevoir.
\end{itemize}
L'\cle{infertilité} désigne une \textbf{diminution de la capacité à concevoir} (difficulté à obtenir une grossesse, délai allongé), sans que celle-ci soit impossible. La stérilité est donc une infertilité totale et définitive sans prise en charge.
\end{definition}

\begin{attention}[Un problème de couple, pas de la femme seule]
Dans l'imaginaire social camerounais, l'absence d'enfant est presque toujours imputée à la femme. C'est scientifiquement faux : les études épidémiologiques montrent que la cause est \textbf{masculine dans environ 30 à 40 \% des cas}, féminine dans 40 à 50 \% des cas, mixte ou inexpliquée dans le reste (10 à 20 \%). Tout bilan de stérilité doit donc être mené \textbf{simultanément chez les deux partenaires}.
\end{attention}

\courssub{Causes de la stérilité}

\textbf{A. Causes masculines}\par

La fertilité masculine exige une spermatogenèse normale (dépendant de la testostérone et de la FSH, vus en sections 1 et 2) et des voies génitales perméables. Les principales causes de stérilité masculine sont :

\begin{itemize}
\item l'\textbf{oligo-asthéno-tératozoospermie} (\cle{OATS}) : association d'une concentration en spermatozoïdes trop faible (\emph{oligo-}), d'une mobilité insuffisante (\emph{asthéno-}) et d'un pourcentage élevé de formes anormales (\emph{térato-}) ;
\item l'\textbf{azoospermie} : absence totale de spermatozoïdes dans l'éjaculat (voir définition ci-dessous) ;
\item le \textbf{varicocèle} : dilatation variqueuse des veines du cordon spermatique, qui élève la température testiculaire et altère la spermatogenèse (cause masculine la plus fréquente, environ 15 \% de la population masculine générale et jusqu'à 40 \% chez les hommes consultant pour infertilité) ;
\item les \textbf{infections} des voies génitales (orchites, épididymites, IST non traitées comme la blennorragie ou l'infection à \emph{Chlamydia trachomatis}), qui peuvent obstruer les canaux déférents ;
\item les \textbf{causes génétiques} : microdélétions du chromosome Y (perte de gènes de la spermatogenèse dans la région AZF), mutation du gène \emph{CFTR} (mucoviscidose) responsable d'une absence bilatérale congénitale des canaux déférents, aneuploïdies (syndrome de Klinefelter, caryotype 47,XXY).
\end{itemize}

\begin{definition}[Azoospermie]
L'\cle{azoospermie} est l'\textbf{absence totale de spermatozoïdes dans l'éjaculat}, confirmée sur deux spermogrammes espacés de 3 mois. On distingue :
\begin{itemize}
\item l'\textbf{azoospermie obstructive} : les testicules produisent des spermatozoïdes normalement, mais la \textbf{voie excrétrice est bouchée} (épididyme, canal déférent) — par exemple après infection, chirurgie ou absence congénitale des canaux déférents ; les spermatozoïdes peuvent être prélevés directement dans le testicule (TESE) ou l'épididyme (MESA) pour une ICSI ;
\item l'\textbf{azoospermie sécrétoire} (ou non obstructive) : \textbf{défaillance de la spermatogenèse} elle-même (testicules atrophiés, cause génétique, irradiation, chimiothérapie, varicocèle sévère non traité) ; le pronostic est beaucoup plus réservé.
\end{itemize}
\end{definition}

\textbf{B. Causes féminines}\par

Chez la femme, la conception exige une ovulation, des trompes perméables (lieu de la fécondation) et un utérus apte à la nidation. Les principales causes sont :

\begin{itemize}
\item les \textbf{troubles de l'ovulation} : syndrome des ovaires polykystiques (\cle{SOPK}, anovulation chronique avec hyperandrogénie), insuffisance ovarienne prématurée (\cle{IOP}, épuisement du stock folliculaire avant 40 ans), dérèglements hypothalamo-hypophysaires (stress, anorexie, hyperprolactinémie) ;
\item les \textbf{atteintes tubaires} : \textbf{salpingites} (infections des trompes, souvent consécutives à des IST mal traitées — cause majeure en Afrique subsaharienne) et \textbf{endométriose} (présence de muqueuse utérine hors de l'utérus, qui déforme et obstrue les trompes, altère la qualité ovocytaire et la réceptivité utérine) ;
\item les \textbf{anomalies utérines} : fibromes (tumeurs bénignes du myomètre, très fréquents chez la femme africaine, pouvant déformer la cavité utérine), malformations congénitales (utérus cloisonné, bicorne), synéchies (cicatrices après curetages répétés ou infections) ;
\item l'\textbf{âge} : la réserve ovarienne et la qualité des ovocytes chutent nettement après 35 ans ; à 42 ans, la probabilité de grossesse par cycle est très faible (< 5 \%) et le risque de fausse couche chromosomique augmente.
\end{itemize}

\begin{savaistu}[Stérilité et infections en Afrique subsaharienne]
En Afrique centrale et de l'Ouest, la proportion de stérilités dues aux \textbf{infections génitales} (blennorragie, chlamydioses) est beaucoup plus élevée qu'en Europe : on parle de « ceinture de l'infertilité » décrite par les démographes. D'où l'importance capitale du \textbf{dépistage et du traitement précoce des IST} et du conseil préconceptionnel dans les centres de santé camerounais : prévenir une salpingite, c'est préserver la fertilité future.
\end{savaistu}

\courssub{Le bilan de fertilité : le spermogramme}

Le premier examen du bilan masculin est le \cle{spermogramme} : analyse macroscopique et microscopique du sperme recueilli par masturbation après \textbf{2 à 7 jours d'abstinence sexuelle}. Il est interprété selon les valeurs de référence de l'OMS (6\up{e} édition, 2021).

\begin{methode}[Lecture d'un spermogramme]
Pour interpréter un spermogramme de façon rigoureuse, procède toujours dans le même ordre :
\begin{enumerate}
\item vérifier les conditions de recueil (abstinence de 2 à 7 jours, recueil complet dans un flacon stérile, analyse dans l'heure qui suit, température maintenue à 20--37 °C) ;
\item contrôler le \textbf{volume} (référence : $\geq$ \SI{1,4}{mL}) et le \textbf{pH} ($\geq$ 7,2) — un volume nul avec pH acide évoque une obstruction des voies excrétrices ou une absence d'éjaculation (anéjaculation) ;
\item évaluer la \textbf{concentration} en spermatozoïdes (référence : $\geq$ \SI{16}{millions/mL}) ;
\item évaluer la \textbf{mobilité progressive} (référence : $\geq$ \SI{32}{\%}) ;
\item évaluer la \textbf{morphologie} (référence : $\geq$ \SI{4}{\%} de formes normales selon les critères stricts de Kruger) et la \textbf{vitalité} ($\geq$ \SI{54}{\%} de spermatozoïdes vivants, test à l'éosine-nigrosine) ;
\item conclure : normal, oligo- / asthéno- / tératozoospermie (isolées ou associées en OATS), ou azoospermie — à confirmer sur un \textbf{deuxième prélèvement} à 3 mois d'intervalle (durée d'un cycle de spermatogenèse : environ \SI{74}{jours}).
\end{enumerate}
\end{methode}

\begin{popfigure}
\begin{center}
\small
\begin{tabular}{|l|c|c|}
\hline
\rowcolor{popblueL} \textbf{Paramètre (OMS 2021)} & \textbf{Valeur normale (5\up{e} centile)} & \textbf{Anomalie si...} \\
\hline
Volume de l'éjaculat & $\geq$ \SI{1,4}{mL} & $<$ \SI{1,4}{mL} (hypospermie) \\
\hline
pH & $\geq$ 7,2 & $<$ 7,2 (obstruction des vésicules séminales ?) \\
\hline
Concentration & $\geq$ \SI{16}{millions/mL} & $<$ \SI{16}{M/mL} (oligozoospermie) \\
\hline
Mobilité progressive & $\geq$ \SI{32}{\%} & $<$ \SI{32}{\%} (asthénozoospermie) \\
\hline
Mobilité totale (prog. + non prog.) & $\geq$ \SI{40}{\%} & $<$ \SI{40}{\%} \\
\hline
Formes normales (Kruger) & $\geq$ \SI{4}{\%} & $<$ \SI{4}{\%} (tératozoospermie) \\
\hline
Vitalité & $\geq$ \SI{54}{\%} & $<$ \SI{54}{\%} (nécrozoospermie) \\
\hline
Leucocytes & $<$ \SI{1}{million/mL} & $\geq$ \SI{1}{M/mL} (leucospermie = infection) \\
\hline
Spermatozoïdes & présents & absents : azoospermie \\
\hline
\end{tabular}
\end{center}
\legende{Valeurs de référence du spermogramme normal selon l'OMS (6\up{e} éd., 2021) et seuils pathologiques. Une anomalie isolée ne signifie pas la stérilité : c'est l'ensemble du tableau, confirmé sur deux examens à 3 mois d'intervalle, qui oriente le diagnostic.}
\end{popfigure}

\begin{exoresolu}[Interprétation d'un spermogramme]
Énoncé. Monsieur N., 34 ans, consulte pour absence de grossesse après 2 ans de mariage. Son spermogramme (après 4 jours d'abstinence) donne : volume = \SI{2,0}{mL} ; pH = 7,4 ; concentration = \SI{8}{millions/mL} ; mobilité progressive = \SI{20}{\%} ; formes normales = \SI{3}{\%} ; vitalité = \SI{60}{\%}. Interpréter ce résultat et proposer une conduite à tenir.
\tcblower
Solution. Appliquons la méthode de lecture :
\begin{enumerate}
\item Volume (\SI{2,0}{mL} $\geq$ \SI{1,4}{mL}) et pH (7,4 $\geq$ 7,2) : \textbf{normaux}. Les voies excrétrices sont vraisemblablement perméables (pas d'obstruction distale majeure).
\item Concentration : \SI{8}{M/mL} $<$ \SI{16}{M/mL} $\Rightarrow$ \textbf{oligozoospermie}.
\item Mobilité progressive : \SI{20}{\%} $<$ \SI{32}{\%} $\Rightarrow$ \textbf{asthénozoospermie}.
\item Morphologie : \SI{3}{\%} $<$ \SI{4}{\%} $\Rightarrow$ \textbf{tératozoospermie}. Vitalité normale (\SI{60}{\%} $\geq$ \SI{54}{\%}).
\end{enumerate}
Conclusion : il s'agit d'une \textbf{oligo-asthéno-tératozoospermie (OATS)} modérée à sévère. Conduite à tenir : confirmer par un second spermogramme dans 3 mois ; rechercher une cause (examen clinique à la recherche d'un varicocèle, bilan infectieux urinaire et sexuel, échographie scrotale, dosage hormonal FSH/LH/Testostérone, caryotype si OATS sévère < 5 M/mL) ; mener en parallèle le bilan de l'épouse (courbe de température ou dosage de progestérone au jour 21 pour confirmer l'ovulation, hystérosalpingographie ou hyfosy pour la perméabilité tubaire). Si l'OATS est confirmée et sévère, une PMA de type FIV-ICSI pourra être proposée.
\end{exoresolu}

\begin{attention}[Pièges du spermogramme]
Trois erreurs fréquentes : (1) conclure sur \textbf{un seul examen} — la spermatogenèse fluctue (fièvre, stress, chaleur, tabac) ; (2) oublier que la spermatogenèse dure environ \textbf{74 jours} : un résultat reflète l'état du testicule deux à trois mois plus tôt ; (3) assimiler OATS et stérilité absolue : une OATS modérée est compatible avec des grossesses spontanées, simplement moins probables.
\end{attention}

\courssub{Les techniques de procréation médicalement assistée}

La \cle{procréation médicalement assistée} (PMA) regroupe l'ensemble des techniques médicales et biologiques permettant de concevoir un enfant lorsque la voie naturelle est impossible ou inefficace. On les classe par ordre de complexité croissante.

\textbf{A. Stimulation ovarienne simple et rapports programmés}\par
On administre à la femme des inducteurs d'ovulation (citrate de clomifène, anti-œstrogène, ou gonadotrophines exogènes FSH/LH) pour provoquer la croissance d'un ou deux follicules, puis on programme les rapports au moment de l'ovulation (repérée par échographie et/ou test LH urinaire). Indication : troubles de l'ovulation (SOPK, hypogonadotrope) avec sperme normal et trompes perméables.

\textbf{B. Insémination intra-utérine (IIU ou IAC)}\par
Le sperme du conjoint est préparé au laboratoire (gradient de densité, swim-up) pour sélectionner les spermatozoïdes les plus mobiles et les concentrer, puis déposé directement dans la cavité utérine via un cathéter souple au moment de l'ovulation. Indications : OATS modérée, anomalies de la glaire cervicale (facteur cervical), infertilité inexpliquée, éjaculation rétrograde. Taux de succès : 10 à 15 \% par cycle.

\textbf{C. Fécondation in vitro (FIV) classique}\par
La \cle{FIV} consiste à réaliser la fécondation \textbf{hors du corps de la femme} (« in vitro » = dans le verre, c'est-à-dire en boîte de Pétri de culture). Ses étapes sont résumées dans le schéma ci-dessous.

\begin{popfigure}
\begin{center}
\begin{tikzpicture}[
  etape/.style={rectangle, rounded corners=3pt, draw=popblue, fill=popblueL, text width=2.8cm, minimum height=1.4cm, align=center, font=\small},
  fleche/.style={-{Stealth[length=3mm]}, very thick, poporange}
]
\node[etape] (e1) at (0,0) {\textbf{1. Stimulation} ovarienne contrôlée (gonadotrophines) : croissance multifolliculaire};
\node[etape] (e2) at (4.5,0) {\textbf{2. Déclenchement} (hCG) puis \textbf{ponction} folliculaire sous écho endovaginale};
\node[etape] (e3) at (9,0) {\textbf{3. Fécondation} : mise en contact ovocytes + spermatozoïdes (FIV) \textbf{OU} ICSI};
\node[etape] (e4) at (4.5,-2.5) {\textbf{4. Culture} embryonnaire en incubateur (3 à 5 jours, stade blastocyste idéal)};
\node[etape] (e5) at (0,-2.5) {\textbf{5. Transfert} de 1 embryon (électif) ou 2 dans l'utérus};
\node[etape, fill=popgreenL, draw=popgreen] (e6) at (9,-2.5) {\textbf{6. Nidation} et test $\beta$-hCG à J+12};
\draw[fleche] (e1) -- (e2);
\draw[fleche] (e2) -- (e3);
\draw[fleche] (e3.south) |- ++(0,-0.5) -| (e4.north);
\draw[fleche] (e4) -- (e5);
\draw[fleche] (e4.east) -- (e6.west);
\end{tikzpicture}
\end{center}
\legende{Étapes de la Fécondation In Vitro (FIV) avec ou sans ICSI : (1) Stimulation ovarienne contrôlée pour obtenir plusieurs ovocytes matures ; (2) Ponction folliculaire transvaginale sous anesthésie ; (3) Fécondation en laboratoire : FIV classique (50 000 à 100 000 spermatozoïdes par ovocyte) ou ICSI (injection d'un spermatozoïde unique) ; (4) Culture des embryons jusqu'au stade blastocyste (J5/J6) ; (5) Transfert embryonnaire indolore sous contrôle échographique ; (6) Attente de la nidation et dosage de la $\beta$-hCG.}
\end{popfigure}

\textbf{D. Injection intracytoplasmique de spermatozoïde (ICSI)}\par

\begin{propriete}[Intérêt de l'ICSI]
L'\cle{ICSI} (\emph{IntraCytoplasmic Sperm Injection}) consiste à \textbf{injecter directement un seul spermatozoïde dans le cytoplasme de l'ovocyte} à l'aide d'une micropipette de verre sous microscope inversé (grossissement $\times$200--400). Elle permet donc la fécondation \textbf{avec un unique spermatozoïde vivant}, même immobile (asthénozoospermie sévère) ou malformé (tératozoospermie sévère). Elle est indiquée en cas d'\textbf{OATS sévère} (concentration $<$ 5 M/mL), d'\textbf{azoospermie obstructive} (les spermatozoïdes sont alors prélevés chirurgicalement : TESE, MESA, PESA) ou d'échecs de FIV classiques.
\end{propriete}

\textbf{E. Techniques avec donneur (don de gamètes ou d'embryons)}\par
Lorsqu'un membre du couple ne produit pas de gamètes utilisables (azoospermie sécrétoire, insuffisance ovarienne prématurée, risque génétique majeur non accessible au DPI) : \textbf{don de spermatozoïdes}, \textbf{don d'ovocytes}, voire \textbf{don d'embryons} (embryons surnuméraires d'un autre couple ayant réussi leur PMA). La \textbf{gestation pour autrui} (GPA, ou « mère porteuse ») consiste à faire porter l'embryon par une tierce femme.

\begin{attention}[Cadre légal et éthique camerounais]
La \textbf{gestation pour autrui (GPA) est interdite par la loi camerounaise} (Loi n° 2016/007 du 12 juillet 2016 relative au Code de la Santé Publique, art. 372), comme dans la grande majorité des pays d'Afrique francophone. Tout projet de PMA doit respecter le cadre légal et éthique national : consentement éclairé écrit des deux partenaires, anonymat et gratuité du don de gamètes, interdiction de toute commercialisation des éléments du corps humain. Au Baccalauréat, ne présente jamais la GPA comme une technique librement accessible au Cameroun.
\end{attention}

\begin{exemplebox}[Efficacité de la FIV/ICSI : des chiffres à connaître]
Le taux de succès d'un cycle de FIV/ICSI (naissance d'un enfant vivant, \emph{baby take-home rate}) est d'environ \textbf{30 à 40 \% chez une femme de moins de 35 ans}, mais il \textbf{chute à moins de 10 \% après 42 ans}, en raison du vieillissement des ovocytes (augmentation des aneuploïdies). Il faut donc souvent plusieurs tentatives (3 à 4 cycles cumulés pour atteindre 60--70 \% de succès cumulé chez les < 35 ans), et l'âge de la femme est le principal facteur pronostique — d'où l'importance de consulter tôt (après 12 mois d'essais, 6 mois si femme > 35 ans).
\end{exemplebox}

\courssub{Avantages, limites et questions éthiques de la PMA}

\textbf{Avantages.} La PMA a rendu la \textbf{parentalité biologique} possible à des millions de couples stériles (plus de 10 millions d'enfants nés par FIV dans le monde depuis Louise Brown en 1978). Elle permet aussi le \textbf{diagnostic préimplantatoire} (DPI) : sur un embryon obtenu par FIV, on prélève quelques cellules (biopsie de blastomères au stade 6-8 cellules J3, ou de trophectoderme au stade blastocyste J5) pour rechercher une anomalie génétique ou chromosomique, et seuls les embryons indemnes sont transférés.

\begin{savaistu}[Le DPI contre la drépanocytose]
La \textbf{drépanocytose} (anémie falciforme), maladie monogénique récessive très fréquente au Cameroun (environ 20 \% de la population est porteuse du trait $S$), peut être évitée grâce au DPI : un couple dont les deux membres sont hétérozygotes $[A//S]$ a un risque de 25 \% d'avoir un enfant drépanocytaire $[S//S]$ à chaque grossesse. La FIV avec DPI permet de sélectionner et de transférer uniquement des embryons $[A//A]$ ou $[A//S]$, exempts de la maladie homozygote. C'est une application majeure de la génétique à la santé reproductive africaine, disponible dans quelques centres spécialisés (ex. Centre Mère-Enfant de la Fondation Chantal Biya à Yaoundé).
\end{savaistu}

\textbf{Limites et inconvénients.}
\begin{itemize}
\item \textbf{Grossesses multiples} (jumeaux, triplés) si plusieurs embryons sont transférés, avec leurs risques (prématurité, prééclampsie, faible poids de naissance) — d'où la tendance actuelle forte au \textbf{transfert d'un embryon unique} (SET, \emph{Single Embryo Transfer}) ;
\item \textbf{Syndrome d'hyperstimulation ovarienne} (SHO) : réponse excessive aux gonadotrophines (ovaires volumineux, ascite, hémoconcentration, risque thromboembolique), potentiellement grave, nécessitant une surveillance étroite ;
\item \textbf{Coût élevé} et \textbf{accès inégal} : au Cameroun, un cycle de FIV/ICSI coûte entre \SI{1,5}{million} et \SI{3}{millions} de francs CFA (hors médicaments), hors de portée de la plupart des familles, et les centres de PMA agréés sont rares (principalement Yaoundé, Douala) ;
\item \textbf{Échecs répétés} et charge psychologique lourde pour le couple (stress, anxiété, dépression, tensions conjugales).
\end{itemize}

\textbf{Questions éthiques.} La PMA soulève des débats de société que la science seule ne tranche pas :
\begin{itemize}
\item Quel \textbf{statut} accorder à l'\textbf{embryon} surnuméraire (conservation cryogénique limitée dans le temps, destruction, don à la recherche ou à un autre couple) ?
\item Faut-il maintenir l'\textbf{anonymat du don} de gamètes ou ouvrir l'accès aux origines pour l'enfant né de don ?
\item Le \textbf{consentement éclairé} des deux partenaires est-il toujours réel, notamment dans les contextes de forte pression familiale pour avoir un enfant ?
\item Le DPI ne risque-t-il pas de dériver vers une « eugénique » de sélection (choix du sexe, traits physiques, intelligence) ?
\end{itemize}
Ces questions n'ont pas de réponse scientifique unique : elles relèvent du dialogue entre science, droit, philosophie et valeurs de la société.

\courssub{Sensibilisation et prévention}

Face à la stérilité, l'action du citoyen et du futur bachelier est triple :
\begin{enumerate}
\item \textbf{Promouvoir le conseil préconceptionnel} : consultation du couple avant le mariage ou le désir d'enfant (dépistage des IST -- VIH, syphilis, hépatites B/C, chlamydies --, sérologies, statut drépanocytaire $[A//A]$, $[A//S]$ ou $[S//S]$, vaccination rubéole/hépatite B si non immunisée, supplémentation en acide folique \SI{0,4}{mg/j} pour la femme) ;
\item \textbf{Lutter contre la stigmatisation} : l'infertilité est une \textbf{maladie du couple} reconnue par l'OMS, pas une faute, ni une malédiction, ni un « mauvais sort » ; les campagnes de sensibilisation doivent rappeler que l'homme est concerné dans près d'un cas sur deux et doit accepter de faire un spermogramme sans que sa virilité ne soit remise en cause ;
\item \textbf{Promouvoir la santé reproductive masculine} : éviter tabac, alcool et chaleur excessive sur les testicules (bains chauds, saunas, sous-vêtements trop serrés, ordinateur portable sur les cuisses), traiter précocement toute infection génitale, consulter pour un varicocèle symptomatique.
\end{enumerate}

\begin{exercice}[Concevoir un outil de sensibilisation]
Rédige le texte d'une affiche destinée à un centre de santé de quartier, intitulée « L'infertilité concerne le couple », comportant : une définition simple de la stérilité, deux causes masculines et deux causes féminines, un message contre la stigmatisation, et un appel à la consultation précoce des deux partenaires.
\end{exercice}

\begin{corrige}[Exercice — éléments de réponse]
L'affiche doit contenir au minimum : (1) « Un couple est dit stérile lorsqu'aucune grossesse n'est obtenue après 12 mois de rapports réguliers non protégés » ; (2) causes masculines : azoospermie, OATS, varicocèle, infections génitales ; causes féminines : troubles de l'ovulation (SOPK), trompes bouchées (salpingite), fibromes utérins, âge > 35 ans ; (3) message : « L'infertilité n'est ni une faute ni une malédiction : c'est un problème médical du couple, dont la cause est masculine dans près d'un cas sur deux » ; (4) appel : « Consultez ensemble, dès 12 mois d'essais infructueux (6 mois si la femme a plus de 35 ans) : le spermogramme est un examen simple, non invasif, et des solutions existent (traitements médicaux, chirurgie, insémination, FIV/ICSI) ». La correction valorise la clarté, l'exactitude scientifique (valeurs OMS 2021) et le ton respectueux, non culpabilisant et inclusif.
\end{corrige}

\begin{aretenir}[L'essentiel de la section 5]
\begin{itemize}
\item Stérilité = absence de grossesse après \textbf{12 mois} de rapports réguliers non protégés ; la cause est masculine dans 30 à 40 \% des cas : le bilan concerne \textbf{toujours les deux partenaires} simultanément.
\item Le \textbf{spermogramme} s'interprète selon l'OMS 2021 : concentration $\geq$ \SI{16}{M/mL}, mobilité progressive $\geq$ \SI{32}{\%}, formes normales $\geq$ \SI{4}{\%} (Kruger), vitalité $\geq$ \SI{54}{\%} ; toute anomalie se confirme sur un second examen à 3 mois (durée de la spermatogenèse).
\item La PMA va du plus simple au plus complexe : stimulation ovarienne + rapports programmés, \textbf{IIU}, \textbf{FIV}, \textbf{ICSI} (un seul spermatozoïde vivant suffit), don de gamètes. La \textbf{GPA est interdite au Cameroun}.
\item Succès de la FIV/ICSI : 30 à 40 \% par cycle avant 35 ans, $<$ 10 \% après 42 ans. Le \textbf{DPI} (biopsie blastocyste J5) permet d'éviter la transmission de maladies monogéniques comme la drépanocytose.
\item La PMA pose des questions éthiques (statut de l'embryon, anonymat du don, inégalités d'accès, dérive eugénique) et impose une démarche de \textbf{sensibilisation} et de lutte contre la stigmatisation, notamment masculine.
\end{itemize}
\end{aretenir}

\newpage

\coursec{Activité d'intégration}

\courssub{Objectifs de l'activité}

À l'issue de cette activité d'intégration, tu seras capable de :
\begin{itemize}
\item mobiliser simultanément plusieurs savoir-faire de la leçon : interpréter des expériences de castration et d'injection, analyser des courbes hormonales, interpréter un spermogramme, choisir une technique de procréation médicalement assistée ;
\item identifier un \cle{rétrocontrôle négatif} (rétroaction négative) à partir de données expérimentales ;
\item expliquer le \cle{mécanisme d'action d'une pilule contraceptive} à partir de l'évolution des taux hormonaux ;
\item proposer une solution médicale adaptée à un cas d'infertilité masculine ;
\item concevoir un outil de sensibilisation destiné au grand public camerounais, en adoptant un vocabulaire scientifique exact et une attitude de respect et de non-stigmatisation.
\end{itemize}

\courssub{Situation}

\textbf{Dossier clinique n\textsuperscript{o} 2024-117 — Centre Hospitalier d'Essos, Yaoundé.}

M. et Mme \textbf{Ngono}, mariés depuis six ans, habitent le quartier Nkol-Eton à Yaoundé. Ils désirent un enfant, mais aucune grossesse n'est survenue malgré des rapports réguliers non protégés depuis plus de deux ans. Le médecin gynécologue du couple a constitué un dossier comportant quatre documents.

\textbf{Document 1 — Courbes hormonales de Mme Ngono.}

Le taux sanguin de LH (hormone hypophysaire) et d'{\oe}stradiol (hormone ovarienne) a été mesuré chez Mme Ngono pendant deux mois consécutifs. Le premier mois, elle ne prenait aucun contraceptif ; le second mois, elle prenait une pilule {\oe}stroprogestative (comprimé quotidien contenant de l'{\oe}stradiol et de la progestérone de synthèse).

\begin{center}
\begin{tabularx}{\linewidth}{|l|X|X|}
\hline
\rowcolor{popblueL} \textbf{Paramètre mesuré} & \textbf{Mois 1 (cycle naturel)} & \textbf{Mois 2 (sous pilule)} \\
\hline
Pic d'{\oe}stradiol & présent vers le 12\textsuperscript{e} jour & absent (taux bas et constant) \\
\hline
Pic de LH (pic ovulatoire) & présent vers le 14\textsuperscript{e} jour & absent \\
\hline
Pic de progestérone (2\textsuperscript{e} partie du cycle) & présent vers le 21\textsuperscript{e} jour & absent \\
\hline
Ovulation (échographie) & oui & non \\
\hline
\end{tabularx}
\end{center}

\textbf{Document 2 — Expériences réalisées chez le rat mâle (laboratoire de physiologie animale).}

\begin{center}
\begin{tabularx}{\linewidth}{|l|X|X|}
\hline
\rowcolor{popblueL} \textbf{Lot} & \textbf{Traitement} & \textbf{Résultats} \\
\hline
A (témoin) & aucun & testicules normaux ; testostéronémie normale \\
\hline
B & castration (ablation des testicules) & testostéronémie effondrée ; taux de LH et FSH \textbf{augmentés} ; atrophie des organes annexes (vésicules séminales) \\
\hline
C & castration + injection quotidienne de testostérone & taux de LH et FSH \textbf{revenus à la normale} ; organes annexes maintenus \\
\hline
D & rat intact + injections massives de testostérone & taux de LH et FSH \textbf{effondrés} ; réduction de la spermatogenèse \\
\hline
\end{tabularx}
\end{center}

\textbf{Document 3 — Spermogramme de M. Ngono.}

\begin{center}
\begin{tabularx}{\linewidth}{|l|X|X|}
\hline
\rowcolor{popblueL} \textbf{Paramètre} & \textbf{Résultat de M. Ngono} & \textbf{Valeurs de référence (OMS)} \\
\hline
Volume de l'éjaculat & \SI{3,0}{mL} & $\geq$ \SI{1,5}{mL} \\
\hline
Concentration en spermatozoïdes & \SI{4}{millions/mL} & $\geq$ \SI{15}{millions/mL} \\
\hline
Mobilité progressive & 18 \% & $\geq$ 32 \% \\
\hline
Formes typiques (morphologie normale) & 3 \% & $\geq$ 4 \% \\
\hline
\end{tabularx}
\end{center}

\textbf{Document 4 — Fiche d'information sur la PMA (service de médecine de la reproduction).}

\begin{itemize}
\item \textbf{Insémination intra-utérine (IIU)} : dépôt de spermatozoïdes sélectionnés dans l'utérus ; indiquée si la concentration en spermatozoïdes mobiles reste suffisante.
\item \textbf{FIV classique} (fécondation in vitro) : mise en contact d'ovocytes et de spermatozoïdes en laboratoire ; il faut au moins quelques centaines de milliers de spermatozoïdes mobiles par ovocyte.
\item \textbf{ICSI} (injection intra-cytoplasmique d'un spermatozoïde) : un seul spermatozoïde est injecté directement dans l'ovocyte ; indiquée dans les oligo-asthéno-tératozoospermies sévères, même avec très peu de spermatozoïdes.
\end{itemize}

Le couple est bouleversé : dans leur entourage, on murmure que « c'est toujours la faute de la femme ». Le médecin souhaite au contraire expliquer scientifiquement la situation au couple et sensibiliser la communauté.

\courssub{Consignes}

Par groupes de quatre élèves, traitez les tâches suivantes et préparez une restitution orale de dix minutes.

\begin{enumerate}
\item À partir du document 2, dégager le type de régulation illustré par les lots B, C et D. Nommer précisément les organes et les hormones impliqués et construire le \cle{schéma fonctionnel} de la régulation du taux de testostérone (avec les flèches $+$ et $-$).
\item À partir du document 1, expliquer pourquoi, sous pilule, le pic de LH disparaît et l'ovulation est bloquée. Justifier en utilisant la notion de rétrocontrôle.
\item Interpréter le spermogramme de M. Ngono (document 3) : comparer chaque paramètre aux valeurs de référence, poser le diagnostic biologique et proposer une cause probable parmi celles étudiées en cours.
\item En vous appuyant sur les documents 3 et 4, choisir la technique de PMA la plus adaptée au couple Ngono et justifier ce choix en écartant les autres techniques.
\item Concevoir une affiche A4 de sensibilisation intitulée « L'infertilité masculine existe : parlons-en ! », destinée aux jeunes couples camerounais. Elle devra comporter : un message corrigeant l'idée reçue selon laquelle l'infertilité est toujours féminine, deux causes évitables d'infertilité masculine, un conseil de prévention et une invitation à consulter.
\end{enumerate}

\courssub{Ressources à mobiliser}

\begin{aretenir}[Ce qu'il faut avoir en tête]
\begin{itemize}
\item Chez l'homme : l'\cle{hypothalamus} libère la \cle{GnRH}, qui stimule l'\cle{hypophyse} ; celle-ci sécrète la \cle{LH} (stimule les cellules de Leydig $\rightarrow$ testostérone) et la \cle{FSH} (agit sur les tubes séminifères $\rightarrow$ spermatogenèse). La testostérone exerce un \cle{rétrocontrôle négatif} sur l'axe hypothalamo-hypophysaire : quand son taux monte, LH et FSH baissent, et inversement.
\item Chez la femme : la \cle{FSH} provoque la croissance folliculaire (sécrétion d'{\oe}strogènes) ; le \cle{pic de LH} déclenche l'ovulation ; la progestérone du corps jaune exerce un rétrocontrôle négatif.
\item La \cle{pilule {\oe}stroprogestative} maintient des taux constants d'{\oe}strogènes et de progestérone qui bloquent, par rétrocontrôle négatif, la sécrétion de FSH et de LH : pas de pic de LH, donc \textbf{pas d'ovulation}.
\item Spermogramme : \cle{oligozoospermie} (concentration faible), \cle{asthénozoospermie} (mobilité faible), \cle{tératozoospermie} (formes anormales).
\item L'\cle{ICSI} est la technique de choix quand les spermatozoïdes sont très rares ou peu mobiles, car un seul spermatozoïde suffit par ovocyte.
\end{itemize}
\end{aretenir}

\courssub{Démarche et corrigé commenté}

\begin{exoresolu}[Corrigé commenté du dossier Ngono]
\textbf{Tâche 1 — Type de régulation et schéma fonctionnel.}

\tcblower

\textbf{Analyse des lots.} Chez le lot B (castré), la testostérone disparaît et les taux de LH/FSH \emph{augmentent} : les testicules exercent donc normalement un frein sur l'hypophyse. Chez le lot C (castré + testostérone), LH et FSH reviennent à la normale : c'est bien la testostérone qui exerce ce frein. Chez le lot D (intact + excès de testostérone), LH et FSH s'effondrent : confirmation. Il s'agit d'un \cle{rétrocontrôle négatif} (rétroaction négative) de la testostérone sur le complexe hypothalamo-hypophysaire.

\textbf{Schéma fonctionnel attendu :}
\begin{popfigure}
\begin{tikzpicture}[node distance=1.1cm, every node/.style={draw, rounded corners, align=center, font=\small}]
\node[fill=popblueL, align=center] (hypo){Hypothalamus\\ (GnRH)};
\node[fill=poppurpleL, below=of hypo, align=center] (hyp){Hypophyse\\ (LH et FSH)};
\node[fill=popgreenL, below=of hyp, align=center] (test){Testicules\\ (testostérone ; spermatogenèse)};
\draw[-{Stealth}, thick, popblue] (hypo) -- node[right, draw=none]{$+$} (hyp);
\draw[-{Stealth}, thick, popblue] (hyp) -- node[right, draw=none]{$+$} (test);
\draw[-{Stealth}, thick, poporange] (test.east) to[bend left=40] node[right, draw=none]{$-$ (rétrocontrôle)} (hyp.east);
\draw[-{Stealth}, thick, poporange] (test.west) to[bend right=40] node[left, draw=none]{$-$} (hypo.west);
\end{tikzpicture}
\legende{Schéma fonctionnel de la régulation du taux de testostérone : la GnRH hypothalamique stimule ($+$) l'hypophyse ; la LH et la FSH stimulent ($+$) les testicules ; la testostérone inhibe ($-$) l'hypophyse et l'hypothalamus.}
\end{popfigure}

\textbf{Tâche 2 — Effet de la pilule sur les courbes de Mme Ngono.}

Sous pilule, les {\oe}strogènes et la progestérone de synthèse sont apportés chaque jour à taux constant. Ces hormones exercent un rétrocontrôle négatif permanent sur l'axe hypothalamo-hypophysaire : la sécrétion de FSH et de LH reste bloquée à un niveau bas. Sans FSH, pas de maturation folliculaire complète (pas de pic d'{\oe}stradiol) ; surtout, \textbf{le pic de LH du 14\textsuperscript{e} jour est supprimé}. Or c'est ce pic qui déclenche l'ovulation : il n'y a donc \textbf{ni ovulation, ni corps jaune}, d'où l'absence du pic de progestérone au 21\textsuperscript{e} jour. La pilule est ainsi un contraceptif \emph{anovulatoire}.

\textbf{Tâche 3 — Interprétation du spermogramme.}

Comparaison aux normes OMS :
\begin{itemize}
\item concentration : \SI{4}{millions/mL} $<$ \SI{15}{millions/mL} $\rightarrow$ \cle{oligozoospermie} ;
\item mobilité progressive : 18 \% $<$ 32 \% $\rightarrow$ \cle{asthénozoospermie} ;
\item formes typiques : 3 \% $<$ 4 \% $\rightarrow$ \cle{tératozoospermie}.
\end{itemize}
Le diagnostic est une \cle{oligo-asthéno-tératozoospermie} (OATS) modérée à sévère : l'infertilité du couple a donc une \textbf{composante masculine}. Causes probables à évoquer : varicocèle (dilatation des veines du cordon spermatique élevant la température du testicule), infection uro-génitale mal traitée (urétrite à \emph{Chlamydia} ou blennorragie négligée), exposition à la chaleur ou à des toxiques (tabac, alcool, pesticides), ou trouble hormonal de l'axe hypothalamo-hypophysaire.

\textbf{Tâche 4 — Choix de la technique de PMA.}

\begin{itemize}
\item \textbf{IIU écartée} : elle exige un nombre suffisant de spermatozoïdes mobiles ; avec 18 \% de mobilité sur seulement \SI{4}{millions/mL}, le réservoir de spermatozoïdes mobiles est trop faible.
\item \textbf{FIV classique écartée} : la fécondation « spontanée » en éprouvette requiert des centaines de milliers de spermatozoïdes mobiles par ovocyte, condition non remplie ici.
\item \textbf{ICSI retenue} : un seul spermatozoïde vivant, même peu mobile, est injecté directement dans le cytoplasme de l'ovocyte. C'est la seule technique compatible avec une OATS sévère.
\end{itemize}
On précisera au couple les limites : coût élevé, taux de réussite par tentative de l'ordre de 30 à 40 \%, questions éthiques (conservation des embryons, don de gamètes) à discuter librement.

\textbf{Tâche 5 — Affiche de sensibilisation (attendus).}

L'affiche doit comporter : un titre accrocheur (« L'infertilité masculine existe : parlons-en ! ») ; le message clé « Dans environ un couple infertile sur deux, un facteur masculin est en cause : ce n'est pas toujours la faute de la femme » ; deux causes évitables (infections sexuellement transmissibles non traitées ; tabac, alcool, chaleur excessive, pesticides) ; un conseil de prévention (se protéger des IST, consulter rapidement en cas d'infection) ; une invitation au dépistage (« Un simple spermogramme au centre hospitalier peut tout changer : consultez ensemble »). Elle doit respecter la dignité des personnes et éviter toute stigmatisation.
\end{exoresolu}

\courssub{Bilan de l'activité}

Cette activité a permis de mettre en évidence que :
\begin{itemize}
\item la régulation hormonale des fonctions reproductrices, chez l'homme comme chez la femme, repose sur un même principe : le \cle{rétrocontrôle négatif} exercé par les hormones gonadiques sur le complexe hypothalamo-hypophysaire — principe que les expériences de castration et d'injection permettent de démontrer rigoureusement ;
\item la contraception hormonale n'est qu'une \emph{application} de ce principe : la pilule « trompe » l'axe hypothalamo-hypophysaire en maintenant un rétrocontrôle permanent qui supprime le pic ovulatoire de LH ;
\item l'infertilité d'un couple peut avoir une cause masculine, féminine ou mixte : seul un bilan biologique (courbes hormonales, spermogramme) permet un diagnostic objectif, loin des préjugés ;
\item la procréation médicalement assistée, notamment l'ICSI, offre des solutions réelles aux couples camerounais concernés, tout en posant des questions éthiques qu'il faut aborder avec honnêteté ;
\item la sensibilisation est un acte de santé publique : informer sans stigmatiser, c'est déjà soigner.
\end{itemize}

\begin{attention}[Erreurs fréquentes à éviter dans ce type de dossier]
Ne confonds pas contraception (empêche l'ovulation ou la fécondation, avant la conception) et contragestion (empêche la nidation, après la fécondation). Ne dis jamais que « la pilule contient de la LH » : elle contient des {\oe}strogènes et de la progestérone de synthèse. Enfin, dans un spermogramme, compare toujours chaque paramètre à la norme avec son unité avant de conclure.
\end{attention}

\newpage

\coursec{Méthodes et exercices résolus}

\courssub{Méthode 1 : Interpréter les expériences de castration, greffes et ablations}

\begin{methode}[Analyser une expérience de castration, d'injection d'extraits ou de greffe]
Face à un tableau de résultats expérimentaux (castration, injection d'extraits testiculaires ou ovariens, ablation ou greffe de gonades ou d'hypophyse, section de la tige pituitaire, parabiose), adopte toujours la même démarche :
\begin{enumerate}
\item \textbf{Identifier le témoin et l'expérience} : repérer l'animal témoin (non opéré ou opéré de façon fictive \emph{sham}) qui sert de référence, et l'animal expérimental ayant subi la manipulation.
\item \textbf{Décrire les résultats quantitativement} : comparer les paramètres mesurés (développement des organes génitaux accessoires, taux d'hormones en ng/mL ou UI/L, comportement sexuel, caractères sexuels secondaires) entre témoin et expérimental, en citant les chiffres avec leurs unités.
\item \textbf{Dégager la relation de cause à effet} : la seule variable qui change entre les deux lots est la manipulation ; la différence observée est donc due à l'organe ou à la substance retirée/injectée.
\item \textbf{Conclure sur le rôle de la structure} : « l'ablation de X supprime Y, l'injection d'extrait de X (ou la greffe de X) restaure Y, donc X sécrète une substance (hormone) responsable de Y ».
\item \textbf{Préciser le mode d'action} : si la greffe fonctionne même en position ectopique (sous la peau, dans la cavité abdominale), la substance agit par voie \cle{sanguine} (générale), ce qui prouve sa nature \cle{hormonale}.
\end{enumerate}
\textbf{Réflexes} : ne jamais conclure sans témoin ; toujours vérifier que l'effet est réversible par l'injection ou la greffe ; pour la section de la tige pituitaire, distinguer ce qui passe par le sang (axes hormonaux : gonades, thyroïde, corticosurrénales) de ce qui passe par les nerfs (médullosurrénale, contraction utérine, éjection lactée).
\end{methode}

\begin{exoresolu}[Castration et injection d'extrait testiculaire chez le coq]
On réalise trois lots de jeunes coqs. Lot A : témoins non opérés. Lot B : castrés (ablation des testicules). Lot C : castrés puis recevant chaque jour une injection d'extrait testiculaire. Après 8 semaines, on mesure la hauteur de la crête et on observe le comportement.

\begin{center}
\begin{tabularx}{\linewidth}{|l|X|X|X|}
\hline
\rowcolor{popblueL} & Lot A (témoin) & Lot B (castré) & Lot C (castré + extrait) \\
\hline
Hauteur de la crête (cm) & 6,2 & 1,8 & 5,9 \\
\hline
Comportement sexuel & normal & aboli & normal \\
\hline
\end{tabularx}
\end{center}

Interpréter ces résultats et conclure sur le rôle du testicule.
\tcblower
\textbf{1. Comparaison des lots.} Le lot B (castré) présente une crête atrophiée (1,8 cm contre 6,2 cm chez le témoin A, soit une réduction d'environ 71 \%) et une abolition du comportement sexuel. La seule différence entre A et B est l'ablation des testicules : la castration est donc la cause de la régression de la crête et de la disparition du comportement sexuel.

\textbf{2. Effet de l'injection.} Le lot C, bien que castré, présente une crête quasi normale (5,9 cm) et un comportement sexuel normal. L'injection d'extrait testiculaire a donc \emph{suppléé} l'absence des testicules : l'extrait contient la substance active normalement produite par le testicule.

\textbf{3. Interprétation.} Le testicule libère dans le sang une substance chimique — une \cle{hormone}, la \cle{testostérone} — qui assure le développement et l'entretien des caractères sexuels secondaires (crête, barbillons, chant) et le comportement sexuel. Comme l'extrait injecté par voie générale agit à distance sur la crête, l'action est bien \cle{humorale} (par voie sanguine).

\textbf{Conclusion} : le testicule exerce, par l'intermédiaire de la testostérone qu'il sécrète, un contrôle hormonal sur les caractères sexuels secondaires et le comportement sexuel du coq.
\end{exoresolu}

\courssub{Méthode 2 : Tracer et analyser les courbes hormonales du cycle ovarien}

\begin{methode}[Analyser les courbes du cycle sexuel de la femme]
\begin{enumerate}
\item \textbf{Repérer les axes} : abscisse = temps en jours (cycle de 28 jours en moyenne, le jour 1 étant le premier jour des règles) ; ordonnée = taux sanguins des hormones (estradiol en pg/mL, progestérone en ng/mL, FSH et LH en UI/L).
\item \textbf{Identifier les pics caractéristiques} : pic de \cle{LH} (et petit pic de FSH) vers le 14\textsuperscript{e} jour $\Rightarrow$ déclenchement de l'\cle{ovulation} ; pic d'\cle{estradiol} juste avant (fin de phase folliculaire, vers J13) ; plateau de \cle{progestérone} en phase lutéale (jours 15 à 26, maximum vers J21).
\item \textbf{Relier hormones et événements ovariens} : FSH $\rightarrow$ croissance des follicules ; follicules $\rightarrow$ sécrétion d'estradiol ; pic de LH $\rightarrow$ ovulation + lutéinisation ; corps jaune $\rightarrow$ progestérone + estradiol.
\item \textbf{Relier hormones et événements utérins} : estradiol $\rightarrow$ phase proliférative de l'endomètre ; progestérone $\rightarrow$ phase sécrétoire (endomètre prêt à la nidation) ; chute brutale des hormones ovariennes $\rightarrow$ règles (désquamation).
\item \textbf{Expliquer les rétrocontrôles} : estradiol faible (début cycle) $\rightarrow$ rétrocontrôle négatif sur FSH/LH ; estradiol élevé et prolongé (fin phase folliculaire) $\rightarrow$ rétrocontrôle \emph{positif} (déclenchement du pic de LH) ; progestérone + estradiol en phase lutéale $\rightarrow$ rétrocontrôle négatif puissant bloquant FSH et LH (empêche une nouvelle ovulation).
\end{enumerate}
\textbf{Réflexe} : pour dater l'ovulation, chercher le pic de LH, pas le pic d'estradiol. Le pic d'estradiol \emph{précède} et \emph{provoque} le pic de LH.
\end{methode}

\begin{exoresolu}[Analyse des courbes hormonales d'un cycle de 28 jours]
Les dosages sanguins quotidiens chez une femme donnent : estradiol maximal au 13\textsuperscript{e} jour (environ 300 pg/mL) ; LH maximale au 14\textsuperscript{e} jour (environ 60 UI/L, contre 8 UI/L en début de cycle) ; progestérone maximale au 21\textsuperscript{e} jour (environ 15 ng/mL) ; chute de toutes les hormones ovariennes au 27\textsuperscript{e} jour. (1) Dater l'ovulation en justifiant. (2) Expliquer l'origine du pic de LH. (3) Expliquer l'apparition des règles.
\tcblower
\textbf{(1) Datation de l'ovulation.} L'ovulation survient au 14\textsuperscript{e} jour, car elle est déclenchée par le pic de LH, qui atteint 60 UI/L ce jour-là (multiplication par 7,5 du taux basal de 8 UI/L). Le pic d'estradiol du 13\textsuperscript{e} jour le précède immédiatement : il en est la cause (voir question 2).

\textbf{(2) Origine du pic de LH.} En fin de phase folliculaire, le follicule de De Graaf mûr sécrète massivement de l'estradiol (300 pg/mL). Ce taux élevé et maintenu pendant environ 48 h exerce un \cle{rétrocontrôle positif} sur le complexe hypothalamo-hypophysaire : l'hypothalamus augmente la fréquence des pulsations de GnRH et l'antéhypophyse répond par une décharge massive de LH (et de FSH). Ce pic de LH provoque la rupture du follicule, c'est-à-dire l'ovulation, environ 24 à 36 h après le début du pic.

\textbf{(3) Apparition des règles.} En l'absence de fécondation, le corps jaune régresse (corps blanc) vers le 26\textsuperscript{e} jour ; il cesse donc de sécréter progestérone et estradiol, dont les taux chutent brutalement au 27\textsuperscript{e} jour. Or la progestérone entretenait l'endomètre sécrétoire : sa chute provoque la vasoconstriction des artères spirales, l'ischémie et la désquamation de la muqueuse utérine, d'où les règles au 28\textsuperscript{e} jour. La chute des hormones ovariennes lève aussi le rétrocontrôle négatif sur l'hypophyse : la FSH remonte légèrement et un nouveau cycle débute.

\textbf{Conclusion} : le cycle ovarien est piloté par les hormones hypophysaires (FSH, LH), elles-mêmes contrôlées par rétroactions des hormones ovariennes sur le complexe hypothalamo-hypophysaire.
\end{exoresolu}

\courssub{Méthode 3 : Réaliser un schéma fonctionnel de régulation hormonale}

\begin{methode}[Construire un schéma de régulation (testostérone ou hormones ovariennes)]
\begin{enumerate}
\item \textbf{Placer les trois étages} dans l'ordre vertical : hypothalamus (neurohormone : GnRH) $\rightarrow$ antéhypophyse (gonadostimulines : LH et FSH) $\rightarrow$ gonade (testicule : testostérone + inhibine ; ovaire : estradiol, progestérone + inhibine).
\item \textbf{Flécher les stimulations} de haut en bas avec le signe « + » et le nom de chaque hormone sur la flèche.
\item \textbf{Ajouter les rétroactions} : flèches remontantes de la gonade vers l'hypophyse \emph{et} vers l'hypothalamus, annotées « $-$ » pour la rétroaction négative (cas général : testostérone, inhibine ; progestérone + estradiol en phase lutéale) ou « + » pour le rétrocontrôle positif de l'estradiol en fin de phase folliculaire.
\item \textbf{Indiquer les effets périphériques} : testostérone $\rightarrow$ spermatogenèse (via cellules de Sertoli stimulées par FSH) et caractères sexuels ; hormones ovariennes $\rightarrow$ cycle utérin, caractères sexuels secondaires, métabolisme osseux.
\item \textbf{Vérifier la cohérence} : toute boucle fermée doit être fonctionnelle (si la gonade est stimulée, son hormone doit freiner l'étage supérieur, sauf rétrocontrôle positif ponctuel).
\end{enumerate}
\textbf{Réflexe} : chez l'homme, la LH stimule les cellules de Leydig (testostérone) et la FSH les cellules de Sertoli (spermatogenèse, sécrétion d'inhibine) ; ne pas inverser. Chez la femme, LH et FSH agissent en synergie sur le follicule (cellules de la thèque : androgènes $\rightarrow$ cellules de la granulosa : estradiol sous action FSH).
\end{methode}

\begin{exoresolu}[Schéma de la régulation du taux de testostérone chez l'homme]
Un homme reçoit par erreur des injections régulières de testostérone. On observe après quelques semaines une baisse des taux sanguins de LH et de FSH et une réduction de la spermatogenèse. (1) Réaliser le schéma fonctionnel de la régulation du taux de testostérone chez l'homme sain. (2) Expliquer les observations à l'aide du schéma.
\tcblower
\textbf{(1) Schéma fonctionnel.}

\begin{popfigure}
\begin{center}
\begin{tikzpicture}[
  box/.style={draw=popblue, thick, rounded corners, fill=popblueL, minimum width=4.5cm, minimum height=1.0cm, align=center, font=\small},
  stim/.style={->, thick, popgreen},
  inhib/.style={->, thick, poporange, dashed}]
\node[box, align=center] (hypo) at (0,4.5){Hypothalamus\\ (GnRH)};
\node[box, align=center] (hyp) at (0,2.5){Antéhypophyse\\ (LH et FSH)};
\node[box, align=center] (test) at (0,0.5){Testicule\\ (Testostérone, Inhibine)};
\draw[stim] (hypo) -- node[right, font=\small]{+ GnRH} (hyp);
\draw[stim] (hyp) -- node[right, font=\small]{+ LH, FSH} (test);
\draw[inhib] (test.west) .. controls (-3.8,1.5) .. node[left, font=\small, align=center]{$-$ Testostérone\\ $-$\ Inhibine\\ (Rétroaction négative)} (hyp.west);
\draw[inhib] (test.west) .. controls (-4.8,2.5) .. (hypo.west);
\node[font=\small, align=left, text=popdark] at (5.0,0.5) {Effets cibles :\\ Spermatogenèse (FSH + Testo)\\ Caractères sexuels secondaires\\ Métabolisme protéique, osseux};
\draw[->, popmuted] (test.east) -- (3.0,0.5);
\end{tikzpicture}
\end{center}
\legende{Schéma fonctionnel de la régulation du taux de testostérone. Flèches vertes pleines : stimulations (+) ; flèches orange pointillées : rétroaction négative ($-$) de la testostérone et de l'inhibine sur l'antéhypophyse et l'hypothalamus.}
\end{popfigure}

\textbf{(2) Explication.} La testostérone injectée s'ajoute à celle produite par le testicule : son taux sanguin devient anormalement élevé. Or la testostérone exerce une \cle{rétroaction négative} sur l'hypothalamus (diminution de la fréquence des pulsations de GnRH) et sur l'antéhypophyse (diminution de la sensibilité au GnRH). Il en résulte une chute des sécrétions de LH et de FSH, ce qui correspond aux dosages observés. La FSH étant indispensable à l'activité des cellules de Sertoli qui nourrissent les spermatozoïdes en formation, sa baisse entraîne le ralentissement de la spermatogenèse ; la baisse de LH réduit en outre la sécrétion testiculaire propre de testostérone (atrophie testiculaire).

\textbf{Conclusion} : les observations s'expliquent par le rétrocontrôle négatif de la testostérone exogène sur le complexe hypothalamo-hypophysaire ; c'est le même principe qui est exploité par la contraception hormonale masculine expérimentale (androgène seul ou androgène + progestatif).
\end{exoresolu}

\begin{exoresolu}[Schéma de la régulation des hormones ovariennes chez la femme]
(1) Réaliser le schéma fonctionnel de la régulation du cycle ovarien chez la femme (phase folliculaire et phase lutéale). (2) Indiquer sur le schéma la nature des rétrocontrôles (négatif / positif) et leur moment d'action.
\tcblower
\textbf{(1) Schéma fonctionnel.}

\begin{popfigure}
\begin{center}
\begin{tikzpicture}[
  box/.style={draw=poppurple, thick, rounded corners, fill=poppurpleL, minimum width=4.5cm, minimum height=1.0cm, align=center, font=\small},
  stim/.style={->, thick, popgreen},
  inhib/.style={->, thick, poporange, dashed},
  posfb/.style={->, thick, poppink}]
\node[box, align=center] (hypo) at (0,5){Hypothalamus\\ (GnRH)};
\node[box, align=center] (hyp) at (0,3){Antéhypophyse\\ (LH et FSH)};
\node[box, align=center] (ov) at (0,1){Ovaire\\ (Follicule / Corps jaune)\\ Estradiol, Progestérone, Inhibine};
\draw[stim] (hypo) -- node[right, font=\small]{+ GnRH} (hyp);
\draw[stim] (hyp) -- node[right, font=\small]{+ FSH, LH} (ov);
% Negative feedback (general)
\draw[inhib] (ov.west) .. controls (-4.0,2) .. node[left, font=\small, align=center]{$-$ Estradiol (faible)\\ $-$\ Progestérone\\ $-$\ Inhibine\\ (Rétrocontrôle négatif)} (hyp.west);
\draw[inhib] (ov.west) .. controls (-5.0,3) .. (hypo.west);
% Positive feedback (specific)
\draw[posfb] (ov.east) .. controls (4.0,2) .. node[right, font=\small, align=center]{$+$ Estradiol (élevé,\\ prolongé J12-J13)\\ (Rétrocontrôle positif)} (hyp.east);
\node[font=\small, align=left, text=popdark] at (5.5,1) {Effets cibles :\\ Endomètre (prolifératif/secrétoire)\\ Caractères sexuels secondaires\\ Os, lipides, coagulation};
\draw[->, popmuted] (ov.east) -- (3.5,1);
% Legend for feedback
\node[font=\tiny, text=poporange] at (-5.5, 2) {Phase lutéale \& début cycle};
\node[font=\tiny, text=poppink] at (5.5, 2) {Fin phase folliculaire (J13)};
\end{tikzpicture}
\end{center}
\legende{Schéma fonctionnel de la régulation ovarienne. Vert : stimulations. Orange pointillé : rétrocontrôle négatif permanent (inhibine, estradiol bas, progestérone). Rose : rétrocontrôle positif transitoire de l'estradiol élevé (J13) déclenchant le pic de LH.}
\end{popfigure}

\textbf{(2) Précisions.} En phase folliculaire (J1-J13), l'estradiol croissant exerce d'abord un rétrocontrôle négatif (FSH baisse). Au seuil critique (environ 200-300 pg/mL pendant > 48h), il bascule en rétrocontrôle \emph{positif} sur l'hypophyse (sensibilisation au GnRH) $\rightarrow$ pic LH/J14. En phase lutéale (J15-J26), le corps jaune sécrète progestérone + estradiol + inhibine : rétrocontrôle négatif puissant $\rightarrow$ FSH/LH bas (pas de nouvelle ovulation). À la regression du corps jaune (J26), chute des hormones $\rightarrow$ levée du frein $\rightarrow$ remontée FSH $\rightarrow$ nouveau cycle.
\end{exoresolu}

\courssub{Méthode 4 : Expliquer l'action des pilules contraceptives et contragestives}

\begin{methode}[Contraception et contragestion : raisonnement type]
\begin{enumerate}
\item \textbf{Identifier la composition} : pilule \cle{estroprogestative} (œstrogène de synthèse \emph{éthinylestradiol} + progestatif, prise 21 jours sur 28) ou pilule \cle{progestative} seule (microprogestatif, prise continue) ; contragestif (RU 486 = \cle{mifépristone}, anti-progestérone compétitif).
\item \textbf{Relier au rétrocontrôle} : les hormones de synthèse maintiennent des taux élevés et constants $\rightarrow$ rétroaction négative permanente sur le complexe hypothalamo-hypophysaire $\rightarrow$ blocage de la sécrétion de FSH et LH.
\item \textbf{Déduire les conséquences ovariennes} : pas de pic de LH $\rightarrow$ pas d'ovulation ; pas de croissance folliculaire (FSH bloquée) $\rightarrow$ pas de sécrétion d'estradiol ovarien ni de corps jaune.
\item \textbf{Citer les effets périphériques} : glaire cervicale épaissie, visqueuse, infranchissable aux spermatozoïdes (effet progestatif) ; atrophie de l'endomètre (effet progestatif) défavorable à la nidation ; motricité tubaire modifiée.
\item \textbf{Expliquer les règles artificielles (pilule estroprogestative)} : l'arrêt de la pilule pendant 7 jours (ou prise de comprimés placebo) provoque une chute brutale des hormones $\rightarrow$ hémorragie de privation (fausses règles), sans ovulation préalable.
\item \textbf{Contragestion (RU 486)} : prise unique (600 mg) dans les 7 semaines d'aménorrhée. Bloque les récepteurs nucléaires de la progestérone $\rightarrow$ l'endomètre n'est plus entretenu $\rightarrow$ décollement de l'œuf (blastocyste) et expulsion par contractions utérines (associée à des prostaglandines : misoprostol). Interruption de grossesse précoce, cadre légal et médical strict (Cameroun : loi restrictive, cas précis).
\end{enumerate}
\textbf{Réflexe} : distinguer clairement \cle{contraception} (empêche la fécondation, en amont : pilule, stérilet, préservatif) et \cle{contragestion} (agit après la fécondation et la nidation, en aval : RU 486, stérilet au cuivre en urgence). Ne jamais écrire que la pilule « détruit l'embryon ».
\end{methode}

\begin{exoresolu}[Courbes hormonales chez une utilisatrice de pilule estroprogestative]
Chez une femme prenant une pilule estroprogestative (21 comprimés actifs puis arrêt de 7 jours), les dosages sanguins montrent : FSH et LH constamment basses (moins de 5 UI/L), aucun pic de LH, estradiol ovarien bas (< 50 pg/mL), pas de corps jaune (progestérone < 1 ng/mL). Des saignements apparaissent 2 à 3 jours après l'arrêt des comprimés. Expliquer ces observations et le mode d'action contraceptif de la pilule.
\tcblower
\textbf{1. Blocage de l'axe hypothalamo-hypophysaire.} Les hormones de synthèse absorbées quotidiennement (éthinylestradiol + progestatif) maintiennent des taux sanguins élevés et stables. Elles exercent en permanence une \cle{rétroaction négative} sur l'hypothalamus (chute de la fréquence des pulsations de GnRH) et sur l'antéhypophyse (baisse de la sensibilité au GnRH). D'où les taux de FSH et LH constamment bas (moins de 5 UI/L), sans variation cyclique.

\textbf{2. Absence d'ovulation.} La FSH étant insuffisante pour stimuler la croissance folliculaire, les follicules ovariens ne se développent pas : l'estradiol ovarien reste bas (< 50 pg/mL). Sans follicule mûr, il n'y a pas de rétrocontrôle positif, donc \emph{pas de pic de LH} ; or le pic de LH est le signal indispensable de l'ovulation. L'ovulation est donc bloquée : c'est le mécanisme contraceptif principal. S'y ajoute l'épaississement de la glaire cervicale (effet du progestatif), qui gêne le passage des spermatozoïdes, et la modification de l'endomètre.

\textbf{3. Saignements d'arrêt.} Pendant les 7 jours sans comprimés actifs, les taux d'hormones de synthèse s'effondrent brutalement (demi-vie courte). L'endomètre, qui était maintenu artificiellement par ces hormones, se desquame : ce sont des hémorragies de privation, appelées « fausses règles », car elles ne succèdent pas à une ovulation ni à une phase lutéale.

\textbf{Conclusion} : la pilule estroprogestative bloque la fécondation en supprimant l'ovulation par rétroaction négative permanente sur le complexe hypothalamo-hypophysaire ; les saignements mensuels sont artificiels et liés à l'interruption de la prise.
\end{exoresolu}

\courssub{Méthode 5 : Interpréter un spermogramme}

\begin{methode}[Lire un spermogramme et conclure (Référentiel OMS 2010/2021)]
\begin{enumerate}
\item \textbf{Comparer chaque paramètre aux valeurs de référence (limites inférieures OMS)} : volume $\geq$ 1,5 mL ; concentration $\geq$ 15 millions de spermatozoïdes/mL ; numération totale $\geq$ 39 millions/éjaculat ; mobilité progressive (a+b) $\geq$ 32 \% ; vitalité $\geq$ 58 \% ; formes typiques (critères stricts de Kruger) $\geq$ 4 \% ; pH $\geq$ 7,2 ; leucocytes $<$ 1 million/mL.
\item \textbf{Nommer l'anomalie avec le vocabulaire exact} : \cle{oligospermie} (concentration $<$ 15 M/mL), \cle{asthénospermie} (mobilité progressive $<$ 32 \%), \cle{tératospermie} (formes typiques $<$ 4 \%), \cle{azoospermie} (absence totale de spermatozoïdes), \cle{hypospermie} (volume $<$ 1,5 mL), \cle{nécrospermie} (vitalité $<$ 58 \%), \cle{leucospermie} (infection/inflammation).
\item \textbf{Rechercher la cause étiologique} : obstruction des voies génitales (azoospermie obstructive), infection (IST, leucospermie), varicocèle (oligo-asthéno-tératospermie), trouble hormonal (hypogonadisme hypogonadotrope : FSH/LH/testostérone basses ; hypergonadotrope : FSH/LH hautes), anomalie génétique (microdélétion Y, caryotype), facteurs toxiques (tabac, alcool, chaleur professionnelle, chimiothérapie), idiopathique.
\item \textbf{Conclure sur la fertilité} et proposer une orientation : traitement de la cause (chirurgie varicocèle, antibiothérapie, traitement hormonal) ou recours à la \cle{procréation médicalement assistée} (PMA) : insémination intra-utérine (IIU), FIV, ICSI.
\end{enumerate}
\textbf{Réflexe} : un seul paramètre anormal suffit à réduire la fertilité ; toujours chiffrer l'écart par rapport à la norme. L'azoospermie impose une exploration hormonale (FSH, testostérone) et génétique avant toute PMA.
\end{methode}

\begin{exoresolu}[Analyse d'un spermogramme au CHU de Yaoundé]
Un homme de 34 ans consulte pour infertilité du couple (3 ans sans grossesse, rapports réguliers non protégés). Son spermogramme (après 3 à 5 jours d'abstinence) donne : volume = 2,0 mL ; concentration = 6 millions de spermatozoïdes/mL ; mobilité progressive = 20 \% ; formes typiques = 8 \% ; vitalité = 65 \% ; pH = 7,8. (1) Analyser chaque paramètre par rapport aux normes OMS. (2) Diagnostiquer les anomalies. (3) Proposer une conduite à tenir étiologique et thérapeutique.
\tcblower
\textbf{(1) Analyse paramètre par paramètre.}
\begin{itemize}
\item Volume : 2,0 mL $\geq$ 1,5 mL : \textbf{normal}.
\item Concentration : 6 millions/mL $<$ 15 millions/mL : \textbf{anormale} (40 \% de la norme).
\item Mobilité progressive : 20 \% $<$ 32 \% : \textbf{anormale}.
\item Formes typiques : 8 \% $\geq$ 4 \% : \textbf{normale}.
\item Vitalité : 65 \% $\geq$ 58 \% : \textbf{normale}.
\item pH : 7,8 $>$ 7,2 : \textbf{légèrement élevé} (suspect d'infection ou obstruction partielle vésicules séminales).
\end{itemize}

\textbf{(2) Diagnostic.} Le patient présente une \cle{oligospermie} (concentration insuffisante) associée à une \cle{asthénospermie} (mobilité insuffisante) : c'est une \cle{oligo-asthénospermie} modérée à sévère. La numération totale de spermatozoïdes mobiles progressifs est de $2,0 \times 6 \times 10^6 \times 0,20 = 2,4 \times 10^6$ (norme $\geq 12,5 \times 10^6$), très insuffisante pour une fécondation naturelle. Le pH élevé oriente vers une infection accessoire ou une dysfonction vésiculaire. Les causes possibles à rechercher : varicocèle (palpation, écho-Doppler), infection des voies génitales (recherche Chlamydia, Gonocoque, Mycoplasmes, culture urines/sperme), trouble hypothalamo-hypophysaire (dosage FSH, LH, testostérone, prolactine), exposition à la chaleur (bains chauds, métier), tabac, anomalie génétique (caryotype, microdélétion Yq si FSH élevée).

\textbf{(3) Conduite à tenir.} Bilan complémentaire : dosage hormonal (FSH, LH, testostérone, prolactine), échographie scrotale (varicocèle), recherche d'infection (PCR urines, culture sperme), caryotype si FSH élevée. Traitement de la cause si identifiée : chirurgie du varicocèle (ligature veine spermatique), antibiothérapie adaptée (doxycycline, azithromycine), arrêt tabac, éviction chaleur. Si l'échec persiste après 6-12 mois ou si cause non traitable : recours à la \cle{PMA}. Ici, oligo-asthénospermie modérée : \textbf{Insémination Intra-Utérine (IIU)} possible si mobilité récupérable après préparation (gradient densité) et trompes perméables (hystérosalpingographie). Si IIU échoue ou trompes bouchées ou oligo-asthénospermie sévère : \textbf{FIV avec ICSI} (injection intracytoplasmique de spermatozoïde) indiquée.

\textbf{Conclusion} : l'infertilité du couple s'explique ici par une oligo-asthénospermie masculine ; l'interprétation chiffrée du spermogramme guide le diagnostic étiologique et le choix de la technique de PMA adaptée.
\end{exoresolu}

\courssub{Méthode 6 : Décrire les techniques de PMA et en discuter les aspects éthiques}

\begin{methode}[Présenter la procréation médicalement assistée (PMA)]
\begin{enumerate}
\item \textbf{Classer les techniques par niveau d'invasivité} :
      \begin{itemize}
      \item \cle{Insémination artificielle (IA / IIU)} : dépôt de sperme préparé (capacitation, sélection mobiles) dans l'utérus au moment de l'ovulation. Indiquée : infertilité masculine modérée (oligo-asthénospermie légère), infertilité cervicale (glaire hostile), infertilité inexpliquée, don de sperme. Nécessite trompes perméables.
      \item \cle{FIV (Fécondation In Vitro)} : mise en contact ovocytes + spermatozoïdes en boîte de Pétri (milieu de culture). Indiquée : trompes bouchées/absentes (obstruction bilatérale), endométriose sévère, échec IA, infertilité masculine modérée.
      \item \cle{ICSI (Injection IntraCytoplasmique de Spermatozoïde)} : injection d'un spermatozoïde unique sélectionné directement dans le cytoplasme de l'ovocyte (micromanipulation). Indiquée : infertilité masculine sévère (oligo-asthénospermie sévère, azoospermie non obstructive avec TESE), échec FIV classique.
      \item \cle{Don de gamètes} (ovocytes ou spermatozoïdes) : insuffisance ovarienne précoce, risque génétique majeur, azoospermie obstructive/non obstructive sans spermatozoïde récupérable.
      \end{itemize}
\item \textbf{Décrire les étapes clés de la FIV/ICSI} : (1) Stimulation ovarienne contrôlée (FSH exogène + antagoniste/agoniste GnRH) $\rightarrow$ (2) Déclenchement ovulation (hCG ou agoniste GnRH) $\rightarrow$ (3) Ponction ovocytaire (écho-guidée, anesthésie) $\rightarrow$ (4) Préparation spermatozoïdes / ICSI $\rightarrow$ (5) Culture embryonnaire (2 à 5 jours, stade 4-8 cellules ou blastocyste) $\rightarrow$ (6) Transfert embryonnaire (1 à 2 embryons, guidé écho) $\rightarrow$ (7) Lutéinisation (progestérone) $\rightarrow$ (8) Test grossesse (bêta-hCG J14).
\item \textbf{Citer les avantages} : permettre à des couples stériles d'avoir un enfant génétiquement lié (sauf don) ; contourner des obstacles mécaniques (trompes) ou gamétiques (sperme) jusque-là irrémédiables ; diagnostic préimplantatoire (DPI) possible pour maladies génétiques graves.
\item \textbf{Discuter les inconvénients et limites éthiques} : destin des embryons surnuméraires (congélation, destruction, don à la recherche) $\rightarrow$ question du statut de l'embryon ; risques de dérives (sélection eugénique par DPI hors pathologie grave, marchandisation des gamètes/ventres porteurs) ; coût très élevé (inégalités d'accès, non remboursé au Cameroun) ; taux de réussite limité par tentative (environ 25-30 \% par transfert, chute après 35-38 ans) ; grossesses multiples (risque prématurité) si transfert > 1 embryon ; syndrome d'hyperstimulation ovarienne (SHSO) ; risque malformatif légèrement majoré (discuté).
\item \textbf{Conclure de façon équilibrée} : la PMA est un progrès médical majeur qui doit rester encadrée par la loi (bioéthique) et l'éthique (respect de l'embryon, non-commercialisation, justice d'accès).
\end{enumerate}
\end{methode}

\begin{exoresolu}[Choisir et justifier une technique de PMA]
Couple 1 : la femme (30 ans) présente une obstruction bilatérale des trompes de Fallope consécutive à une salpingite (infection) ; l'homme a un spermogramme normal. Couple 2 : la femme (32 ans) est normale (ovulation, trompes perméables) ; l'homme présente une oligo-asthénospermie sévère (concentration 2 millions/mL, mobilité progressive 10 \%, formes typiques 3 \%). Pour chaque couple, proposer la technique de PMA la plus adaptée en justifiant, puis citer deux avantages et deux inconvénients éthiques majeurs de la PMA.
\tcblower
\textbf{Couple 1 : FIV (Fécondation In Vitro) classique.} L'obstacle est purement mécanique : les trompes bouchées empêchent la rencontre des gamètes et le transport de l'embryon, mais les ovocytes (femme 30 ans, réserve ovarienne présumée bonne) et les spermatozoïdes (normaux) sont fonctionnels. La FIV court-circuite les trompes : après stimulation ovarienne, les ovocytes sont ponctionnés, mis en contact avec les spermatozoïdes en laboratoire (fécondation « naturelle » in vitro), puis l'embryon obtenu est transféré directement dans la cavité utérine. L'insémination (IIU) serait ici inutile car les spermatozoïdes déposés dans l'utérus ne pourraient jamais atteindre l'ovocyte dans la trompe, ni l'embryon redescendre.

\textbf{Couple 2 : FIV avec ICSI.} Le problème est gamétique et masculin : les spermatozoïdes, trop peu nombreux (2 millions/mL au lieu de 15) et trop peu mobiles (10 \% au lieu de 32 \%), avec une tératospermie associée (3 \% formes typiques $<$ 4 \%), ne peuvent ni atteindre l'ovocyte dans la trompe, ni le pénétrer spontanément (défaillance de la réaction acrosomique ou de la fusion). L'ICSI consiste à injecter un spermatozoïde sélectionné (morphologie, mobilité) directement dans le cytoplasme de l'ovocyte, contournant toutes les barrières de la fécondation naturelle. Les trompes étant fonctionnelles, c'est la qualité du sperme qui impose cette technique de micromanipulation.

\textbf{Avantages de la PMA} : (1) Elle permet à des couples stériles de concrétiser leur désir d'enfant biologique, source d'équilibre familial et social ; (2) Elle contourne des causes d'infertilité jusque-là irrémédiables (trompes détruites par infections pelviennes fréquentes en zone tropicale, azoospermie partielle par TESE).

\textbf{Inconvénients éthiques majeurs} : (1) Le sort des embryons surnuméraires (congélation prolongée, destruction, don à la recherche) pose la question fondamentale du statut de l'embryon humain et du début de la personne ; (2) Le coût très élevé (souvent 1,5 à 3 millions FCFA par tentative au Cameroun, non pris en charge) crée de facto de fortes inégalités d'accès selon le niveau socio-économique, et le risque de dérives marchandes (location d'utérus, vente de gamètes) impose un cadre légal strict et une régulation nationale.

\textbf{Conclusion} : le choix de la technique de PMA dépend du diagnostic précis de la cause de stérilité (bilan complet obligatoire) ; progrès médical remarquable, la PMA doit être pratiquée dans un cadre éthique, légal et médical rigoureux, accessible et transparent.
\end{exoresolu}

\begin{exemplebox}[Dérèglements hormonaux sexuels : exemples cliniques, diagnostic et traitement]
\textbf{1. Hypogonadisme masculin.}
\begin{itemize}
\item \textbf{Hypogonadisme hypogonadotrope (secondaire)} : cause hypothalamo-hypophysaire (tumeur, hyperprolactinémie, syndrome de Kallmann, prise de stéroïdes anabolisants). \emph{Bilan} : Testostérone basse, LH et FSH basses ou inappropriément normales. \emph{Traitement} : cause (chirurgie tumeur, arrêt dopage) ou substitution (testostérone transdermique/injectable) si désir de fertilité $\rightarrow$ hCG + FSH recombinante ou GnRH pulsatile.
\item \textbf{Hypogonadisme hypergonadotrope (primaire)} : cause testiculaire (orchidie bilatérale, séquelles oreillons, chimiothérapie, radiothérapie, syndrome de Klinefelter 47,XXY). \emph{Bilan} : Testostérone basse, LH et FSH \emph{élevées} (perte rétrocontrôle). \emph{Traitement} : substitution androgénique à vie ; fertilité : TESE + ICSI possible si foyers spermatogéniques résiduels (sauf Klinefelter souvent azoospermie complète).
\end{itemize}

\textbf{2. Syndrome des Ovaires Polykystiques (SOPK) — 1ère cause d'infertilité féminine par anovulation.}
\begin{itemize}
\item \textbf{Critères de Rotterdam (2 sur 3)} : (a) Anovulation chronique (oligo-aménorrhée) ; (b) Hyperandrogénie clinique (hirsutisme, acné) ou biologique (testostérone libre élevée, androstènedione) ; (c) Ovaires polykystiques à l'échographie ($>$ 20 follicules/ovaire ou volume $>$ 10 mL).
\item \textbf{Physiopathologie} : Résistance à l'insuline $\rightarrow$ hyperinsulinémie compensatrice $\rightarrow$ stimulation thécaïque (LH relativement élevée / FSH normale/basse $\rightarrow$ ratio LH/FSH $>$ 2-3) $\rightarrow$ surproduction androgènes $\rightarrow$ blocage maturation folliculaire (arrêt stade antral) $\rightarrow$ pas d'ovulation $\rightarrow$ pas de corps jaune $\rightarrow$ pas de progestérone $\rightarrow$ endomètre prolifératif permanent (risque hyperplasie/endométriose).
\item \textbf{Prise en charge} : Hygiène de vie (perte poids 5-10 \% restaure souvent ovulation) ; \textbf{Induction ovulation} : Citrate de clomifène (anti-œstrogène $\rightarrow$ lève rétrocontrôle négatif $\rightarrow$ montée FSH) ou Letrozole (inhibiteur aromatase, 1\textsuperscript{re} ligne actuelle) ; Gonadotrophines (FSH/LH) si échec ; \textbf{Metformine} si résistance insuline avérée ; Chirurgie (drilling ovarien cœlioscopique) en 2\textsuperscript{e} ligne.
\end{itemize}

\textbf{3. Ménopause et Traitement Hormonal Substitutif (THS).}
\begin{itemize}
\item \textbf{Physiologie} : Épuisement folliculaire $\rightarrow$ chute estradiol/inhibine $\rightarrow$ perte rétrocontrôle négatif $\rightarrow$ FSH/LH très élevées (marqueur biologique). Symptômes : bouffées chaleur, troubles sommeil, atrophie urogénitale, ostéoporose, risque cardiovasculaire.
\item \textbf{THS} : Œstrogènes (voie transdermique privilégiée : pas d'effet hépatique 1\textsuperscript{er} passage $\rightarrow$ moins thrombose) + Progestatif (si utérus présent, protection endomètre). Indication : symptômes invalidants, ostéoporose. Contre-indications : cancer sein/endomètre, AVC/IDM antérieur, thrombose active, saignement inexpliqué. Réévaluation annuelle.
\end{itemize}
\end{exemplebox}

\begin{savaistu}[Contexte camerounais : Infertilité, PMA et Santé Reproductive]
\begin{itemize}
\item \textbf{Épidémiologie} : L'infertilité touche 15 à 20 \% des couples en âge de procréer au Cameroun. La cause masculine est retrouvée dans 30-40 \% des cas (infections sexuellement transmissibles mal soignées, varicocèle, expositions professionnelles), la cause féminine dans 30-40 \% (séquelles d'infections pelviennes / salpingites $\rightarrow$ trompes bouchées, SOPK, endométriose), causes mixtes ou inexpliquées pour le reste.
\item \textbf{Accès à la PMA} : Existe dans quelques centres spécialisés (CHU Yaoundé, Douala, cliniques privées). Coût prohibitif pour la majorité (FIV/ICSI $\approx$ 1,5 - 3 millions FCFA). Pas de remboursement par la CNPS ou assurances. Nécessité de développer des protocoles « low-cost » (stimulation douce, FIV naturelle) et une prise en charge publique partielle.
\item \textbf{Prévention} : Lutte contre les IST (dépistage/traitement précoce partenaire, préservatifs), prise en charge rapide des infections pelviennes (éviter séquelles tubaires), éducation sexuelle (âge 1\textsuperscript{re} grossesse, espacement), lutte contre le tabagisme/alcoolisme, dépistage varicocèle à l'adolescence.
\item \textbf{Aspects socioculturels} : Stigmatisation forte de la femme stérile (« femme sans enfant »), pression familiale, recours aux tradipraticiens (retard diagnostic médical). Sensibilisation nécessaire : infertilité = pathologie du couple, prise en charge médicale moderne efficace.
\end{itemize}
\end{savaistu}

\begin{attention}[Les 5 erreurs qui coûtent des points au Bac]
\begin{enumerate}
\item \textbf{Conclure une expérience sans comparer au témoin.} Toute interprétation expérimentale (castration, greffe, injection, parabiose) doit s'appuyer sur la comparaison chiffrée (avec unités) avec le lot témoin ; une conclusion sans témoin est scientifiquement invalide.
\item \textbf{Dater l'ovulation sur le pic d'estradiol.} L'ovulation est déclenchée par le \cle{pic de LH} (vers le 14\textsuperscript{e} jour) ; le pic d'estradiol le précède (J13) et le provoque par rétrocontrôle positif. Inverser les deux ou dire « l'estradiol déclenche l'ovulation » est l'erreur la plus fréquente.
\item \textbf{Oublier le rétrocontrôle positif de l'estradiol.} L'estradiol n'exerce pas toujours une rétroaction négative : à taux élevé et prolongé en fin de phase folliculaire, il \emph{stimule} la décharge de LH (rétrocontrôle positif). Écrire « les hormones ovariennes inhibent toujours l'hypophyse » est faux et pénalisé.
\item \textbf{Confondre contraception et contragestion.} La pilule (estroprogestative ou progestative) empêche l'ovulation (avant la fécondation) ; le RU 486 (mifépristone) agit après la nidation en bloquant les récepteurs de la progestérone. Ne jamais écrire que la pilule « détruit l'embryon » ou « empêche la nidation » comme mécanisme principal (effet secondaire possible sur endomètre, mais principal = blocage ovulation).
\item \textbf{Utiliser un vocabulaire approximatif dans le spermogramme.} « Peu de spermatozoïdes » ne vaut rien : il faut écrire \cle{oligospermie} avec la valeur chiffrée comparée à la norme OMS (15 millions/mL). De même : \cle{asthénospermie} (mobilité), \cle{tératospermie} (morphologie), \cle{azoospermie} (absence) — chaque terme correspond à un paramètre précis. Ne pas confondre concentration et numération totale.
\end{enumerate}
\end{attention}

\newpage

\coursec{Exercices}

\courssub{Je vérifie mes connaissances}

\begin{exercice}[QCM : la régulation hormonale]
Pour chaque question, choisis la (ou les) bonne(s) réponse(s).
\begin{enumerate}
\item Le complexe hypothalamo-hypophysaire qui régule la sécrétion de testostérone comprend :
\begin{enumerate}
\item la GnRH hypothalamique ;
\item la FSH et la LH hypophysaires ;
\item la thyroxine ;
\item les cellules de Leydig et les cellules de Sertoli.
\end{enumerate}
\item Le rétrocontrôle négatif exercé par la testostérone s'exerce sur :
\begin{enumerate}
\item l'hypothalamus et l'hypophyse ;
\item les spermatozoïdes ;
\item les cellules de Leydig uniquement ;
\item la prostate.
\end{enumerate}
\item Le pic de LH au milieu du cycle ovarien :
\begin{enumerate}
\item déclenche l'ovulation ;
\item est provoqué par le rétrocontrôle positif des œstrogènes ;
\item inhibe la sécrétion de progestérone ;
\item survient vers le 14\ieme{} jour d'un cycle de 28 jours.
\end{enumerate}
\item La pilule œstroprogestative agit principalement en :
\begin{enumerate}
\item détruisant les ovocytes ;
\item bloquant l'ovulation par rétrocontrôle négatif permanent ;
\item épaississant la glaire cervicale ;
\item stérilisant définitivement la femme.
\end{enumerate}
\end{enumerate}
\end{exercice}

\begin{exercice}[Vrai ou faux, à justifier]
Indique si chaque affirmation est vraie ou fausse, puis justifie en une ou deux phrases.
\begin{enumerate}
\item La castration d'un rat mâle entraîne une augmentation du taux sanguin de LH.
\item La progestérone est sécrétée par le follicule de De Graaf avant l'ovulation.
\item La section de la tige pituitaire supprime la régulation hypothalamique de l'hypophyse antérieure.
\item La contragestion (pilule du lendemain) est une méthode contraceptive d'usage régulier recommandée.
\end{enumerate}
\end{exercice}

\begin{exercice}[Définitions]
Définis précisément, en une phrase chacun, les termes suivants : \cle{rétrocontrôle négatif}, \cle{ovulation}, \cle{contraception}, \cle{contragestion}, \cle{stérilité}, \cle{procréation médicalement assistée}.
\end{exercice}

\begin{exercice}[Calculs directs : le spermogramme]
Un spermogramme réalisé au Centre Hospitalier d'Essos (Yaoundé) donne les résultats suivants : volume de sperme $V = \SI{3,0}{mL}$ ; concentration en spermatozoïdes $C = \SI{40e6}{spermatozoides/mL}$ ; mobilité progressive $= 45\,\%$.
\begin{enumerate}
\item Calcule le nombre total de spermatozoïdes dans l'éjaculat.
\item Calcule le nombre de spermatozoïdes progressivement mobiles.
\item Sachant que les valeurs de référence de l'OMS sont : concentration $\geq \SI{15e6}{/mL}$ et mobilité progressive $\geq 32\,\%$, ce spermogramme est-il normal ? Justifie.
\end{enumerate}
\end{exercice}

\courssub{Je m'entraîne}

\begin{exercice}[Schéma fonctionnel de la régulation de la testostérone]
Réalise le schéma fonctionnel complet de la régulation du taux de testostérone chez l'homme : tu feras figurer l'hypothalamus, l'hypophyse antérieure, le testicule (cellules de Leydig et de Sertoli), les hormones GnRH, LH, FSH, testostérone et inhibine, ainsi que les flèches de stimulation ($+$) et de rétroaction négative ($-$) correctement placées.
\end{exercice}

\begin{exercice}[Analyse de courbes hormonales]
On a mesuré chez une femme non enceinte, âgée de 26 ans, les taux sanguins d'œstradiol et de progestérone au cours d'un cycle de 28 jours :
\begin{center}
\begin{tabular}{|l|c|c|c|c|c|c|c|}
\hline
\rowcolor{popblueL} Jour du cycle & 1 & 7 & 12 & 14 & 18 & 24 & 28 \\
\hline
Œstradiol (pg/mL) & 30 & 80 & 350 & 120 & 150 & 100 & 40 \\
\hline
Progestérone (ng/mL) & 0,5 & 0,6 & 1,0 & 2,0 & 12 & 10 & 1,0 \\
\hline
\end{tabular}
\end{center}
\begin{enumerate}
\item Trace sur un même graphique les courbes d'évolution des deux hormones (échelles différentes sur deux axes des ordonnées).
\item Identifie les deux phases du cycle et situe approximativement le jour de l'ovulation. Justifie.
\item Quelle structure ovarienne sécrète la progestérone entre les jours 15 et 26 ?
\end{enumerate}
\end{exercice}

\begin{exercice}[Interprétation d'expériences]
Chez le rat, on réalise quatre lots expérimentaux : lot A (témoin) ; lot B (castration) ; lot C (castration puis injection quotidienne d'extrait testiculaire) ; lot D (ablation de l'hypophyse).
\begin{enumerate}
\item Prédis, en le justifiant, l'évolution du taux de LH sanguin dans les lots B et C.
\item Prédis l'état des testicules (taille, spermatogenèse) dans le lot D. Quelle conclusion tires-tu sur le rôle de l'hypophyse ?
\item Quelle expérience complémentaire permettrait de montrer que l'hypophyse est elle-même contrôlée par l'hypothalamus ?
\end{enumerate}
\end{exercice}

\begin{exercice}[Pilule contraceptive et courbes]
Une femme prend une pilule œstroprogestative pendant 21 jours, suivis de 7 jours d'arrêt.
\begin{enumerate}
\item Explique, en t'appuyant sur le rétrocontrôle négatif, pourquoi l'ovulation est bloquée pendant la prise de la pilule.
\item Représente qualitativement l'allure des courbes de FSH et de LH chez cette utilisatrice, et compare-les à celles d'un cycle naturel.
\item Pourquoi des saignements surviennent-ils pendant les 7 jours d'arrêt ?
\end{enumerate}
\end{exercice}

\courssub{Je me prépare au Bac}

\begin{exercice}[Type Bac : courbes hormonales et analogue de la GnRH]
Chez une femme volontaire, on a dosé pendant un cycle de 28 jours les taux plasmatiques de FSH, de LH, d'œstradiol et de progestérone. Les valeurs moyennes sont résumées ci-dessous :
\begin{center}
\begin{tabular}{|l|c|c|c|c|c|c|c|c|}
\hline
\rowcolor{popblueL} Jour & 2 & 6 & 10 & 13 & 14 & 16 & 22 & 27 \\
\hline
FSH (UI/L) & 6 & 8 & 7 & 5 & 9 & 3 & 2 & 4 \\
\hline
LH (UI/L) & 5 & 6 & 8 & 40 & 20 & 5 & 4 & 5 \\
\hline
Œstradiol (pg/mL) & 40 & 90 & 200 & 400 & 150 & 120 & 90 & 40 \\
\hline
Progestérone (ng/mL) & 0,5 & 0,5 & 0,8 & 1,5 & 3 & 10 & 12 & 1 \\
\hline
\end{tabular}
\end{center}
\begin{enumerate}
\item Trace, sur un même repère muni de deux axes des ordonnées, les courbes de la LH et de l'œstradiol en fonction du temps.
\item Détermine le jour de l'ovulation. Justifie ta réponse à partir des données du tableau.
\item Explique, en t'appuyant sur la notion de rétrocontrôle positif, le mécanisme qui relie l'évolution du taux d'œstradiol entre les jours 10 et 13 au pic de LH du jour 13.
\item On administre à cette femme, au jour 10 d'un autre cycle, un analogue de la GnRH à action bloquante sur l'hypophyse. Prédis l'évolution des taux de FSH, de LH et d'œstradiol dans les jours suivants, puis l'effet sur l'ovulation. Justifie.
\item Cite une application médicale de ce type d'analogue dans la prise en charge de l'infertilité ou des troubles hormonaux.
\end{enumerate}
\end{exercice}

\begin{exercice}[Type Bac : infertilité masculine et PMA au Cameroun]
M. et Mme N., mariés depuis 5 ans à Douala, consultent pour absence de grossesse malgré des rapports réguliers non protégés. Le spermogramme de M. N. donne : volume $\SI{2,5}{mL}$ ; concentration $\SI{5e6}{spermatozoides/mL}$ ; mobilité progressive $15\,\%$ ; formes typiques $2\,\%$. Les valeurs de référence de l'OMS (2021) sont : concentration $\geq \SI{16e6}{/mL}$ ; mobilité progressive $\geq 30\,\%$ ; formes typiques $\geq 4\,\%$. Le bilan de Mme N. (hormones, échographie, perméabilité tubaire) est normal.
\begin{enumerate}
\item Analyse le spermogramme de M. N. en comparant chaque paramètre aux valeurs de référence. Conclus en nommant précisément les anomalies (utilise le vocabulaire médical : oligo-, asthéno-, tératozoospermie).
\item Propose deux causes possibles de ces anomalies, l'une d'origine infectieuse et l'autre d'origine hormonale, en précisant le mécanisme de chacune.
\item Parmi les techniques de PMA (insémination intra-utérine, FIV classique, FIV avec ICSI), indique celle qui est la plus indiquée ici et justifie ton choix.
\item Décris en étapes numérotées le protocole de la technique choisie.
\item Dégage un avantage et une limite (médicale, éthique ou économique) de la PMA dans le contexte camerounais.
\end{enumerate}
\end{exercice}

\begin{exercice}[Type Bac : parabiose et régulation testiculaire]
On réalise chez le rat une parabiose (union chirurgicale des circulations sanguines) entre un mâle castré et une femelle intacte. Résultats observés après quelques semaines : chez la femelle, les ovaires se développent fortement, avec de nombreux follicules matures ; chez le mâle castré, le taux de LH sanguin est très élevé, alors que la testostérone reste indétectable. Chez un mâle castré isolé (témoin), le taux de LH est également élevé.
\begin{enumerate}
\item Rappelle, à l'aide d'un schéma fonctionnel légendé, la régulation normale de la sécrétion de testostérone chez le mâle intact.
\item Explique pourquoi le taux de LH est élevé chez le mâle castré isolé.
\item Interprète le développement des ovaires de la femelle en parabiose avec le mâle castré. Quelle déduction tires-tu sur la nature (chimique, voie de transport) du message hypophysaire ?
\item La testostérone reste indétectable chez le mâle castré malgré l'excès de LH : explique pourquoi.
\item À partir de l'ensemble de ces résultats, démontre l'existence d'un rétrocontrôle négatif exercé par les gonades sur l'axe hypothalamo-hypophysaire.
\end{enumerate}
\end{exercice}

\newpage

\coursec{Corrigés des exercices}

\begin{corrige}[Exercice 1 — QCM : la régulation hormonale]
\begin{enumerate}
\item \textbf{Bonnes réponses : a, b et d.} Le complexe hypothalamo-hypophysaire comprend l'\cle{hypothalamus} (qui sécrète la \cle{GnRH}), l'\cle{hypophyse antérieure} (qui sécrète les gonadostimulines \cle{FSH} et \cle{LH}) et leurs cellules cibles testiculaires (cellules de \cle{Leydig}, cibles de la LH, et cellules de \cle{Sertoli}, cibles de la FSH). La thyroxine (c) est une hormone thyroïdienne sans rôle dans cet axe.
\item \textbf{Bonne réponse : a.} La testostérone exerce un \cle{rétrocontrôle négatif} sur l'hypothalamus (freination de la GnRH) et sur l'hypophyse antérieure (freination de la LH). Elle n'agit pas directement sur les spermatozoïdes ni sur la prostate dans ce rétrocontrôle ; les cellules de Leydig sont les productrices, non la cible du rétrocontrôle.
\item \textbf{Bonnes réponses : a, b et d.} Le pic de \cle{LH} déclenche l'\cle{ovulation} (a) ; il est provoqué par le \cle{rétrocontrôle positif} des œstrogènes, sécrétés en forte quantité par le follicule mûr (b) ; il survient vers le 14\ieme{} jour d'un cycle de 28 jours (d). La réponse (c) est fausse : le pic de LH provoque au contraire la transformation du follicule rompu en corps jaune, qui sécrète la progestérone.
\item \textbf{Bonnes réponses : b et c.} La pilule œstroprogestative maintient des taux élevés d'œstrogènes et de progestérone qui exercent un \cle{rétrocontrôle négatif permanent} sur l'axe hypothalamo-hypophysaire : la FSH et la LH restent basses, il n'y a ni maturation folliculaire ni pic de LH, donc pas d'ovulation (b). Elle épaissit aussi la glaire cervicale, ce qui freine le passage des spermatozoïdes (c). Elle ne détruit pas les ovocytes (a) et son effet est réversible à l'arrêt (d).
\end{enumerate}
\end{corrige}

\begin{attention}[Points de vigilance — Exercice 1]
\begin{itemize}
\item Distinguer clairement les \emph{cellules cibles} (Leydig pour LH, Sertoli pour FSH) des \emph{cellules sources} d'hormones (Leydig pour testostérone, Sertoli pour inhibine).
\item Le rétrocontrôle négatif de la testostérone porte sur la \textbf{GnRH} (hypothalamus) et la \textbf{LH} (hypophyse) ; l'inhibine (cellules de Sertoli) cible spécifiquement la \textbf{FSH}.
\item Pour le cycle féminin, le rétrocontrôle \emph{positif} des œstrogènes (seuil élevé maintenu ~48 h) est l'unique cause du pic de LH ovulatoire.
\item La pilule contraceptive n'est pas abortive : elle empêche l'ovulation (et modifie la glaire), elle n'interrompt pas une grossesse déjà implantée.
\end{itemize}
\end{attention}

\begin{corrige}[Exercice 2 — Vrai ou faux, à justifier]
\begin{enumerate}
\item \textbf{Vrai.} La castration supprime la sécrétion de testostérone, donc le rétrocontrôle négatif qu'elle exerce sur l'hypophyse : la sécrétion de \cle{LH} n'est plus freinée et son taux sanguin augmente fortement.
\item \textbf{Faux.} Avant l'ovulation, le follicule de De Graaf sécrète surtout des \cle{œstrogènes} (œstradiol). La \cle{progestérone} est sécrétée massivement après l'ovulation par le \cle{corps jaune}, issu de la transformation du follicule rompu.
\item \textbf{Vrai.} La tige pituitaire contient les vaisseaux du système porte hypothalamo-hypophysaire qui acheminent la \cle{GnRH} jusqu'à l'hypophyse antérieure. Sa section supprime donc le contrôle hypothalamique : les sécrétions de FSH et LH s'effondrent.
\item \textbf{Faux.} La contragestion (pilule du lendemain) est une méthode d'\emph{urgence}, à usage exceptionnel : elle est moins efficace qu'une contraception régulière, plus dosée en hormones (effets secondaires) et ne protège pas des infections sexuellement transmissibles. Elle ne doit en aucun cas remplacer une contraception régulière.
\end{enumerate}
\end{corrige}

\begin{attention}[Points de vigilance — Exercice 2]
\begin{itemize}
\item La castration \emph{mâle} lève le frein sur la LH (et FSH) ; la castration \emph{femelle} (ovariectomie) lève le frein sur FSH et LH (hausse des deux).
\item La section de la tige pituitaire isole l'hypophyse de la GnRH : FSH et LH chutent, entraînant l'atrophie des gonades.
\item Vocabulaire précis : \emph{contraception} (prévention réversible avant fécondation/nidation) vs \emph{contragestion} (urgence post-coïtale, avant nidation confirmée) vs \emph{interruption volontaire de grossesse} (IVG, grossesse installée).
\end{itemize}
\end{attention}

\begin{corrige}[Exercice 3 — Définitions]
\begin{itemize}
\item \cle{Rétrocontrôle négatif} : mécanisme de régulation par lequel une hormone, lorsque son taux sanguin augmente, freine la sécrétion des hormones (hypothalamiques et hypophysaires) qui commandent sa propre production, maintenant ainsi son taux constant.
\item \cle{Ovulation} : libération de l'ovocyte (entouré de cellules folliculaires) par rupture du follicule de De Graaf à la surface de l'ovaire, déclenchée par le pic de LH, vers le 14\ieme{} jour d'un cycle de 28 jours.
\item \cle{Contraception} : ensemble des méthodes visant à empêcher de façon \emph{réversible} la fécondation ou la nidation, sans interruption d'une grossesse déjà commencée.
\item \cle{Contragestion} : méthode d'urgence visant à empêcher la nidation ou à interrompre une toute première étape de la gestation (ex. pilule du lendemain à forte dose de lévonorgestrel).
\item \cle{Stérilité} : incapacité d'un couple à obtenir une grossesse après au moins 12 mois de rapports sexuels réguliers non protégés.
\item \cle{Procréation médicalement assistée (PMA)} : ensemble des techniques médicales (insémination artificielle, fécondation in vitro, ICSI\ldots) permettant à un couple stérile ou infertile d'obtenir une grossesse.
\end{itemize}
\end{corrige}

\begin{corrige}[Exercice 4 — Calculs directs : le spermogramme]
\begin{enumerate}
\item Nombre total de spermatozoïdes :
\[ N = V \times C = \SI{3,0}{mL} \times \SI{40e6}{/mL} = \SI{1,2e8}{} \]
soit \textbf{120 millions de spermatozoïdes} dans l'éjaculat.
\item Nombre de spermatozoïdes progressivement mobiles :
\[ N_{m} = N \times \frac{45}{100} = \num{1,2e8} \times 0,45 = \SI{5,4e7}{} \]
soit \textbf{54 millions de spermatozoïdes progressivement mobiles}.
\item Comparaison aux valeurs de référence de l'OMS (2021) :
\begin{itemize}
\item concentration : $\SI{40e6}{/mL} \geq \SI{16e6}{/mL}$ : valeur \textbf{normale} ;
\item mobilité progressive : $45\,\% \geq 30\,\%$ : valeur \textbf{normale}.
\end{itemize}
\textbf{Conclusion :} ce spermogramme est \textbf{normal} (normozoospermie) pour les deux paramètres étudiés.
\end{enumerate}
\end{corrige}

\begin{attention}[Point de vigilance — Exercice 4]
Ne pas confondre la \emph{concentration} (par mL) avec le \emph{nombre total} (dans tout l'éjaculat) : il faut multiplier par le volume. Respecter les puissances de 10 : $\num{40e6} \times 3 = \num{120e6} = \num{1,2e8}$. Les valeurs de référence OMS 2021 sont : concentration $\ge \SI{16e6}{/mL}$, mobilité progressive $\ge 30\,\%$, formes typiques $\ge 4\,\%$.
\end{attention}

\begin{corrige}[Exercice 5 — Schéma fonctionnel de la régulation de la testostérone]
Le schéma attendu fait apparaître trois étages (hypothalamus, hypophyse antérieure, testicule) et deux boucles de rétroaction négative : la testostérone sur l'hypothalamus et l'hypophyse ; l'inhibine (sécrétée par les cellules de Sertoli) sur la sécrétion de FSH uniquement.

\begin{popfigure}
\begin{center}
\begin{tikzpicture}[
  box/.style={draw=popblue, thick, rounded corners, fill=popblueL, minimum width=3.4cm, minimum height=0.9cm, align=center, font=\small},
  stim/.style={->, thick, popgreen},
  inhib/.style={->, thick, poporange, dashed},
  lbl/.style={font=\scriptsize, fill=white, inner sep=1pt}]
\node[box, align=center] (hypo) at (0,0){\textbf{Hypothalamus}\\ GnRH};
\node[box, align=center] (hyp) at (0,-2.2){\textbf{Hypophyse antérieure}\\ LH \; et \; FSH};
\node[box, align=center] (ley) at (-3,-4.6){\textbf{Cellules de Leydig}\\ testostérone};
\node[box, align=center] (ser) at (3,-4.6){\textbf{Cellules de Sertoli}\\ spermatogenèse, inhibine};
\draw[stim] (hypo) -- node[lbl, right]{$+$ GnRH} (hyp);
\draw[stim] (hyp.west) -| node[lbl, above, pos=0.25]{$+$ LH} (ley.north);
\draw[stim] (hyp.east) -| node[lbl, above, pos=0.25]{$+$ FSH} (ser.north);
\draw[inhib] (ley.west) -- ++(-1.6,0) |- node[lbl, left, pos=0.75]{$-$ testostérone} (hypo.west);
\draw[inhib] (ley.north west) ++(-0.4,0.1) -- ++(-0.9,0) |- node[lbl, right, pos=0.8]{$-$ testostérone} (hyp.west);
\draw[inhib] (ser.east) -- ++(1.6,0) |- node[lbl, right, pos=0.75]{$-$ inhibine} (hyp.east);
\end{tikzpicture}
\end{center}
\legende{Schéma fonctionnel de la régulation du taux de testostérone chez l'homme. Flèches vertes pleines : stimulations ($+$) ; flèches orange pointillées : rétroactions négatives ($-$). La GnRH hypothalamique stimule l'hypophyse antérieure, qui libère LH (cible : cellules de Leydig, productrices de testostérone) et FSH (cible : cellules de Sertoli, siège de la spermatogenèse et productrices d'inhibine).}
\end{popfigure}

\begin{attention}[Points de vigilance — Exercice 5]
\begin{itemize}
\item L'inhibine ne freine \textbf{que} la FSH (pas la LH) : c'est une erreur classique de la faire agir sur toute l'hypophyse.
\item Les flèches de rétroaction doivent \emph{remonter} vers l'hypothalamus et l'hypophyse, jamais descendre vers le testicule.
\item Ne pas oublier de légender chaque flèche avec le nom de l'hormone transportée et son signe ($+$ ou $-$).
\end{itemize}
\end{attention}
\end{corrige}

\begin{corrige}[Exercice 6 — Analyse de courbes hormonales]
\begin{enumerate}
\item Tracé des deux courbes (deux axes des ordonnées : œstradiol en pg/mL à gauche, progestérone en ng/mL à droite) :

\begin{popfigure}
\begin{center}
\begin{tikzpicture}
\begin{axis}[
  width=0.85\linewidth, height=6cm,
  axis y line*=left, axis x line*=bottom,
  xlabel={Jour du cycle}, ylabel={Œstradiol (pg/mL)},
  xmin=0, xmax=29, ymin=0, ymax=400,
  xtick={0,7,14,21,28},
  ylabel style={popblue}, yticklabel style={popblue},
  legend style={at={(0.02,0.98)}, anchor=north west, font=\scriptsize}]
\addplot[popcurve, color=popblue, mark=*] coordinates {(1,30) (7,80) (12,350) (14,120) (18,150) (24,100) (28,40)};
\addlegendentry{Œstradiol}
\end{axis}
\begin{axis}[
  width=0.85\linewidth, height=6cm,
  axis y line*=right, axis x line*=none,
  ylabel={Progestérone (ng/mL)},
  xmin=0, xmax=29, ymin=0, ymax=16,
  ylabel style={poporange}, yticklabel style={poporange},
  legend style={at={(0.98,0.98)}, anchor=north east, font=\scriptsize}]
\addplot[popcurve, color=poporange, mark=square*] coordinates {(1,0.5) (7,0.6) (12,1.0) (14,2.0) (18,12) (24,10) (28,1.0)};
\addlegendentry{Progestérone}
\end{axis}
\end{tikzpicture}
\end{center}
\legende{Évolution des taux d'œstradiol (bleu, axe de gauche) et de progestérone (orange, axe de droite) au cours d'un cycle de 28 jours.}
\end{popfigure}

\item \textbf{Phase folliculaire} (jours 1 à 13 environ) : l'œstradiol augmente fortement (pic à \SI{350}{pg/mL} au jour 12) sous l'effet de la maturation du follicule, tandis que la progestérone reste très basse. \textbf{Phase lutéale} (jours 15 à 28) : la progestérone s'élève nettement (maximum \SI{12}{ng/mL} au jour 18) puis chute en fin de cycle. L'\textbf{ovulation} se situe entre le jour 12 (pic d'œstradiol, qui déclenche le pic de LH) et le jour 14 : on la situe vers le \textbf{jour 13--14}, juste avant l'élévation franche de la progestérone (jour 18 : \SI{12}{ng/mL}), signe de la présence d'un corps jaune fonctionnel.
\item Entre les jours 15 et 26, la progestérone est sécrétée par le \cle{corps jaune} (corps lutéum), structure issue de la transformation du follicule rompu après l'ovulation.
\end{enumerate}
\end{corrige}

\begin{attention}[Points de vigilance — Exercice 6]
\begin{itemize}
\item L'ovulation n'est pas \emph{au} pic d'œstradiol, mais \emph{après} le pic de LH qu'il a provoqué (décalage de 24 à 36 h).
\item La chute de l'œstradiol au jour 14 correspond à la luteinisation du follicule ; la reprise modérée (jours 18--24) est due au corps jaune.
\item La progestérone reste basale ($< \SI{1,5}{ng/mL}$) tant qu'il n'y a pas de corps jaune fonctionnel.
\end{itemize}
\end{attention}

\begin{corrige}[Exercice 7 — Interprétation d'expériences]
\begin{enumerate}
\item \textbf{Lot B (castration) :} la suppression des testicules fait disparaître la testostérone, donc le rétrocontrôle négatif sur l'hypophyse : le taux de \textbf{LH augmente fortement}. \textbf{Lot C (castration + extrait testiculaire) :} l'extrait apporte de la testostérone exogène qui rétablit le rétrocontrôle négatif : le taux de \textbf{LH redescend vers la normale}. Conclusion : le testicule sécrète une substance (la testostérone) qui freine la sécrétion hypophysaire de LH.
\item \textbf{Lot D (ablation de l'hypophyse) :} en l'absence de LH et de FSH, les testicules ne sont plus stimulés : ils \textbf{s'atrophient} (diminution de taille), la \textbf{spermatogenèse s'arrête} et la sécrétion de testostérone s'effondre. \textbf{Conclusion :} l'hypophyse antérieure, par ses gonadostimulines FSH et LH, est indispensable au fonctionnement du testicule (endocrine et exocrine).
\item Pour montrer que l'hypophyse est elle-même contrôlée par l'hypothalamus, on peut réaliser la \textbf{section de la tige pituitaire} (ou une greffe d'hypophyse loin de l'hypothalamus) : les sécrétions de FSH et LH s'effondrent alors que l'hypophyse est intacte, ce qui prouve qu'un message d'origine hypothalamique (la GnRH, véhiculée par le système porte) commande l'hypophyse. On peut aussi injecter de la GnRH à un animal hypophysectomisé partiellement ou à tige sectionnée : la sécrétion de LH repart.
\end{enumerate}
\end{corrige}

\begin{attention}[Points de vigilance — Exercice 7]
\begin{itemize}
\item L'expérience de castration + greffe/extrait démontre l'origine \emph{testiculaire} du facteur inhibiteur (testostérone).
\item L'hypophysectomie démontre la \emph{dépendance} du testicule vis-à-vis de l'hypophyse (trophie et fonction).
\item La section de la tige pituitaire (ou greffe ectopique) démontre la \emph{dépendance} de l'hypophyse vis-à-vis de l'hypothalamus (message humoral : GnRH).
\end{itemize}
\end{attention}

\begin{corrige}[Exercice 8 — Pilule contraceptive et courbes]
\begin{enumerate}
\item La pilule œstroprogestative apporte quotidiennement des œstrogènes et un progestatif à taux sensiblement constants. Ces hormones exercent un \textbf{rétrocontrôle négatif permanent} sur l'hypothalamus et l'hypophyse : la sécrétion de GnRH, puis de FSH et de LH, reste bloquée à un niveau bas. Sans FSH, pas de maturation folliculaire complète ; sans pic de LH, pas d'ovulation. L'ovulation est donc inhibée pendant toute la prise.
\item Allure qualitative des courbes :

\begin{popfigure}
\begin{center}
\begin{tikzpicture}
\begin{axis}[
  width=0.85\linewidth, height=5.5cm,
  xlabel={Jour}, ylabel={Taux de FSH et LH (unités arbitraires)},
  xmin=0, xmax=29, ymin=0, ymax=12,
  xtick={0,7,14,21,28},
  legend style={at={(0.5,1.02)}, anchor=south, font=\scriptsize, legend columns=2}]
\addplot[popcurve, color=popblue, domain=1:28, samples=100] {3 + 6*exp(-((x-14)^2)/2)};
\addlegendentry{LH, cycle naturel}
\addplot[popcurve, color=popgreen, domain=1:28, samples=100] {2.5 + 2*exp(-((x-13)^2)/4) + 1.5*exp(-((x-14)^2)/2)};
\addlegendentry{FSH, cycle naturel}
\addplot[popcurve, color=poporange, dashed, domain=1:28, samples=60] {1.2};
\addlegendentry{LH sous pilule}
\addplot[popcurve, color=poppurple, dashed, domain=1:28, samples=60] {1.5};
\addlegendentry{FSH sous pilule}
\end{axis}
\end{tikzpicture}
\end{center}
\legende{Comparaison qualitative des taux de FSH et LH : cycle naturel (courbes pleines, avec pic pré-ovulatoire vers le jour 14) et cycle sous pilule œstroprogestative (courbes pointillées, taux bas et constants, sans pic).}
\end{popfigure}

Chez l'utilisatrice, les taux de FSH et de LH restent \textbf{bas et sans pic} pendant les 21 jours de prise, contrairement au cycle naturel qui présente un pic de LH (et un pic plus modeste de FSH) au 14\ieme{} jour.
\item Pendant les 7 jours d'arrêt, l'apport en hormones exogènes cesse brutalement : la chute des taux d'œstrogènes et de progestérone prive l'endomètre (qui s'était épaissi) de son soutien hormonal. Il se desquame : ce sont des \textbf{hémorragies de privation}, qui ressemblent à des règles mais ne sont pas de vraies règles (il n'y a pas eu d'ovulation).
\end{enumerate}
\end{corrige}

\begin{attention}[Points de vigilance — Exercice 8]
\begin{itemize}
\item La pilule \emph{bloque} l'axe hypothalamo-hypophysaire : pas de pic de LH = pas d'ovulation. C'est l'effet principal.
\item L'épaississement de la glaire cervicale (effet progestatif) est un effet secondaire contraceptif (barrière mécanique).
\item Les saignements sous pilule sont des \emph{hémorragies de privation} (chute brutale des hormones exogènes), non des menstruations physiologiques (qui font suite à la regression du corps jaune).
\end{itemize}
\end{attention}

\begin{corrige}[Exercice 9 — Type Bac : courbes hormonales et analogue de la GnRH]
\textbf{Barème indicatif (sur 10) :} tracé correct des courbes : 2 pts ; jour de l'ovulation justifié : 2 pts ; explication du rétrocontrôle positif : 2,5 pts ; prédictions avec l'analogue : 2,5 pts ; application médicale : 1 pt.

\begin{enumerate}
\item Tracé des courbes de LH et d'œstradiol :

\begin{popfigure}
\begin{center}
\begin{tikzpicture}
\begin{axis}[
  width=0.85\linewidth, height=6cm,
  axis y line*=left, axis x line*=bottom,
  xlabel={Jour du cycle}, ylabel={LH (UI/L)},
  xmin=0, xmax=29, ymin=0, ymax=45,
  xtick={0,7,14,21,28},
  ylabel style={popblue}, yticklabel style={popblue},
  legend style={at={(0.02,0.98)}, anchor=north west, font=\scriptsize}]
\addplot[popcurve, color=popblue, mark=*] coordinates {(2,5) (6,6) (10,8) (13,40) (14,20) (16,5) (22,4) (27,5)};
\addlegendentry{LH}
\end{axis}
\begin{axis}[
  width=0.85\linewidth, height=6cm,
  axis y line*=right, axis x line*=none,
  ylabel={Œstradiol (pg/mL)},
  xmin=0, xmax=29, ymin=0, ymax=450,
  ylabel style={poporange}, yticklabel style={poporange},
  legend style={at={(0.98,0.98)}, anchor=north east, font=\scriptsize}]
\addplot[popcurve, color=poporange, mark=square*] coordinates {(2,40) (6,90) (10,200) (13,400) (14,150) (16,120) (22,90) (27,40)};
\addlegendentry{Œstradiol}
\end{axis}
\end{tikzpicture}
\end{center}
\legende{Évolution comparée de la LH (bleu, axe de gauche) et de l'œstradiol (orange, axe de droite) au cours du cycle.}
\end{popfigure}

\item L'\textbf{ovulation} a lieu vers le \textbf{jour 14}. Justification à partir du tableau : le pic de LH est atteint au jour 13 (\SI{40}{UI/L}, valeur très supérieure aux autres) ; or l'ovulation suit le pic de LH de quelques heures à 24 h. De plus, la progestérone commence à s'élever nettement à partir du jour 14--16 (passage de \SI{1,5}{} à \SI{3}{} puis \SI{10}{ng/mL}), signe de la formation du corps jaune, et l'œstradiol chute après son maximum du jour 13 (\SI{400}{pg/mL}).
\item Entre les jours 10 et 13, le taux d'œstradiol passe de \SI{200}{} à \SI{400}{pg/mL} : il dépasse un \textbf{seuil} et reste élevé pendant plus de 48 h. À cette concentration élevée et prolongée, l'œstradiol n'exerce plus un rétrocontrôle négatif mais un \textbf{rétrocontrôle positif} sur l'axe hypothalamo-hypophysaire : il stimule la sécrétion de GnRH et rend l'hypophyse très sensible à la GnRH, ce qui provoque la libération massive de LH (pic ovulatoire du jour 13). Ce pic de LH déclenche alors l'ovulation.
\item L'analogue bloquant de la GnRH empêche l'action de la GnRH sur l'hypophyse. Prédictions :
\begin{itemize}
\item \textbf{FSH et LH} : leurs taux chutent fortement dès les jours suivants (l'hypophyse n'est plus stimulée) ;
\item \textbf{Œstradiol} : sans FSH, la maturation folliculaire s'arrête, donc le taux d'œstradiol diminue et reste bas ;
\item \textbf{Ovulation} : elle est \textbf{bloquée}, car il n'y a ni maturation folliculaire complète ni pic de LH.
\end{itemize}
\item Applications médicales : les analogues bloquants (antagonistes de la GnRH) sont utilisés en \textbf{PMA} (protocoles de FIV) pour empêcher une ovulation spontanée prématurée avant le prélèvement des ovocytes ; les analogues agonistes sont utilisés dans le traitement de l'\textbf{endométriose}, des \textbf{fibromes utérins} ou des pubertés précoces.
\end{enumerate}
\end{corrige}

\begin{attention}[Points de vigilance — Exercice 9]
\begin{itemize}
\item Justifier le jour de l'ovulation par \emph{au moins deux indices} du tableau (pic de LH + début d'élévation de la progestérone), pas par la seule affirmation « c'est le 14\ieme{} jour ».
\item Pour le rétrocontrôle positif, préciser impérativement la condition : taux d'œstradiol \textbf{élevé et maintenu} (seuil dépassé pendant environ 48 h).
\item Bien distinguer analogue \emph{bloquant} (antagoniste : inhibition immédiate) et analogue \emph{agoniste} (stimulation puis désensibilisation).
\end{itemize}
\end{attention}

\begin{corrige}[Exercice 10 — Type Bac : infertilité masculine et PMA au Cameroun]
\textbf{Barème indicatif (sur 10) :} analyse du spermogramme et diagnostic : 3 pts ; deux causes mécanistiques : 2 pts ; choix de la technique justifié : 2 pts ; protocole en étapes : 2 pts ; avantage et limite : 1 pt.

\begin{enumerate}
\item Analyse paramètre par paramètre :
\begin{center}
\begin{tabularx}{\linewidth}{|l|X|X|X|}
\hline
\rowcolor{popblueL} \textbf{Paramètre} & \textbf{M. N.} & \textbf{Référence OMS (2021)} & \textbf{Conclusion} \\
\hline
Concentration & \SI{5e6}{/mL} & $\geq \SI{16e6}{/mL}$ & \textbf{oligozoospermie} (sévère) \\
\hline
Mobilité progressive & $15\,\%$ & $\geq 30\,\%$ & \textbf{asthénozoospermie} \\
\hline
Formes typiques & $2\,\%$ & $\geq 4\,\%$ & \textbf{tératozoospermie} \\
\hline
\end{tabularx}
\end{center}
\textbf{Conclusion :} M. N. présente un syndrome \textbf{oligo-asthéno-tératozoospermie (OATS)} : les trois paramètres sont anormaux, ce qui explique l'infertilité du couple, d'autant que le bilan de Mme N. est normal.
\item Causes possibles :
\begin{itemize}
\item \textbf{Cause infectieuse :} une infection sexuellement transmissible non traitée (ex. à \emph{Chlamydia trachomatis} ou \emph{Neisseria gonorrhoeae}) peut provoquer une \textbf{orchi-épididymite} : l'inflammation de l'épididyme perturbe la maturation et la survie des spermatozoïdes (diminution de la mobilité, formes anormales), et des lésions des tubes séminifères réduisent la production.
\item \textbf{Cause hormonale :} un \textbf{déficit de l'axe hypothalamo-hypophysaire} (insuffisance de FSH et/ou de LH, par exemple par tumeur ou dysfonctionnement hypophysaire) prive les cellules de Sertoli et de Leydig de leurs stimulations : la spermatogenèse ralentit (oligozoospermie) et le déficit en testostérone intratesticulaire altère la qualité des spermatozoïdes.
\end{itemize}
(On accepte aussi : varicocèle, cryptorchidie, tabac/alcool, exposition à la chaleur.)
\item La \textbf{FIV avec ICSI} (injection intra-cytoplasmique d'un spermatozoïde) est la plus indiquée : les spermatozoïdes sont trop peu nombreux et trop peu mobiles pour l'insémination intra-utérine ou la FIV classique, qui exigent des millions de spermatozoïdes mobiles. L'ICSI ne requiert qu'\textbf{un seul spermatozoïde vivant} par ovocyte, choisi pour sa morphologie.
\item Protocole de la FIV-ICSI :
\begin{enumerate}
\item Stimulation ovarienne contrôlée de Mme N. par des gonadotrophines (FSH) pour obtenir plusieurs follicules matures, sous surveillance échographique et hormonale ;
\item déclenchement de l'ovulation (injection de hCG) puis \textbf{ponction folliculaire} : prélèvement des ovocytes par voie vaginale sous échographie ;
\item recueil du sperme de M. N. (ou prélèvement testiculaire si nécessaire) et sélection d'un spermatozoïde vivant de morphologie correcte ;
\item \textbf{micro-injection} d'un spermatozoïde directement dans le cytoplasme de chaque ovocyte (ICSI) ;
\item culture \emph{in vitro} des embryons obtenus pendant 3 à 5 jours ;
\item \textbf{transfert} d'un ou deux embryons dans l'utérus de Mme N. ;
\item soutien lutéal (progestérone) puis test de grossesse environ 14 jours après le transfert.
\end{enumerate}
\item \textbf{Avantage :} la PMA permet à des couples stériles d'avoir un enfant génétiquement leur ; des centres existent désormais au Cameroun (Yaoundé, Douala), évitant les coûteux voyages à l'étranger. \textbf{Limites :} coût élevé (souvent plus d'un million de FCFA par tentative) non pris en charge par l'assurance maladie, taux de réussite limité (environ 30 à 40\,\% par cycle), risques médicaux (hyperstimulation ovarienne, grossesses multiples) et questions éthiques (destinée des embryons surnuméraires, pression sociale).
\end{enumerate}
\end{corrige}

\begin{attention}[Points de vigilance — Exercice 10]
\begin{itemize}
\item Employer le vocabulaire médical exact : \emph{oligozoospermie} (concentration), \emph{asthénozoospermie} (mobilité), \emph{tératozoospermie} (morphologie) ; comparer \emph{chaque} valeur chiffrée à la référence.
\item Justifier le choix de l'ICSI par la \emph{gravité} des anomalies (un seul spermatozoïde suffit), pas seulement le nommer.
\item Pour les causes, exiger le \textbf{mécanisme} (comment la cause produit l'anomalie), pas seulement le nom de la cause.
\end{itemize}
\end{attention}

\begin{corrige}[Exercice 11 — Type Bac : parabiose et régulation testiculaire]
\textbf{Barème indicatif (sur 10) :} schéma fonctionnel : 2,5 pts ; LH élevée chez le castré : 1,5 pts ; interprétation de la parabiose et nature du message : 2,5 pts ; testostérone indétectable : 1,5 pts ; démonstration du rétrocontrôle : 2 pts.

\begin{enumerate}
\item Schéma fonctionnel de la régulation normale chez le mâle intact :

\begin{popfigure}
\begin{center}
\begin{tikzpicture}[
  box/.style={draw=popblue, thick, rounded corners, fill=popblueL, minimum width=3.2cm, minimum height=0.85cm, align=center, font=\small},
  stim/.style={->, thick, popgreen},
  inhib/.style={->, thick, poporange, dashed},
  lbl/.style={font=\scriptsize, fill=white, inner sep=1pt}]
\node[box, align=center] (hypo) at (0,0){\textbf{Hypothalamus}\\ GnRH};
\node[box, align=center] (hyp) at (0,-2){\textbf{Hypophyse antérieure}\\ LH + FSH};
\node[box, align=center] (test) at (0,-4){\textbf{Testicule}\\ Leydig : testostérone \\ Sertoli : spermatogenèse};
\draw[stim] (hypo) -- node[lbl, right]{$+$ GnRH} (hyp);
\draw[stim] (hyp) -- node[lbl, right]{$+$ LH, $+$ FSH} (test);
\draw[inhib] (test.west) -- ++(-2.2,0) |- node[lbl, left, pos=0.75]{$-$ testostérone} (hypo.west);
\draw[inhib] (test.east) -- ++(2.2,0) |- node[lbl, right, pos=0.75]{$-$ testostérone, $-$ inhibine} (hyp.east);
\end{tikzpicture}
\end{center}
\legende{Régulation de la sécrétion de testostérone chez le mâle intact : stimulations ($+$) en vert ; rétroactions négatives ($-$) de la testostérone et de l'inhibine en orange pointillé.}
\end{popfigure}

\item Chez le mâle castré isolé, l'ablation des testicules supprime la production de testostérone. Le rétrocontrôle négatif que celle-ci exerçait sur l'hypothalamus et l'hypophyse disparaît : la GnRH et surtout la LH ne sont plus freinées, d'où un \textbf{taux de LH sanguin très élevé}.
\item En parabiose, les sangs des deux animaux se mélangent : la LH (et la FSH) produites en excès par l'hypophyse du mâle castré passent dans la circulation de la femelle. Ces gonadostimulines stimulent ses ovaires : forte \textbf{folliculogenèse} et développement de nombreux follicules matures. \textbf{Déduction :} le message hypophysaire est de nature \textbf{chimique} (hormonale : FSH et LH, des glycoprotéines) et il est transporté par \textbf{voie sanguine} (il franchit la circulation commune de la parabiose) ; ce n'est pas un message nerveux.
\item La LH ne peut agir que sur ses cellules cibles, les \textbf{cellules de Leydig}, qui sont situées dans les testicules. Or le mâle est castré : les testicules sont absents. Même en très grande quantité, la LH ne trouve aucune cible : la testostérone reste donc \textbf{indétectable}. Cela confirme que la LH n'est pas elle-même la testostérone et qu'elle agit indirectement, via le testicule.
\item Démonstration du rétrocontrôle négatif :
\begin{itemize}
\item Chez le mâle \textbf{intact} (situation normale), la testostérone est présente et le taux de LH est normal : la gonade freine l'hypophyse.
\item Chez le mâle \textbf{castré} (isolé ou en parabiose), la testostérone disparaît et le taux de LH s'élève fortement : la suppression de la gonade lève la freination de l'axe hypothalamo-hypophysaire.
\item Le passage de cette LH en excès dans le sang de la femelle (parabiose) et l'effet observé sur ses ovaires prouvent que cette hyperactivité hypophysaire est bien réelle et d'origine humorale.
\end{itemize}
L'ensemble montre que les gonades exercent, par leurs hormones (testostérone chez le mâle), un \textbf{rétrocontrôle négatif} sur l'axe hypothalamo-hypophysaire : quand l'hormone gonadique manque, les gonadostimulines s'accumulent ; quand elle est présente, elle les maintient à un niveau bas.
\end{enumerate}
\end{corrige}

\begin{attention}[Points de vigilance — Exercice 11]
\begin{itemize}
\item Dans l'interprétation de la parabiose, conclure explicitement sur les \textbf{deux} aspects demandés : nature \emph{chimique} du message ET transport par \emph{voie sanguine}.
\item Ne pas écrire que « la LH se transforme en testostérone » : la LH est une \emph{stimulation}, la testostérone est le \emph{produit} des cellules de Leydig.
\item La démonstration finale doit articuler les trois observations (intact, castré, parabiose) en un raisonnement comparatif, et non les juxtaposer.
\end{itemize}
\end{attention}

\newpage

\coursec{Fiche bilan}

\courssub{Carte mentale de la leçon}

\begin{popfigure}
\begin{tikzpicture}[mindmap, grow cyclic, every node/.style={concept, text=white, font=\small},
  root concept/.append style={concept color=popink, font=\bfseries},
  level 1/.append style={level distance=3.4cm, sibling angle=72},
  level 2/.append style={level distance=2.2cm, sibling angle=45, font=\scriptsize}]
\node[root concept]{Hormones sexuelles : régulation et troubles}
  child[concept color=popblue]{ node {Régulation de la testostérone}
    child { node {GnRH hypothalamique} }
    child { node {LH et FSH hypophysaires} }
    child { node {Rétroaction négative} }
  }
  child[concept color=poporange]{ node {Cycles sexuels de la femme}
    child { node {Phases folliculaire et lutéale} }
    child { node {\OE stradiol et progestérone} }
    child { node {Pic de LH : ovulation} }
  }
  child[concept color=popgreen]{ node {Maîtrise de la reproduction}
    child { node {Contraception (pilule)} }
    child { node {Contragestion (RU 486)} }
  }
  child[concept color=poppurple]{ node {Dérèglements hormonaux}
    child { node {Causes et conséquences} }
    child { node {Traitement et prévention} }
  }
  child[concept color=poppink]{ node {Stérilité et PMA}
    child { node {Causes chez l'homme et la femme} }
    child { node {FIV, IAC, ICSI} }
    child { node {Enjeux éthiques} }
  };
\end{tikzpicture}
\legende{Carte mentale de la leçon : le fonctionnement de l'axe hypothalamo-hypophyso-gonadique (bleu, orange) fonde la maîtrise de la reproduction (vert), la compréhension des dérèglements (violet) et la prise en charge de la stérilité (rose).}
\end{popfigure}

\courssub{L'essentiel à retenir}

\begin{bilan}
\begin{itemize}
  \item \cle{Testostérone} : hormone stéroïde sécrétée par les \textbf{cellules de Leydig} (tissu interstitiel du testicule) ; responsable des caractères sexuels secondaires et du maintien de la spermatogenèse.
  \item \cle{GnRH} : neurohormone hypothalamique qui stimule la sécrétion des gonadotrophines hypophysaires \textbf{LH} et \textbf{FSH}.
  \item \cle{Rétroaction négative} (rétrocontrôle) : la testostérone (chez l'homme), l'\oe stradiol et la progestérone (chez la femme) inhibent l'hypothalamus et l'hypophyse, ce qui stabilise leur propre taux sanguin.
  \item \cle{Cycle sexuel féminin} : succession de variations cycliques (environ \num{28} jours) de l'ovaire (cycle ovarien : phases folliculaire puis lutéale) et de l'utérus (cycle utérin : règles, prolifération, sécrétion).
  \item \cle{Ovulation} : libération de l'ovocyte vers le \num{14}\ieme{} jour, déclenchée par le \textbf{pic de LH}.
  \item \cle{Contraception} : ensemble des méthodes empêchant la fécondation (pilule \oe stroprogestative : blocage de l'ovulation par rétroaction négative permanente).
  \item \cle{Contragestion} : interruption d'un début de grossesse (ex. RU 486, antiprogestérone qui bloque les récepteurs de la progestérone et provoque le décollement de l'endomètre).
  \item \cle{Stérilité} : impossibilité pour un couple d'obtenir une grossesse après 12 à 24 mois de rapports non protégés ; on parle d'\textbf{infertilité} quand la grossesse n'aboutit pas.
  \item \cle{PMA} : procréation médicalement assistée (insémination artificielle, FIVETE, ICSI).
\end{itemize}
\end{bilan}

\begin{bilan}
\begin{center}
\begin{tabularx}{\linewidth}{|l|X|X|}
\hline
\rowcolor{popblueL} \textbf{Étape} & \textbf{Chez l'homme} & \textbf{Chez la femme} \\
\hline
Hypothalamus & GnRH (sécrétion continue) & GnRH (sécrétion pulsatile) \\
\hline
Hypophyse & LH $\to$ cellules de Leydig ; FSH $\to$ cellules de Sertoli & FSH $\to$ croissance folliculaire ; LH $\to$ ovulation et corps jaune \\
\hline
Gonade & Testostérone (taux stable) & \OE stradiol (phase folliculaire) puis progestérone (phase lutéale) \\
\hline
Rétrocontrôle & Négatif, permanent & Négatif ; \textbf{positif} de l'\oe stradiol juste avant l'ovulation (pic de LH) \\
\hline
\end{tabularx}
\end{center}
\end{bilan}

\begin{bilan}
\begin{itemize}
  \item \textbf{Interpréter une expérience} : comparer témoin et opéré (castration, hypophysectomie, section de la tige pituitaire) ; conclure sur l'organe responsable de la sécrétion et sur le sens du rétrocontrôle.
  \item \textbf{Analyser les courbes du cycle} : repérer le pic de LH (ovulation), la chute de progestérone en fin de cycle (règles), les deux pics d'\oe stradiol.
  \item \textbf{Expliquer la pilule} : apport permanent d'\oe stroprogestatifs $\to$ rétroaction négative continue $\to$ pas de pic de LH $\to$ pas d'ovulation.
  \item \textbf{Lire un spermogramme} : normes OMS : volume $\geq$ \SI{1,5}{mL}, concentration $\geq$ \num{15} millions de spermatozoïdes/mL, mobilité $\geq$ \SI{40}{\percent}, formes typiques $\geq$ \SI{4}{\percent}.
\end{itemize}
\end{bilan}

\begin{aretenir}[Checklist avant le Bac]
\begin{itemize}
  \item $\square$ Je sais citer les organes de l'axe hypothalamo-hypophyso-testiculaire et le rôle de chacun.
  \item $\square$ Je sais interpréter une castration, une injection d'extrait testiculaire et une hypophysectomie.
  \item $\square$ Je sais dessiner le schéma fonctionnel de la régulation du taux de testostérone avec les flèches $+$ et $-$.
  \item $\square$ Je connais les deux phases du cycle ovarien et les hormones dominantes de chacune.
  \item $\square$ Je sais repérer le pic de LH et la date d'ovulation sur une courbe.
  \item $\square$ Je sais expliquer le rétrocontrôle positif de l'\oe stradiol avant l'ovulation.
  \item $\square$ Je sais distinguer contraception et contragestion et expliquer leur mode d'action.
  \item $\square$ Je connais les principales causes de stérilité chez l'homme et chez la femme.
  \item $\square$ Je sais interpréter un spermogramme avec les normes de l'OMS.
  \item $\square$ Je sais décrire l'IAC, la FIVETE et l'ICSI et discuter leurs enjeux éthiques.
  \item $\square$ Je connais les conséquences des dérèglements hormonaux (ménopause, andropause, troubles du cycle) et leur prévention.
  \item $\square$ Je rédige mes réponses avec le vocabulaire exact : GnRH, LH, FSH, rétroaction, corps jaune, endomètre.
\end{itemize}
\end{aretenir}

\findecours

\end{document}
